\chapter*{导数专题练习}
\begin{enumerate}
    \item 瞬时速度、平均速度
    \begin{enumerate}
        \item 若质点 $A$ 运动的位移 $S$ (单位: ${m}$ )与时间 $t$ (单位: ${s}$ )之间的函数关系是 $S\left( t\right)  =  - \dfrac{2}{t}\left( {t \geq  1}\right)$ ,那么该质点在 $t = 3{s}$ 时的瞬时速度和从 $t = 1{s}$ 到 $t = 3{s}$ 这两秒内的平均速度分别为\blkda{$\dfrac{2}{9},\dfrac{2}{3}$}

        \begin{jieda}
            $\dfrac{\Delta S}{\Delta t} = \dfrac{S\left( {3 + {\Delta t}}\right)  - S\left( 3\right) }{\Delta t} = \dfrac{\dfrac{2}{3 + {\Delta t}} + \dfrac{2}{3}}{\Delta t} = \dfrac{2}{3\left( {3 + {\Delta t}}\right) }$ ,
            
            所以 $\mathop{\lim }\limits_{{{\Delta t} \rightarrow  0 }}\dfrac{\Delta S}{\Delta t} = \mathop{\lim }\limits_{{{\Delta t} \rightarrow  0 }}\dfrac{2}{3\left( {3 + {\Delta t}}\right) } = \dfrac{2}{9}$ 即该质点在 $t = {3s}$ 时的瞬时速度为 $\dfrac{2}{9}$ ;
            
            从 $t = 1{s}$ 到 $t = 3{s}$ 这两秒内的平均速度为 $\dfrac{S\left( 3\right)  - S\left( 1\right) }{3 - 1} = \dfrac{-\dfrac{2}{3} + 2}{2} = \dfrac{2}{3}$
        \end{jieda}
        
        \item 某物体的运动方程为 $s\left( t\right)  = 4{t}^{2} + 2$ (位移单位:m,时间单位:s), \\若 $v = \mathop{\lim }\limits_{{{\Delta t} \rightarrow  0 }}\dfrac{s\left( {3 + {\Delta t}}\right)  - s\left( 3\right) }{\Delta t} = {24}{m}/{s}$ ,则下列说法中正确的是\dotfill{(\kh{C})}
        \xx{${24}{m}/{s}$ 是物体从开始到 $3{s}$ 这段时间内的平均速度}{${24}{m}/{s}$ 是物体从 $3{s}$ 到 $\left( {3 + {\Delta t}}\right) {s}$ 这段时间内的速度}{${24}{m}/{s}$ 是物体在 $3{s}$ 这一时刻的瞬时速度}{${24}{m}/{s}$ 是物体从 $3{s}$ 到 $\left( {3 + {\Delta t}}\right) {s}$ 这段时间内的平均速度}
        \begin{jieda}
            由 $v = \mathop{\lim }\limits_{{{\Delta t} \rightarrow  0 }}\dfrac{s\left( {3 + {\Delta t}}\right)  - s\left( 3\right) }{\Delta t} = {24}{m}/{s}$，可知，$24m/s$ 是物体在 ${{3s}}$ 这一时刻的瞬时速度.
        \end{jieda}
        
        \item 如果说某物体做直线运动的时间与距离满足 $s\left( t\right)  = 2{\left( 1 - t\right) }^{2}$ , 则其在 $t = {0.5}$ 时的瞬时速度为\blkda{$-2$}

        \begin{jieda}
            根据导数的定义可得,在 $t = {0.5}$ 时的瞬时速度为${s}^{\prime }\left( t\right)  = \mathop{\lim }\limits_{{{\Delta t} \rightarrow  0 }}\dfrac{s\left( {{0.5} + {\Delta t}}\right)  - s\left( {0.5}\right) }{\Delta t} = \mathop{\lim }\limits_{{{\Delta t} \rightarrow  0 }}\dfrac{2{\left( 1 - {0.5} - \Delta t\right) }^{2} - 2{\left( 1 - {0.5}\right) }^{2}}{\Delta t} =  - 2$
        \end{jieda}
    \end{enumerate}
    \item 平均变化率、瞬时变化率
    \begin{enumerate}
        \item 若函数 $f\left( x\right)  = {x}^{2} - x$ ,则函数 $f\left( x\right)$ 从 $x = 1$ 到 $x = 3$ 的平均变化率为\blkda{3}
        \begin{jieda}
            因为 $f\left( x\right)  = {x}^{2} - x$ ,所以 $f\left( 1\right)  = {1}^{2} - 1 = 0,f\left( 3\right)  = {3}^{2} - 3 = 6$，故函数 $f\left( x\right)$ 从 $x = 1$ 到 $x = 3$ 的平均变化率为 $\dfrac{\Delta y}{\Delta x} = \dfrac{f\left( 3\right)  - f\left( 1\right) }{3 - 1} = \dfrac{6 - 0}{2} = 3$ .
        \end{jieda}
        
        \item 函数 $f\left( x\right)  = {x}^{2}$ 在 $x = 2$ 处的瞬时变化率为\blkda{4}
        \begin{jieda}
            因为 $\dfrac{f\left( {2 + {\Delta x}}\right)  - f\left( 2\right) }{\Delta x} = \dfrac{{\left( 2 + \Delta x\right) }^{2} - 4}{\Delta x} = \dfrac{4 + {4\Delta x} + {\left( \Delta x\right) }^{2} - 4}{\Delta x} = {\Delta x} \neq  2$，所以,函数 $f\left( x\right)  = {x}^{2}$ 在 $x = 2$ 处的瞬时变化率为 ${f}^{\prime }\left( 2\right)  = \mathop{\lim }\limits_{{{\Delta x} \rightarrow  0 }}\dfrac{f\left( {2 + {\Delta x}}\right)  - f\left( 2\right) }{\Delta x}$
            
            $= \mathop{\lim }\limits_{{{\Delta x} \rightarrow  0 }}\left( {{\Delta x} + 4}\right)  = 4.$
        \end{jieda}
        
        \item 下列函数中，在区间$[-1,0]$上的平均变化率最大的是\dotfill{(\kh{B})}
        \xx{$y = {x}^{2}$}{$y = {x}^{3}$}{$y = {\left( \dfrac{1}{2}\right) }^{x}$}{$y = {2}^{x}$}
        
        \begin{jieda}
            对于 ${A},y = {x}^{2}$ 在 $\left\lbrack  {-1,0}\right\rbrack$ 上的平均变化率为 $\dfrac{0 - 1}{0 - \left( {-1}\right) } =  - 1$ ,
            
            对于 ${B},y = {x}^{3}$ 在 $\left\lbrack  {-1,0}\right\rbrack$ 上的平均变化率为 $\dfrac{0 - \left( {-1}\right) }{0 - \left( {-1}\right) } = 1$ ,
            
            对于 ${C},y = {\left( \dfrac{1}{2}\right) }^{x}$ 在 $\left\lbrack  {-1,0}\right\rbrack$ 上的平均变化率为 $\dfrac{1 - {\left( \dfrac{1}{2}\right) }^{-1}}{0 - \left( {-1}\right) } =  - 1$ ,
            
            对于 ${D},y = {2}^{x}$ 在 $\left\lbrack  {-1,0}\right\rbrack$ 上的平均变化率为 $\dfrac{1 - {2}^{-1}}{0 - \left( {-1}\right) } = \dfrac{1}{2}$ ,
            
            故 $y = {x}^{3}$ 在 $\left\lbrack  {-1,0}\right\rbrack$ 上的平均变化率最大
        \end{jieda}
        
        \item 已知某容器的高度为 20cm，现在向容器内注入液体，且容器内液体的高度 $h$ (单位: ${{cm}}$ ) 与时间 $t$ (单位: ${s}$ ) 的函数关系式为 $h = \dfrac{1}{3}{t}^{3} + {t}^{2}$ ,当 $t = {t}_{0}$ 时,液体上升高度的瞬时变化率为 $3{{cm}}/{s}$ ,则当 $t = {t}_{0} + 1$ 时，液体上升高度的瞬时变化率为\dotfill{(\kh{C})}
        \xx{$5{{cm}}/{s}$}{$6{{cm}}/{s}$}{$8{{cm}}/{s}$}{${10}{{cm}}/{s}$}

        \begin{jieda}
            由 $h = \dfrac{1}{3}{t}^{3} + {t}^{2}$ ,求导得: ${h}^{\prime } = {t}^{2} + {2t}$ .
            
            当 $t = {t}_{0}$ 时， ${h}^{\prime } = {t}_{0}^{2} + 2{t}_{0} = 3$ ，解得 ${t}_{0} = 1\left( {{t}_{0} =  - 3}\right.$ 舍去 $)$ .
            
            故当 $t = {t}_{0} + 1 = 2$ 时,液体上升高度的瞬时变化率为 ${2}^{2} + 2 \times  2 = 8{{cm}}/{s}$ .
        \end{jieda}
        \item 一个球的半径以 2 厘米/秒的速度在增长，当此球的半径为 6 厘米时，其体积的变化率为\dotfill{(\kh{B})}
        \xx{$576\pi {cm}^{3}/$ 秒}{$288\pi {cm}^{3}/$ 秒}{$144\pi {cm}^{3}/$ 秒}{${96\pi } {cm}^{3}/$ 秒}
        
        \item 一个圆的半径以 0.1 厘米/秒的速度在增长，若圆的周长为 $C$ 厘米，则其面积的变化率为\dotfill{(\kh{B})}
        \xx{$0.2\pi C {cm}^{2}/$ 秒}{$0.1C {cm}^{2}/$ 秒}{$\dfrac{0.1C}{2\pi} {cm}^{2}/$ 秒}{$-\dfrac{0.1C}{2\pi} {cm}^{2}/$ 秒}
    \end{enumerate}
    
    \item 利用导数的定义解题
    \begin{enumerate}
        \item 若 ${f}^{\prime }\left( 2\right)  = 4$ ，则 $\mathop{\lim }\limits_{{{\Delta x} \rightarrow  0 }}\dfrac{f\left( {2 + {\Delta x}}\right)  - f\left( {2 - {\Delta x}}\right) }{\Delta x} =$ \blkda{8}
        \begin{jieda}
            由题意知, ${f}^{\prime }\left( 2\right)  = 4$，则 $\mathop{\lim }\limits_{{{\Delta x} \rightarrow  0 }}\dfrac{f\left( {2 + {\Delta x}}\right)  - f\left( {2 - {\Delta x}}\right) }{\Delta x} = 2\mathop{\lim }\limits_{{{\Delta x} \rightarrow  0 }}\dfrac{f\left( {2 + {\Delta x}}\right)  - f\left( {2 - {\Delta x}}\right) }{2\Delta x} = 2{f}^{\prime }\left( 2\right)  = 8$ .
        \end{jieda}
        
        \item 设 $f\left( x\right)$ 在 $R$ 上的导函数为 ${f}^{\prime }\left( x\right)$ ,若 $\mathop{\lim }\limits_{{{\Delta x} \rightarrow  0 }}\dfrac{f\left( {3 - {\Delta x}}\right)  - f\left( 3\right) }{3\Delta x} = 2$，则 ${f}^{\prime }\left( 3\right)  =$\blkda{$-6$}
        \begin{jieda}
            由于 $\mathop{\lim }\limits_{{{\Delta x} \rightarrow  0 }}\dfrac{f\left( {3 - {\Delta x}}\right)  - f\left( 3\right) }{3\Delta x} =  - \dfrac{1}{3}\mathop{\lim }\limits_{{{\Delta x} \rightarrow  0 }}\dfrac{f\left( {3 - {\Delta x}}\right)  - f\left( 3\right) }{-{\Delta x}} =  - \dfrac{1}{3}{f}^{\prime }\left( 3\right)  = 2$ ,则 ${f}^{\prime }\left( 3\right)  =  - 6$ .
        \end{jieda}
        
        \item 若 $f\left( x\right)  = \ln x + {x}^{3}$ ，则 $\mathop{\lim }\limits_{{{\Delta x} \rightarrow  0 }}\dfrac{f\left( {1 + {\Delta x}}\right)  - f\left( 1\right) }{2\Delta x} =$\blkda{2}
        \begin{jieda}
            根据导数的定义可得 $\mathop{\lim }\limits_{{{\Delta x} \rightarrow  0 }}\dfrac{f\left( {1 + {\Delta x}}\right)  - f\left( 1\right) }{2\Delta x} = \dfrac{1}{2}\mathop{\lim }\limits_{{{\Delta x} \rightarrow  0 }}\dfrac{f\left( {1 + {\Delta x}}\right)  - f\left( 1\right) }{\Delta x} = \dfrac{1}{2}{f}^{\prime }\left( 1\right)$ ,
            
            由 $f\left( x\right)  = \ln x + {x}^{3}$ 可得 ${f}^{\prime }\left( x\right)  = \dfrac{1}{x} + 3{x}^{2}$ ,
            
            可得 $\dfrac{1}{2}{f}^{\prime }\left( 1\right)  = \dfrac{1}{2}\left( {\dfrac{1}{1} + 3 \times  {1}^{2}}\right)  = 2$ ,即 $\mathop{\lim }\limits_{{{\Delta x} \rightarrow  0 }}\dfrac{f\left( {1 + {\Delta x}}\right)  - f\left( 1\right) }{2\Delta x} = 2$ .
        \end{jieda}

    \end{enumerate}

    \item 求曲线切线的斜率 (倾斜角)
    \begin{enumerate}
        \item 设 $f\left( x\right)$ 是可导函数,且 $\mathop{\lim }\limits_{{{\Delta x} \rightarrow  0 }}\dfrac{f\left( {1 + {3\Delta x}}\right)  - f\left( 1\right) }{\Delta x} = 2$ ,则 $f\left( x\right)$ 在 $x = 1$ 处的切线的斜率等于\blkda{$\dfrac{2}{3}$}
        \begin{jieda}
            因为 $\mathop{\lim }\limits_{{{\Delta x} \rightarrow  0 }}\dfrac{f\left( {1 + {3\Delta x}}\right)  - f\left( 1\right) }{\Delta x} = 2$ ,所以 ${f}^{\prime }\left( 1\right)  = \mathop{\lim }\limits_{{{\Delta x} \rightarrow  0 }}\dfrac{f\left( {1 + {3\Delta x}}\right)  - f\left( 1\right) }{3\Delta x} = \dfrac{1}{3}\mathop{\lim }\limits_{{{\Delta x} \rightarrow  0 }}\dfrac{f\left( {1 + {3\Delta x}}\right)  - f\left( 1\right) }{\Delta x} = \dfrac{2}{3}$ , 由导数的几何意义知, $f\left( x\right)$ 在 $x = 1$ 处的切线的斜率为 $\dfrac{2}{3}$
        \end{jieda}
        
        \item 曲线 $f\left( x\right)  = \dfrac{9}{x}$ 在点 $\left( {3,3}\right)$ 处的切线的倾斜角 $\alpha$ 等于\blkda{${135}^{ \circ  }$}
        \begin{jieda}
            ${f}^{\prime }\left( x\right)  = \mathop{\lim }\limits_{{{\Delta x} \rightarrow  0 }}\dfrac{f\left( {x + {\Delta x}}\right)  - f\left( x\right) }{\Delta x} = 9\mathop{\lim }\limits_{{{\Delta x} \rightarrow  0 }}\dfrac{\dfrac{1}{x + {\Delta x}}}{\Delta x} =  - 9\mathop{\lim }\limits_{{{\Delta x} \rightarrow  0 }}\dfrac{1}{\left( {x + {\Delta x}}\right) {\Delta x}} =  - \dfrac{9}{{x}^{2}}$ ,
            
            所以 ${f}^{\prime }\left( 3\right)  =  - 1$ . 又切线的倾斜角 $\alpha$ 的范围为 ${0}^{ \circ  } \leq  \alpha  \leq  {180}^{ \circ  }$ ,所以所求倾斜角为 ${135}^{ \circ  }$ .
        \end{jieda}

        \item 函数 $y = f\left( x\right)$ 的图象如图所示，下列关系正确的是\dotfill{(\kh{C})}
        
        \begin{center}
        \includegraphics[width=0.25\textwidth]{bodyxb/src/bo_d5nlf73ef24c73bmm510_4_143_1619_379_287_0.jpg}
        \end{center}

        \xx{$0 < {f}^{\prime }\left( 4\right)  < {f}^{\prime }\left( 5\right)  < f\left( 5\right)  - f\left( 4\right)$}{$0 < {f}^{\prime }\left( 4\right)  < f\left( 5\right)  - f\left( 4\right)  < {f}^{\prime }\left( 5\right)$}{$0 < {f}^{\prime }\left( 5\right)  < f\left( 5\right)  - f\left( 4\right)  < {f}^{\prime }\left( 4\right)$}{$0 < f\left( 5\right)  - f\left( 4\right)  < {f}^{\prime }\left( 5\right)  < {f}^{\prime }\left( 4\right)$}
        \begin{jieda}
            根据题意 ${f}^{\prime }\left( 4\right)$ 的几何意义为 $y = f\left( x\right)$ 在点 $B$ 处切线的斜率,
            
            ${f}^{\prime }\left( 5\right)$ 的几何意义为 $y = f\left( x\right)$ 在点 $A$ 处切线的斜率,
            
            $f\left( 5\right)  - f\left( 4\right)  = \dfrac{f\left( 5\right)  - f\left( 4\right) }{5 - 4}$ ，其几何意义为割线 ${AB}$ 的斜率，
            
            则有 ${f}^{\prime }\left( 5\right)  < f\left( 5\right)  - f\left( 4\right)  < {f}^{\prime }\left( 4\right)$ .
        \end{jieda}
        
        \item 设 $f\left( x\right)$ 为可导函数，且满足 $\mathop{\lim }\limits_{{{\Delta x} \rightarrow  0 }}\dfrac{f\left( {3 + {\Delta x}}\right)  - f\left( 3\right) }{3\Delta x} = 2$ ，则曲线 $y = f\left( x\right)$ 在点 $\left( {3,f\left( 3\right) }\right)$ 处的切线的斜率是\blkda{6}
        \begin{jieda}
            由题意, $\mathop{\lim }\limits_{{{\Delta x} \rightarrow  0 }}\dfrac{f\left( {3 + {\Delta x}}\right)  - f\left( 3\right) }{3\Delta x} = \dfrac{1}{3}\mathop{\lim }\limits_{{{\Delta x} \rightarrow  0 }}\dfrac{f\left( {3 + {\Delta x}}\right)  - f\left( 3\right) }{\Delta x} = 2$ ,
            
            即 $\dfrac{1}{3}{f}^{\prime }\left( 3\right)  = 2$ ,故 ${f}^{\prime }\left( 3\right)  = 6$ ,即曲线 $y = f\left( x\right)$ 在点 $\left( {3,f\left( 3\right) }\right)$ 处的切线的斜率是 6 .
        \end{jieda}

    \end{enumerate}
    \item 函数图象与导函数的关系
    \begin{enumerate}
        \item 函数 $f\left( x\right)$ 的图象如图所示,下列数值排序正确的是\dotfill{(\kh{D})}
        
        \begin{center}
        \includegraphics[width=0.2\textwidth]{bodyxb/src/bo_d5nlf73ef24c73bmm510_10_142_1485_349_326_0.jpg}
        \end{center}

        \xx{${f}^{\prime }\left( 1\right)  > {f}^{\prime }\left( 2\right)  > {f}^{\prime }\left( 3\right)  > 0$}{${f}^{\prime }\left( 1\right)  < {f}^{\prime }\left( 2\right)  < {f}^{\prime }\left( 3\right)  < 0$}{$0 < {f}^{\prime }\left( 1\right)  < {f}^{\prime }\left( 2\right)  < {f}^{\prime }\left( 3\right)$}{$0 > {f}^{\prime }\left( 1\right)  > {f}^{\prime }\left( 2\right)  > {f}^{\prime }\left( 3\right)$}
        
        \begin{jieda}
            作曲线在点 $\left( {1,f\left( 1\right) }\right) ,\left( {2,f\left( 2\right) }\right) ,\left( {3,f\left( 3\right) }\right)$ 处的切线,记其斜率依次为 ${k}_{1},{k}_{2},{k}_{3}$ , 结合图象可得 $0 > {k}_{1} > {k}_{2} > {k}_{3}$ 
                        
            \begin{center}
            \includegraphics[width=0.2\textwidth]{bodyxb/src/bo_d5nlf73ef24c73bmm510_11_143_674_344_317_0.jpg}
            \end{center}
            
            由导数的几何意义可得 ${k}_{1} = {f}^{\prime }\left( 1\right) ,{k}_{2} = {f}^{\prime }\left( 2\right) ,{k}_{3} = {f}^{\prime }\left( 3\right)$
            
            所以 $0 > {f}^{\prime }\left( 1\right)  > {f}^{\prime }\left( 2\right)  > {f}^{\prime }\left( 3\right)$ .
        \end{jieda}
        \item 已知函数 $f\left( x\right)$ 的部分图象如图所示, ${f}^{\prime }\left( x\right)$ 为 $f\left( x\right)$ 的导函数, 则\dotfill{(\kh{D})}
        
        \begin{center}
        \includegraphics[width=0.2\textwidth]{bodyxb/src/bo_d5nlf73ef24c73bmm510_8_140_1394_322_292_0.jpg}
        \end{center}

        \xx{$f\left( 1\right)  - f\left( 0\right)  > {f}^{\prime }\left( 1\right)  > {f}^{\prime }\left( 0\right)$}{${f}^{\prime }\left( 1\right)  > {f}^{\prime }\left( 0\right)  > f\left( 1\right)  - f\left( 0\right)$}{${f}^{\prime }\left( 0\right)  > f\left( 1\right)  - f\left( 0\right)  > {f}^{\prime }\left( 1\right)$}{${f}^{\prime }\left( 1\right)  > f\left( 1\right)  - f\left( 0\right)  > {f}^{\prime }\left( 0\right)$}

        \begin{jieda}
            由导数的几何意义可知, ${f}^{\prime }\left( 0\right)$ 表示曲线 $y = f\left( x\right)$ 在 $x = 0$ 处的切线斜率，${f}^{\prime }\left( 1\right)$ 表示曲线 $y = f\left( x\right)$ 在 $x = 1$ 处的切线斜率，
            
            $\dfrac{f\left( 1\right)  - f\left( 0\right) }{1 - 0}$ 表示 $\left( {0,f\left( 0\right) }\right) ,\left( {1,f\left( 1\right) }\right)$ 两点连线的斜率,
            
            由图可知，当 $x$ 从 0 变化到 1 时， $f\left( x\right)$ 切线斜率越来越大，
            
            所以 ${f}^{\prime }\left( 1\right)  > f\left( 1\right)  - f\left( 0\right)  > {f}^{\prime }\left( 0\right)$ ,对比选项可知, ${D}$ 正确.
        \end{jieda}
        
        \item 函数 $f\left( x\right)$ 的图象如图所示，下列数值排序正确的是\dotfill{(\kh{B})}
        
        \begin{center}
        \includegraphics[width=0.2\textwidth]{bodyxb/src/bo_d5nlf73ef24c73bmm510_9_138_577_303_330_0.jpg}
        \end{center}

        \xx{$0 < {f}^{\prime }\left( 1\right)  < {f}^{\prime }\left( 2\right)  < f\left( 2\right)  - f\left( 1\right)$}{$0 < {f}^{\prime }\left( 2\right)  < f\left( 2\right)  - f\left( 1\right)  < {f}^{\prime }\left( 1\right)$}{$0 < {f}^{\prime }\left( 2\right)  < {f}^{\prime }\left( 1\right)  < f\left( 2\right)  - f\left( 1\right)$}{$0 < f\left( 2\right)  - f\left( 1\right)  < {f}^{\prime }\left( 1\right)  < {f}^{\prime }\left( 2\right)$}
        \begin{jieda}
            如图,设函数 $f\left( x\right)$ 的图象上有两点 $A\left( {1,f\left( 1\right) }\right) ,B\left( {2,f\left( 2\right) }\right)$ ,经过点 $A,B$ 的切线分别为 ${l}_{A},{l}_{B}$
            \begin{center}
            \includegraphics[width=0.2\textwidth]{bodyxb/src/bo_d5nlf73ef24c73bmm510_9_140_1196_349_329_0.jpg}
            \end{center}
            则直线 ${AB},{l}_{A},{l}_{B}$ 的斜率依次为 ${k}_{AB} = \dfrac{f\left( 2\right)  - f\left( 1\right) }{2 - 1} = f\left( 2\right)  - f\left( 1\right) ,{k}_{{l}_{A}} = {f}^{\prime }\left( 1\right) ,{k}_{{l}_{B}} = {f}^{\prime }\left( 2\right)$ ,
            
            由图知直线 ${AB},{l}_{A},{l}_{B}$ 的倾斜角 ${\alpha }_{AB},{\alpha }_{{l}_{A}},{\alpha }_{{l}_{B}}$ 满足, $0 < {\alpha }_{{l}_{B}} < {\alpha }_{AB} < {\alpha }_{{l}_{A}} < \dfrac{{\pi }}{2}$ ,
            
            因函数 $y = \tan x$ 在 $\left( {0,\dfrac{\pi }{2}}\right)$ 上递增，故 $0 < \tan {\alpha }_{{l}_{B}} < \tan {\alpha }_{AB} < \tan {\alpha }_{{l}_{A}}$ ,
            
            即 $0 < {f}^{\prime }\left( 2\right)  < f\left( 2\right)  - f\left( 1\right)  < {f}^{\prime }\left( 1\right)$ .
        \end{jieda}
        
        
        \item 如图是函数 $f\left( x\right)$ 的部分图象，记 $f\left( x\right)$ 的导数为 ${f}^{\prime }\left( x\right)$ ，则下列选项中值最大的是\dotfill{(\kh{A})}
        
        \begin{center}
        \includegraphics[width=0.3\textwidth]{bodyxb/src/bo_d5nlf73ef24c73bmm510_10_143_271_358_241_0.jpg}
        \end{center}

        \xx{$f\left( 3\right)$}{$3{f}^{\prime }\left( 3\right)$}{$f\left( {-{14}}\right)$}{${f}^{\prime }\left( 8\right)$}
        \begin{jieda}
            由图可知, $f\left( {-{14}}\right) ,{f}^{\prime }\left( 8\right)$ 为负数, $f\left( 3\right) ,3{f}^{\prime }\left( 3\right)$ 为正数,故不选 $f\left( {-{14}}\right) ,{f}^{\prime }\left( 8\right)$ 
            
            \begin{center}
            \includegraphics[width=0.3\textwidth]{bodyxb/src/bo_d5nlf73ef24c73bmm510_10_143_756_358_244_0.jpg}
            \end{center}
            
            设 $f\left( x\right)$ 在 $x = 3$ 处的点为 $A$ ，显然 ${OA}$ 的斜率 ${k}_{OA}$ 大于 ${f}^{\prime }\left( 3\right)$ ，
            
            则 $\dfrac{f\left( 3\right)  - 0}{3 - 0} > {f}^{\prime }\left( 3\right)$ ，可转化为 $f\left( 3\right)  > 3{f}^{\prime }\left( 3\right)$ ，
            
            所以 $f\left( 3\right)$ 的值最大.
        \end{jieda}
    \end{enumerate}
    
    \item 基本初等函数的导数

    \begin{enumerate}
    \item 下列求导运算结果错误的是\dotfill{(\kh{A})}
    \xx{${\left( \dfrac{1}{x}\right) }^{\prime } = \dfrac{1}{{x}^{2}}$}{${\left( \ln x\right) }^{\prime } = \dfrac{1}{x}$}{${\left( {{e}}^{x}\right) }^{\prime } = {{e}}^{x}$}{${\left( \sin x\right) }^{\prime } = \cos x$}

    \begin{jieda}
        对于 ${A},{\left( \dfrac{1}{x}\right) }^{\prime } =  - \dfrac{1}{{x}^{2}}$ ,故 ${A}$ 错误;
        
        对于 ${B},{\left( \ln x\right) }^{\prime } = \dfrac{1}{x}$ ,故 ${B}$ 正确;
        
        对于 ${C},{\left( {{e}}^{x}\right) }^{\prime } = {{e}}^{x}$ ,故 ${C}$ 正确;
        
        对于 ${D},{\left( \sin x\right) }^{\prime } = \cos x$ ,故 ${D}$ 正确.
    \end{jieda}
    
    \item 下列导数运算正确的是\dotfill{(\kh{C})}
    \xx{${\left( \cos 3\right) }^{\prime } =  - \sin 3$}{${\left( {e}^{3}\right) }^{\prime } = 3{e}^{3}$}{${\left( \dfrac{1}{x}\right) }^{\prime } =  - \dfrac{1}{{x}^{2}}$}{${\left( {\log }_{2}x\right) }^{\prime } = \dfrac{1}{x}$}

    \begin{jieda}
        因为 ${\left( \cos 3\right) }^{\prime } = 0,{\left( {{e}}^{3}\right) }^{\prime } = 0,{\left( \dfrac{1}{x}\right) }^{\prime } =  - \dfrac{1}{{x}^{2}},{\left( {\log }_{2}x\right) }^{\prime } = \dfrac{1}{x\ln 2}$ .
    \end{jieda}
    
    \item 下列求导运算正确的是\dotfill{(\kh{D})}
    \xx{${\left( \ln 2\right) }^{\prime } = \dfrac{1}{2}$}{${\left( {e}^{3}\right) }^{\prime } = {3e}$}{${\left( {a}^{x}\right) }^{\prime } = x \cdot  {a}^{x - 1}$}{${\left( \cos x\right) }^{\prime } =  - \sin x$}
    \begin{jieda}
        对于 ${A} : {\left( \ln 2\right) }^{\prime } = 0$ ,故 ${A}$ 错误;
        
        对于 ${B} : {\left( {{e}}^{3}\right) }^{\prime } = 0$ ,故 ${B}$ 错误;
        
        对于 ${C} : {\left( {a}^{x}\right) }^{\prime } = {a}^{x} \cdot  \ln a$ ,故 ${C}$ 错误;
        
        对于 ${D} : {\left( \cos x\right) }^{\prime } =  - \sin x$ ，故 ${D}$ 正确
    \end{jieda}
    
    \end{enumerate}

    \item 导数的四则运算法则
    \begin{enumerate}
        \item 下列导数运算正确的是\dotfill{(\kh{B})}
        \xx{${\left( \sin x\right) }^{\prime } =  - \cos x$}{${\left( \dfrac{{{e}}^{x}}{x}\right) }^{\prime } = \dfrac{x{{e}}^{x} - {{e}}^{x}}{{x}^{2}}$}{${\left( {2}^{x}\right) }^{\prime } = {2}^{x}$}{${\left( \dfrac{1}{\sqrt{x}}\right) }^{\prime } =  - \dfrac{1}{2\sqrt{x}}$}
        \begin{jieda}
            对于 ${A},{\left( \sin x\right) }^{\prime } = \cos x$ ,故 ${A}$ 错误;
            
            对于 ${B},{\left( \dfrac{{{e}}^{x}}{x}\right) }^{\prime } = \dfrac{x{{e}}^{x} - {{e}}^{x}}{{x}^{2}}$ ，故 ${B}$ 正确；
            
            对于 ${C},{\left( {2}^{x}\right) }^{\prime } = {2}^{x} \cdot  \ln 2$ ,故 ${C}$ 错误;
            
            对于 ${D},{\left( \dfrac{1}{\sqrt{x}}\right) }^{\prime } = {\left( {x}^{-\frac{1}{2}}\right) }^{\prime } =  - \dfrac{1}{2}{x}^{-\frac{3}{2}}$ ,故 ${D}$ 错误.

        \end{jieda}

        \item 下列求导运算正确的是\dotfill{(\kh{D})}
        \xx{${\left( \cos x\right) }^{\prime } = \sin x$}{${\left( {x}^{2} + \ln 3\right) }^{\prime } = {2x} + \dfrac{1}{3}$}{${\left( {x}^{3}{{e}}^{x}\right) }^{\prime } = \sin x$}{${\left( \sqrt{x}\right) }^{\prime } = \dfrac{1}{2\sqrt{x}}$}

        \begin{jieda}
            对于 ${A},{\left( \cos x\right) }^{\prime } =  - \sin x$ ,所以 ${A}$ 错误,
            
            对于 ${B},\left( {{x}^{2} + \ln 3}\right)  = {2x}$ ,所以 ${B}$ 错误,
            
            对于 ${C},{\left( {x}^{3}{{e}}^{x}\right) }^{\prime } = 3{x}^{2}{{e}}^{x} + {x}^{3}{{e}}^{x}$ ,所以 ${C}$ 错误,
            
            对于 ${D},{\left( \sqrt{x}\right) }^{\prime } = \dfrac{1}{2}{x}^{-\frac{1}{2}} = \dfrac{1}{2\sqrt{x}}$ ，所以 ${D}$ 正确.
        \end{jieda}
        
        \item 已知函数 $f\left( x\right)  = {{e}}^{x} - f'\left( 1\right) x$ ，则\dotfill{(\kh{C})}
        \xx{$f\left( 1\right)  =  - \dfrac{{e}}{2}$}{$f'\left( 1\right)  =  - \dfrac{{e}}{2}$}{$f\left( 2\right)  = {{e}}^{2} - {e}$}{$f'\left( 2\right)  = {{e}}^{2} - {e}$}
        \begin{jieda}
            因为 $f\left( x\right)  = {{e}}^{x} - f'\left( 1\right) x$ ,所以 $f'\left( x\right)  = {{e}}^{x} - f'\left( 1\right)$，则 $f'\left( 1\right)  = {e} - f'\left( 1\right)$ ，所以 $f'\left( 1\right)  = \dfrac{{e}}{2}$
            
            则 $f\left( x\right)  = {{e}}^{x} - \dfrac{{e}}{2}x$ ,所以 $f\left( 1\right)  = \dfrac{{e}}{2},f\left( 2\right)  = {{e}}^{2} - \dfrac{{e}}{2},f'\left( 2\right)  = {{e}}^{2} - {e}$ .
        \end{jieda}
        
        \item 已知函数 $f\left( x\right)  = {{e}}^{x - 1} + f'\left( 1\right) {x}^{2} + f\left( 1\right) x + 1$ ,则 $f'\left( 2\right)  =$\blkda{$e-7$}
        \begin{jieda}
            令 $x = 1$ ,得 $f\left( 1\right)  = 1 + f'\left( 1\right)  + f\left( 1\right)  + 1$ ,解得 $f'\left( 1\right)  =  - 2$ ,
            
            所以 $f'\left( x\right)  = {{e}}^{x - 1} - {4x} + f\left( 1\right)$ ,
            
            令 $x = 1$ ,得 $f'\left( 1\right)  = 1 - 4 + f\left( 1\right)  =  - 2$ ,解得 $f\left( 1\right)  = 1$ ,
            
            所以 $f\left( x\right)  = {{e}}^{x - 1} - 2{x}^{2} + x + 1$ ,所以 $f'\left( x\right)  = {{e}}^{x - 1} - {4x} + 1$ ,所以 $f'\left( 2\right)  = {e} - 7$ .
        \end{jieda}
        
        \item 已知函数 $f\left( x\right)  = x\left( {x - 1}\right) \left( {x - 2}\right) \left( {x - 3}\right) \left( {x - 4}\right) \left( {x - 5}\right)$ ,求 $f'\left( 2\right)  =$\blkda{$-12$}
        \begin{jieda}
            令 $g\left( x\right)  = x\left( {x - 1}\right) \left( {x - 3}\right) \left( {x - 4}\right) \left( {x - 5}\right)$ ,则 $f\left( x\right)  = \left( {x - 2}\right) g\left( x\right)$ ,两边求导得到$f'\left( x\right)  = g\left( x\right)  + \left( {x - 2}\right) {g}^{\prime }\left( x\right)$ ,令 $x = 2$ ,得到 $f'\left( 2\right)  = g\left( 2\right)  =  - {12}$ .
        \end{jieda}
    \end{enumerate}
    \item 复合函数的求导方法
    \begin{enumerate}
        \item 下列求导运算不正确的是\dotfill{(\kh{A})}
        \xx{${\left( {x}^{2}\sin x + 1\right) }^{\prime } = {2x}\cos x$}{${\left( \sqrt{{2x} - 1}\right) }^{\prime } = \dfrac{1}{\sqrt{{2x} - 1}}$}{${\left\lbrack  \ln \left( 1 - 2x\right) \right\rbrack  }^{\prime } = \dfrac{2}{{2x} - 1}$}{${\left( \dfrac{x}{{{e}}^{x}}\right) }^{\prime } = \dfrac{1 - x}{{{e}}^{x}}$}
        \begin{jieda}
            对于 ${A} : {\left( {x}^{2}\sin x + 1\right) }^{\prime } = {2x}\sin x + {x}^{2}\cos x$ , ${A}$ 错误;
            
            对于 ${B}$ ,令 ${\left( \sqrt{{2x} - 1}\right) }^{\prime } = \dfrac{1}{2}{\left( 2x - 1\right) }^{-\frac{1}{2}} \cdot  2 = \dfrac{1}{\sqrt{{2x} - 1}}$ , B 正确;
            
            对于 ${C} : {\left\lbrack  \ln \left( 1 - 2x\right) \right\rbrack  }^{\prime } = \dfrac{1}{1 - {2x}} \cdot  \left( {-2}\right)  = \dfrac{2}{{2x} - 1},{C}$ 正确;
            
            对于 ${D} : {\left( \dfrac{x}{{{e}}^{x}}\right) }^{\prime } = \dfrac{{{e}}^{x} - x{{e}}^{x}}{{\left( {{e}}^{x}\right) }^{2}} = \dfrac{1 - x}{{{e}}^{x}}$ , ${D}$ 正确.
        \end{jieda}
        
        \item 下列求导不正确的是\dotfill{(\kh{C})}
        \xx{${\left\lbrack  {\left( 3x + 5\right) }^{3}\right\rbrack  }^{\prime } = 9{\left( 3x + 5\right) }^{2}$}{${\left( {{e}}^{2}\right) }^{\prime } = 0$}{${\left( \dfrac{2\sin x}{{x}^{2}}\right) }^{\prime } = \dfrac{{2x}\cos x + 4\sin x}{{x}^{3}}$}{${\left( {2}^{x} + \cos x\right) }^{\prime } = {2}^{x}\ln 2 - \sin x$}
        \begin{jieda}
            对于 ${A},{\left( {\left( 3x + 5\right) }^{3}\right) }^{\prime } = {\left( 3x + 5\right) }^{\prime } \cdot  3{\left( 3x + 5\right) }^{2} = 9{\left( 3x + 5\right) }^{2}$ ,故 ${A}$ 正确；
            
            对于 ${B},{{e}}^{2}$ 是常数,导数为 0,故 ${B}$ 正确;
            
            对于 ${C},{\left( \dfrac{2\sin x}{{x}^{2}}\right) }^{\prime } = \dfrac{{\left( 2\sin x\right) }^{\prime }{x}^{2} - 2\sin x \cdot  {\left( {x}^{2}\right) }^{\prime }}{{x}^{4}} = \dfrac{{2x}\cos x - 4\sin x}{{x}^{3}}$ ,故 ${C}$ 错误;
            
            对于 ${D},{\left( {2}^{x} + \cos x\right) }^{\prime } = {\left( {2}^{x}\right) }^{\prime } + {\left( \cos x\right) }^{\prime } = {2}^{x}\ln 2 - \sin x$ ,故 ${D}$ 正确.
        \end{jieda}
        
        \item 设 ${f}_{0}\left( x\right)  = \sin {2x} + \cos {2x},{f}_{1}\left( x\right)  = {f}_{0}^{\prime }\left( x\right) ,{f}_{2}\left( x\right)  = {f}_{1}^{\prime }\left( x\right) ,\cdots  {f}_{n + 1}\left( x\right)  = {f}_{n}^{\prime }\left( x\right) ,n \in  {N},$ 则 ${f}_{2020}\left( x\right)  =$\dotfill{(\kh{C})}
        \xx{${2}^{2020}\left( {\cos {2x} - \sin {2x}}\right)$}{ ${2}^{2021}\left( {-\sin {2x} - \cos {2x}}\right)$}{${2}^{2020}\left( {\sin {2x} + \cos {2x}}\right)$}{${2}^{2021}\left( {-\cos {2x} + \sin {2x}}\right)$}
        \begin{jieda}
            因为 ${f}_{0}\left( x\right)  = \sin {2x} + \cos {2x}$ ,
            
            所以 ${f}_{1}\left( x\right)  = {f}_{0}\left( x\right)  = 2\cos {2x} - 2\sin {2x} = 2\left( {\cos {2x} - \sin {2x}}\right)$ ,
            
            ${f}_{2}\left( x\right)  = {f}_{1}\left( x\right)  = {2}^{2}\left( {-\sin {2x} - \cos {2x}}\right)$ ,
            
            ${f}_{3}\left( x\right)  = {f}_{2}\left( x\right)  = {2}^{3}\left( {-\cos {2x} + \sin {2x}}\right)$ ,
            
            ${f}_{4}\left( x\right)  = {f}_{3}\left( x\right)  = {2}^{4}\left( {\sin {2x} + \cos {2x}}\right) ,$
            
            ...
            
            因为 ${2020} = {505} \times  4$ ,
            
            所以 ${f}_{2020}\left( x\right)  = {2}^{2020}\left( {\sin {2x} + \cos {2x}}\right)$ .
        \end{jieda}
        
        \item 设定义在 $R$ 上的函数 $f\left( x\right)$ 的导函数为 $f'\left( x\right)$ ，且 $f\left( {x + 1}\right)  = {2f}\left( x\right) ,f\left( x\right)  > 0$ ，则 $\dfrac{f'\left( 5\right) }{f'\left( 0\right) } =$ \blkda{$32$}

        \begin{jieda}
            $f\left( {x + 1}\right)  = {2f}\left( x\right)$ 两边对 $x$ 求导,得 $f'\left( {x + 1}\right)  \cdot  {\left( x + 1\right) }^{\prime } = {2f}'\left( x\right)$ ,即 $\dfrac{f'\left( {x + 1}\right) }{f'\left( x\right) } = 2$ 
            
            所以 $\dfrac{f'\left( 1\right) }{f'\left( 0\right) } = 2,\dfrac{f'\left( 2\right) }{f'\left( 1\right) } = 2,\cdots ,\dfrac{f'\left( 5\right) }{f'\left( 4\right) } = 2$ ,累乘可得 $\dfrac{f'\left( 5\right) }{f'\left( 0\right) } = {2}^{5} = {32}$ .
        \end{jieda}
    \end{enumerate}

    \item 函数图象的判断及应用
    \begin{enumerate}
        \item 函数 $f\left( x\right)  = x\sin x + \cos x$ 的导数 $f'\left( x\right)$ 的部分图象大致为\dotfill{(\kh{D})}
        \xx{\includegraphics[width=0.35\textwidth]{bodyxb/src/bo_d5o6oibef24c73bmtjs0_7_247_1844_285_275_0.jpg}}
        {\includegraphics[width=0.35\textwidth]{bodyxb/src/bo_d5o6oibef24c73bmtjs0_7_763_1839_295_279_0.jpg}}
        {\includegraphics[width=0.35\textwidth]{bodyxb/src/bo_d5o6oibef24c73bmtjs0_8_243_210_286_272_0.jpg}}
        {\includegraphics[width=0.35\textwidth]{bodyxb/src/bo_d5o6oibef24c73bmtjs0_8_767_206_292_276_0.jpg}}
        \begin{jieda}
            因为 $f\left( x\right)  = x\sin x + \cos x$ ,所以 $f'\left( x\right)  = \sin x + x\cos x - \sin x = x\cos x$ ,
            
            令 $g\left( x\right)  = f'\left( x\right)  = x\cos x,x \in  {R}$ ，则 $g\left( {-x}\right)  =  - x\cos x =  - g\left( x\right)$ ，
            
            所以函数 $g\left( x\right)  = x\cos x$ 是奇函数，故 A，C 错误；
            
            又 $g\left( {\pi }\right)  = {\pi }\cos {\pi } =  - {\pi } < 0$ ,故 ${B}$ 错误.
        \end{jieda}

        \item 若函数 $f\left( x\right)  = a{x}^{2} + {bx} + c\left( {a > 0}\right)$ 的图象的顶点在第三象限，则函数 $f'\left( x\right)$ 的图象大致是\dotfill{(\kh{C})}
        \xx{\includegraphics[width=0.35\textwidth]{bodyxb/src/bo_d5o6oibef24c73bmtjs0_8_243_1106_270_256_0.jpg}}
        {\includegraphics[width=0.35\textwidth]{bodyxb/src/bo_d5o6oibef24c73bmtjs0_8_764_1103_276_259_0.jpg}}
        {\includegraphics[width=0.35\textwidth]{bodyxb/src/bo_d5o6oibef24c73bmtjs0_8_241_1373_270_256_0.jpg}}
        {\includegraphics[width=0.35\textwidth]{bodyxb/src/bo_d5o6oibef24c73bmtjs0_8_766_1369_274_258_0.jpg}}
        \begin{jieda}
            因为函数 $f\left( x\right)  = a{x}^{2} + {bx} + c\left( {a > 0}\right)$ 的图象的顶点在第三象限，则 $\left\{  \begin{matrix} a > 0 \\   - \dfrac{b}{2a} < 0 \\  \dfrac{{4ac} - {b}^{2}}{4a} < 0 \end{matrix}\right.$ ，所以 $\left\{  \begin{array}{l} a > 0 \\  b > 0 \end{array}\right.$ ，又 $f\left( x\right)  = {2ax} + b$ ，则 $f\left( x\right)$ 的图象单调递增，且与 $y$ 轴交于正半轴.
        \end{jieda}
        
        \item 设函数 $f\left( x\right)  =  - \cos x - {x}^{4}$ 的导函数为 $g\left( x\right)$ ,则 $\left| {g\left( x\right) }\right|$ 图象大致是\dotfill{(\kh{D})}
        \xx{\includegraphics[width=0.35\textwidth]{bodyxb/src/bo_d5o6oibef24c73bmtjs0_9_254_284_215_178_0.jpg}}{\includegraphics[width=0.35\textwidth]{bodyxb/src/bo_d5o6oibef24c73bmtjs0_9_778_278_208_183_0.jpg}}{\includegraphics[width=0.35\textwidth]{bodyxb/src/bo_d5o6oibef24c73bmtjs0_9_253_506_233_175_0.jpg}}{\includegraphics[width=0.35\textwidth]{bodyxb/src/bo_d5o6oibef24c73bmtjs0_9_770_500_261_185_0.jpg}}
        
        \begin{jieda}
            因为 $f\left( x\right)  =  - \cos x - {x}^{4}$ ,所以 $f'\left( x\right)  = \sin x - 4{x}^{3}$ ,所以 $g\left( x\right)  = \sin x - 4{x}^{3}$ ,
            
            所以函数 $g\left( x\right)$ 是奇函数，其图象关于原点成中心对称，而函数 $\left| {g\left( x\right) }\right|$ 为偶函数，其图象关于 $y$ 轴对称，所以选项 $B,C$ 错误; 又因为其图象过原点 $O$ ,所以舍去 $A$
        \end{jieda}
        
        \item 下列四个图象中，有一个图象是函数 $f\left( x\right)  = \dfrac{1}{3}{x}^{3} - a{x}^{2} + \left( {{a}^{2} - 4}\right) x + 8\left( {a \neq  0}\right)$ 的导数的图象，则 $f\left( {-2}\right)$ 的值为\blkda{$-\dfrac{8}{3}$}
        \fourtu{\includegraphics[width=0.4\textwidth]{bodyxb/src/bo_d5o6oibef24c73bmtjs0_9_142_1210_222_258_0.jpg}}{\includegraphics[width=0.4\textwidth]{bodyxb/src/bo_d5o6oibef24c73bmtjs0_9_392_1210_222_258_0.jpg}}{\includegraphics[width=0.4\textwidth]{bodyxb/src/bo_d5o6oibef24c73bmtjs0_9_649_1210_222_258_0.jpg}}{\includegraphics[width=0.4\textwidth]{bodyxb/src/bo_d5o6oibef24c73bmtjs0_9_922_1210_222_258_0.jpg}}
        \begin{jieda}
            函数 $f\left( x\right)  = \dfrac{1}{3}{x}^{3} - a{x}^{2} + \left( {{a}^{2} - 4}\right) x + 8$ ,求导得 $f\left( x\right)  = {x}^{2} - {2ax} + {a}^{2} - 4 = {\left( x - a\right) }^{2} - 4$，于是函数 $y = f\left( x\right)$ 的图象是开口向上,对称轴为 $x = a$ 的抛物线, \ding{172}\ding{173} 不满足,
            
            又 $a \neq  0$ ，即函数 $y = f\left( x\right)$ 的图象对称轴不是 $y$ 轴， \ding{175} 不满足，因此符合条件的是 \ding{174}，
            
            函数 $y = f\left( x\right)$ 的图象过原点,且 $a > 0$ ,显然 $f'\left( 0\right)  = 0$ ,从而 $a = 2$，$f\left( x\right)  = \dfrac{1}{3}{x}^{3} - 2{x}^{2} + 8$ ，所以 $f\left( {-2}\right)  = \dfrac{1}{3} \times  {\left( -2\right) }^{3} - 2 \times  {\left( -2\right) }^{2} + 8 =  - \dfrac{8}{3}$ .
        \end{jieda}
    \end{enumerate}
    
    \item 求在曲线上一点的切线方程、过一点的曲线方程
    \begin{enumerate}
        \item 曲线 $y = \ln x$ 在点 $\left( {{e},1}\right)$ 处的切线方程为\blkda{$y = \dfrac{1}{{e}}x$}
        \begin{jieda}
            已知 $y = \ln x$ ,函数定义域为 $\left( {0, + \infty }\right)$，可得 ${y}^{\prime } = \dfrac{1}{x}$，此时 ${y}^{\prime }{|}_{x = {e}} = \dfrac{1}{{e}}$，所以曲线 $y = \ln x$ 在点 $\left( {{e},7}\right)$ 处的切线方程为 $y - 1 = \dfrac{1}{{e}}\left( {x - {e}}\right)$，即 $y = \dfrac{1}{{e}}x$
        \end{jieda}
        
        \item 已知曲线$y = \dfrac{x^3}{3}$，求过点$P\left(2, \dfrac{8}{3}\right)$的切线方程.
        \begin{jieda}
            设切点的坐标为 $\left( {{x}_{0},{y}_{0}}\right)$ ,切点在曲线上，在点 $\left( {{x}_{0},{y}_{0}}\right)$ 处切线的斜率为 $k = {y}^{\prime } = {x}_{0}^{2}$，切线方程为 $y - {y}_{0} = {x}_{0}^{2}\left( {x - {x}_{0}}\right)$ .
            
            又因为切线过点 $P\left( {2,\dfrac{8}{3}}\right)$ ,且切点 $\left( {{x}_{0},{y}_{0}}\right)$ 在曲线 $y = \dfrac{1}{3}{x}^{3}$ 上，故$\left\{  \begin{array}{l} \dfrac{8}{3} - {y}_{0} = {x}_{0}^{2}\left( {2 - {x}_{0}}\right) , \\  {y}_{0} = \dfrac{1}{3}{x}_{0}^{3}, \end{array}\right.$ 解得 ${x}_{0} = 2$ 或 ${x}_{0} =  - 1$ .
            \begin{dingautolist}{172}
                \item 当 ${x}_{0} = 2,{y}_{0} = \dfrac{8}{3}$ ,即切线斜率为 4 时,切线的方程为 ${12x} - {3y} - {16} = 0$
                \item 当 ${x}_{0} =  - 1,{y}_{0} =  - \dfrac{1}{3}$ ,即切线斜率为 1 时,切线的方程为 ${3x} - {3y} + 2 = 0$
            \end{dingautolist}
            综上,所求切线方程为 ${12x} - {3y} - {16} = 0$ 或 ${3x} - {3y} + 2 = 0$ .
        \end{jieda}
        
        \item 函数 $f\left( x\right)  = {x}^{2} + \ln x$ 在点 $\left( {1,1}\right)$ 处的切线方程是\blkda[3]{${3x} - y - 2 = 0$}
        \begin{jieda}
            由 $f\left( x\right)  = {x}^{2} + \ln x$ ,得到 ${f}^{\prime }\left( x\right)  = {2x} + \dfrac{1}{x}$ ,所以 ${f}^{\prime }\left( 1\right)  = 2 + 1 = 3$，所以 $f\left( x\right)  = {x}^{2} + \ln x$ 在点 $\left( {1,1}\right)$ 处的切线方程是 $y - 1 = 3\left( {x - 1}\right)$ ,即 ${3x} - y - 2 = 0$
        \end{jieda}
        
        \item 过点$(1,4)$且与曲线 $f\left( x\right)  = {x}^{3} + x + 2$ 相切的直线方程为\blkda[5]{${4x} - y = 0$ 或 ${7x} - {4y} + 9 = 0$}
        \begin{jieda}
            设过点 $\left( {1,4}\right)$ 的曲线 $y = f\left( x\right)$ 的切线为: $l : y - {y}_{0} = \left( {3{x}_{0}^{2} + 1}\right) \left( {x - {x}_{0}}\right)$ ,有 $\left\{  \begin{matrix} \left( {3{x}_{0}^{2} + 1}\right) \left( {1 - {x}_{0}}\right)  = 4 - {y}_{0} \\  {y}_{0} = {x}_{0}^{3} + {x}_{0} + 2 \end{matrix}\right.$，解得 $\left\{  \begin{array}{l} {x}_{0} = 1 \\  {y}_{0} = 4 \end{array}\right.$ 或 $\left\{  \begin{array}{l} {x}_{0} =  - \dfrac{1}{2} \\  {y}_{0} = \dfrac{9}{8} \end{array}\right.$，代入 $l$ 可得 ${4x} - y = 0$ 或 ${7x} - {4y} + 9 = 0$ .
        \end{jieda}
        
        \item 已知函数 $f\left( x\right)  = {x}^{2}{f}^{\prime }\left( \dfrac{1}{3}\right)  + \ln x - 9$ ，则函数在 $x = 1$ 处的切线方程是\blkda[3]{$y = {19x} - {19}$}
        \begin{jieda}
            ${f}^{\prime }\left( x\right)  = 2{f}^{\prime }\left( \dfrac{1}{3}\right) x + \dfrac{1}{x}$ ,令 $x = \dfrac{1}{3}$ ,可得 ${f}^{\prime }\left( \dfrac{1}{3}\right)  = 9$ ,
            
            $\therefore f\left( 1\right)  = 9 + 0 - 9 = 0,{f}^{\prime }\left( 1\right)  = {19}$ ,
            
            所以 $f\left( x\right)$ 在 $x = 1$ 处的切线方程为 $y = {19x} - {19}$ .
        \end{jieda}
        
        \item 曲线 $y = \dfrac{x}{{2x} - 1}$ 在点 $\left( {1,1}\right)$ 处的切线方程为\blkda[3]{$x + y - 2 = 0$}
        \begin{jieda}
            求导可得 ${y}^{\prime } = \dfrac{{2x} - 1 - {2x}}{{\left( 2x - 1\right) }^{2}} = \dfrac{-1}{{\left( 2x - 1\right) }^{2}}$，所以切线的斜率为 $k =  - 1$ ,
            
            所以切线方程为 $y - 1 =  - \left( {x - 1}\right)$ ,即 $x + y - 2 = 0$
        \end{jieda}
        
        \item 曲线 $f\left( x\right)  = \left( {x - 2}\right) {{e}}^{x} - {x}^{2} + {2x}$ 在 $x = 2$ 处的切线方程是\blkda[3]{$y = \left( {{{e}}^{2} - 2}\right) \left( {x - 2}\right)$}
        \begin{jieda}
            由函数 $f\left( x\right)  = \left( {x - 2}\right) {{e}}^{x} - {x}^{2} + {2x}$ ,可得 $f\left( x\right)  = \left( {x - 1}\right) {{e}}^{x} - {2x} + 2$ ,
            
            则 $f\left( 2\right)  = 0$ 且 $f\left( 2\right)  = {{e}}^{2} - 2$ ，即切点为 $\left( {2,0}\right)$ ，切线的斜率为 ${{e}}^{2} - 2$ ，
            
            所以函数 $f\left( x\right)$ 在 $x = 2$ 处的切线方程为 $y = \left( {{{e}}^{2} - 2}\right) \left( {x - 2}\right)$ .
        \end{jieda}
        
        \item 若函数 $f\left( x\right)  = \ln x - {ax} + 1$ 的图象在 $x = 2$ 处的切线与 $y$ 轴垂直，则函数 $f\left( x\right)$ 的图象在 $x = 1$ 处的切线方程为\blkda[3]{$x - {2y} = 0$}
        \begin{jieda}
            因为 $f\left( x\right)  = \ln x - {ax} + 1,x > 0$ ,

            所以 $f\left( x\right)  = \dfrac{1}{x} - a$ ,
            
            又因为函数在 $x = 2$ 处的切线与 $y$ 轴垂直，
            
            所以 $f\left( 2\right)  = \dfrac{1}{2} - a = 0$ ,
            
            解得 $a = \dfrac{1}{2}$ ,
            
            所以 $f\left( x\right)  = \ln x - \dfrac{1}{2}x + 1,x > 0$ ,
            
            所以 $f\left( 1\right)  = \dfrac{1}{2}$ ,
            
            即切点为 $\left( {1,\dfrac{1}{2}}\right)$ ,
            
            又因为 $f\left( x\right)  = \dfrac{1}{x} - \dfrac{1}{2}$ ，
            
            所以在 $x = 1$ 处的切线的斜率 $k = f\left( 1\right)  = \dfrac{1}{2}$ ,
            
            所以在 $x = 1$ 处的切线方程为: $y - \dfrac{1}{2} = \dfrac{1}{2}\left( {x - 1}\right)$，即 $y = \dfrac{1}{2}x,x - {2y} = 0$ .
        \end{jieda}
        
        \item 已知 $f\left( x\right)  = \ln \left( {-{2x}}\right)$ ，则 $f\left( x\right)$ 在点 $\left( {{x}_{0},f\left( {x}_{0}\right) }\right)$ 处切线的斜率为\blkda{$\dfrac{1}{{x}_{0}}$}
    \end{enumerate}
    
    \item 已知切线 (斜率) 求参数
    \begin{enumerate}
        \item 已知函数$f(x) = \dfrac{1}{2}x - \dfrac{1}{4}sinx-\dfrac{\sqrt{3}}{4}cosx$的图象在$A(x_0, y_0)$的切线斜率为$1$，则$tanx_0 = $\blkda{$-\sqrt{3}$}
        \begin{jieda}
            函数的导数$f'(x) = \dfrac{1}{2}-\dfrac{1}{4}cosx+\dfrac{\sqrt{3}}{4}sinx$，令$f'(x_0) = 1$，即$-\dfrac{1}{2}cosx_0 + \dfrac{\sqrt{3}}{2}sinx_0 = 1$，根据辅助角公式有$sin\left(x_0 - \dfrac{\pi}{6}\right) = 1$，因此$x_0 = 2k\pi + \dfrac{2\pi}{3}, k \in \mathbf{Z}$，所以$tanx_0 = -\sqrt{3}$
        \end{jieda}
        \item 已知函数$f(x) = lnx - \dfrac{a}{x} - 1$，若曲线$y = f(x)$存在斜率为$-1$的切线，求实数$a$的取值范围.
        \begin{jieda}
            $f'(x) = \dfrac{1}{x} + \dfrac{a}{x^2} (x > 0)$，令$f'(x) = -1$，即$x^2+x+a = 0$存在大于$0$的实数根. 易知$y = x^2+x+a$在$(0, +\infty)$单调递增，因此$a \in (-\infty, 0)$
        \end{jieda}
        \item 设函数$f(x) = ax+\dfrac{1}{x+b}(a, b \in \mathbf{Z})$，曲线$y = f(x)$在点$(2, f(2))$处的切线方程为$y = 3$
        \begin{enumerate}
            \item 求$y = f(x)$的表达式.
            \item 证明：曲线$y = f(x)$上任意一点的切线与直线$x = 1$和直线$y = x$所围成的三角形面积为定值，求此定值.
        \end{enumerate}
        \begin{jieda}
            \begin{enumerate}
                \item $\left \{ \begin{array}{l}
                     f(2) = 2a+\dfrac{1}{2+b} = 3 \\
                     f'(2) = a-\dfrac{1}{(2+b)^2} = 0
                \end{array} \right.$，解得$\left \{ \begin{array}{l}
                     a = 1 \\
                     b = -1
                \end{array} \right.$或$\left \{ \begin{array}{l}
                     a = \dfrac{9}{4} \\
                     b = -\dfrac{8}{3}
                \end{array} \right.$，又因为$a, b \in \mathbf{Z}$，所以$f(x) = x + \dfrac{1}{x-1}$
                \item $f'(x) = 1 - \dfrac{1}{(x-1)^2}$，故$\left(x_0, x_0+\dfrac{1}{x_0-1}\right)$处的切线为$y = \left(1 - \dfrac{1}{(x_0-1)^2}\right)(x-x_0)+\left(x_0+\dfrac{1}{x_0-1}\right)$ \\
                通过联立方程可知切线与$x = 1$与$y = x$的交点为$A\left(1, \dfrac{x_0+1}{x_0-1}\right)$和$B(2x_0-1, 2x_0-1)$，直线$x=1$与$y = x$的交点为$C(1,1)$，故$\vv{CA} = \left(0, \dfrac{2}{x_0-1}\right)$，$\vv{CB} = (2x_0-2, 2x_0 - 2)$，故三角形面积$S = \dfrac{1}{2}\left|(2x_0-2)\cdot\left(\dfrac{2}{x_0-1}\right) - 0 \cdot (2x_0-2)\right| = 2$
            \end{enumerate}
        \end{jieda}
        
        \item 若曲线 $y = f\left( x\right)$ 在 $x = 1$ 处的切线的斜率为 -3,则 $\mathop{\lim }\limits_{{{\Delta x} \rightarrow  0 }}\dfrac{f\left( 1\right)  - f\left( {1 + {2\Delta x}}\right) }{\Delta x} =$\blkda{6}
        \begin{jieda}
             $\mathop{\lim }\limits_{{{\Delta x} \rightarrow  0 }}\dfrac{f\left( 1\right)  - f\left( {1 + {2\Delta x}}\right) }{\Delta x} =  - 2\mathop{\lim }\limits_{{{\Delta x} \rightarrow  0 }}\dfrac{f\left( 1\right)  - f\left( {1 + {2\Delta x}}\right) }{-{2\Delta x}} =  - 2{f}^{\prime }\left( 1\right)  = 6$
        \end{jieda}
        
        \item 已知函数 $f\left( x\right)  = \left( {x - a}\right) \left( {x - 2}\right) \left( {x - 3}\right) \left( {x - 4}\right)$ ，若 $f\left( x\right)$ 的图象在 $x = 2$ 处的切线方程为 $y = {6x} + b$ ，则 $a + b =$\blkda{$- {13}$}
        \begin{jieda}
            由题意知 $f\left( 2\right)  = \left( {2 - a}\right)  \times  \left( {2 - 2}\right)  \times  \left( {2 - 3}\right)  \times  \left( {2 - 4}\right)  = 0$ ,
            
            所以 $0 = 6 \times  2 + b$ ,解得 $b =  - {12}$ ,
            
            又 ${f}^{\prime }\left( x\right)  = \left( {x - a}\right) \left( {x - 3}\right) \left( {x - 4}\right)  + \left( {x - 2}\right) {\left\lbrack  \left( x - a\right) \left( x - 3\right) \left( x - 4\right) \right\rbrack  }^{\prime }$ ,
            
            所以 ${f}^{\prime }\left( 2\right)  = \left( {2 - a}\right)  \times  \left( {2 - 3}\right)  \times  \left( {2 - 4}\right)  = 6$ ,解得 $a =  - 1$ ,所以 $a + b =  - 1 + \left( {-{12}}\right)  =  - {13}$ .
        \end{jieda}
        
        \item 已知函数 $f\left( x\right)$ 在点 $x =  - 1$ 处的切线方程为 $x + y - 1 = 0$ , 则 ${f}^{\prime }\left( {-1}\right)  + f\left( {-1}\right)  =$ \blkda{1}
        \begin{jieda}
            由切线方程 $x + y - 1 = 0$ ,得 ${f}^{\prime }\left( {-1}\right)  = k =  - 1$ ,
            
            将 $x =  - 1$ 代入切线方程 $x + y - 1 = 0$ ，得 $y = 2$ ，所以 $f\left( {-1}\right)  = 2$ ，
            
            则 ${f}^{\prime }\left( {-1}\right)  + f\left( {-1}\right)  =  - 1 + 2 = 1$ .
        \end{jieda}
        
        \item 已知直线 $y = {ax} - 1$ 与曲线 $y = \dfrac{\ln x}{x}$ 相切，则 $a$ 的值为\blkda{1}
        \begin{jieda}
            $y = \dfrac{\ln x}{x}$ 的导函数 ${y}^{\prime } = \dfrac{1 - \ln x}{{x}^{2}}$ ,设切点为 $\left( {{x}_{0},{y}_{0}}\right)$ ,则 $\left\{  \begin{matrix} {y}_{0} = a{x}_{0} - 1 \\  {y}_{0} = \dfrac{\ln {x}_{0}}{{x}_{0}} \\  a = \dfrac{1 - \ln {x}_{0}}{{x}_{0}^{2}} \end{matrix}\right.$ ,故 $\left\{  \begin{matrix} {y}_{0} = \dfrac{1 - \ln {x}_{0}}{{x}_{0}} - 1 \\  {y}_{0} = \dfrac{\ln {x}_{0}}{{x}_{0}} \end{matrix}\right.$ , 即 $\dfrac{\ln {x}_{0}}{{x}_{0}} = \dfrac{1 - \ln {x}_{0}}{{x}_{0}} - 1$ ，则 $2\ln {x}_{0} + {x}_{0} - 1 = 0$ .

            易得函数 $f\left( x\right)  = 2\ln x + x - 1$ 为增函数，且 $f\left( 1\right)  = 0$ ，故 ${x}_{0} = 1$ .故 $a = \dfrac{1 - \ln 1}{1} = 1$ .
        \end{jieda}
        
        \item 若直线 $y = \dfrac{1}{4}x + b$ 是曲线 $y = \ln x\left( {x > 0}\right)$ 的一条切线，则实数 $b$ 的值为\blkda{$ln4-1$}
        \begin{jieda}
            $\because y = \ln x$ 的导数 ${y}^{\prime } = \dfrac{1}{x}$ , $\therefore$ 令 $\dfrac{1}{x} = \dfrac{1}{4}$ ,得 $x = 4$ , $\therefore$ 切点为 $\left( {4,\ln 4}\right)$ .
            
            代入直线 $y = \dfrac{1}{4}x + b$ ，得 $\ln 4 = \dfrac{1}{4} \times  4 + b$ ，即 $b = \ln 4 - 1$ .
        \end{jieda}
        \item 若直线 $y = {2x}$ 是曲线 $f\left( x\right)  = x\left( {{{e}}^{2x} - a}\right)$ 的切线，则 $a =$\blkda{$-1$}
        \begin{jieda}
             $f\left( x\right)  = x\left( {{{e}}^{2x} - a}\right) ,f'\left( x\right)  = \left( {1 + {2x}}\right) {{e}}^{2x} - a$ ,
            
            依题意,直线 $y = {2x}$ 是曲线 $f\left( x\right)  = x\left( {{{e}}^{2x} - a}\right)$ 的切线,
            
            设切点为 $\left( {t,{2t}}\right)$ ,则 $\left\{  \begin{matrix} t\left( {{{e}}^{2t} - a}\right)  = {2t} \\  \left( {1 + {2t}}\right) {{e}}^{2t} - a = 2 \end{matrix}\right.$ , $\left\{  \begin{matrix} t{{e}}^{2t} = \left( {2 + a}\right) t \\  \left( {1 + {2t}}\right) {{e}}^{2t} = 2 + a \end{matrix}\right.$ ,
            
            通过对比系数可得 $\left( {1 + {2t}}\right) t = t,2{t}^{2} = 0,t = 0$ ,则 $a =  - 1$ .
        \end{jieda}
        
        \item 已知函数 $f\left( x\right)  = \dfrac{1}{x} + a\ln x$ ，在点 $\left( {1,f\left( 1\right) }\right)$ 处的切线方程为 $x - y = 0$ ,则 $a =$\blkda{2}
        \begin{jieda}
            因为 $f\left( x\right)  = \frac{1}{x} + a\ln x$ ,则 $f\left( x\right)  = \dfrac{a}{x} - \dfrac{1}{{x}^{2}}$
            
            因为函数 $f\left( x\right)$ 在点 $\left( {1,f\left( 1\right) }\right)$ 处的切线方程为 $x - y = 0$ ,则 $f\left( 1\right)  = a - 1 = 1$ ,解得 $a = 2$ ,
            
            且 $f\left( 1\right)  = 1$ ,则点 $\left( {1,f\left( 1\right) }\right)$ 在直线 $x - y = 0$ 上,合乎题意.
            
            因此, $a = 2$ .
        \end{jieda}
        
    \end{enumerate}
    \item 两条切线平行、垂直、重合 (公切线) 问题
    \begin{enumerate}
        \item 若曲线 $y = {x}^{2} - a\ln x$ 在点 $P\left( {1,1}\right)$ 处的切线与直线 $y = x - 2$ 垂直, 则实数 $a$ 的值为\blkda{3}
        \begin{jieda}
            由 $y = {x}^{2} - a\ln x$ ,求导 ${y}^{\prime } = {2x} - \dfrac{a}{x}$ ,
            
            则 $y = {x}^{2} - a\ln x$ 在点 $P\left( {1,1}\right)$ 处的切线的斜率为 ${\left. {y}^{\prime }\right| }_{x = 1} = 2 - a$ ,
            
            而 $y = {x}^{2} - a\ln x$ 在点 $P\left( {1,1}\right)$ 处的切线与直线 $y = x - 2$ 垂直，
            
            则 $2 - a =  - 1$ ，故 $a = 3$ .
        \end{jieda}
        
        \item 若直线 $x + y + a = 0$ 是曲线 $f\left( x\right)  = {x}^{3} + {bx} - {14}$ 与曲线 $g\left( x\right)  =  {x}^{2} - {3\ln }x$ 的公切线，则 $a - b =$\blkda{11}
        \begin{jieda}
            因为 $g\left( x\right)  = {x}^{2} - 3\ln x$，所以 ${g}^{\prime }\left( x\right)  = {2x} - \dfrac{3}{x}$ ,由 ${2x} - \dfrac{3}{x} =  - 1$ ,解得 $x = 1$ 或 $x =  - \dfrac{3}{2}$ (舍去)，所以切点为 $\left( {1,1}\right)$
            
            因为切点在切线 $x + y + a = 0$ 上,解得 $a =  - 2$，所以切线方程为 $x + y - 2 = 0$
            
            ${f}^{\prime }\left( x\right)  = 3{x}^{2} + b$ 设切点为 $\left( {t,{t}^{3} + {bt} - {14}}\right)$ ,
            
            由题意得 $\left\{  \begin{matrix} 3{t}^{2} + b =  - 1 \\  t + {t}^{3} + {bt} - {14} - 2 = 0 \end{matrix}\right.$ ,解得 $\left\{  \begin{matrix} b =  - {13} \\  t =  - 2 \end{matrix}\right.$ ,
            
            所以 $a - b = {11}$
        \end{jieda}

        
        \item 已知 $f\left( x\right)  = {x}^{3} + {nx} - {52},g\left( x\right)  = {x}^{2} - 3\ln x$ ，若直线 $x + y + m = 0$ 是曲线 $y = f\left( x\right)$ 与曲线 $y = g\left( x\right)$ 的公切线,则 $m - n =$\blkda{26}
        \begin{jieda}
            设直线 $x + y + m = 0$ 与曲线 $y = f\left( x\right)$ 相切于点 $\left( {a, - a - m}\right)$ ,与曲线 $y = g\left( x\right)$ 相切于点 $\left( {b, - b - m}\right) ,b > 0$ .
            
            由 $g\left( x\right)  = {x}^{2} - 3\ln x$ 知 ${g}^{\prime }\left( x\right)  = {2x} - \dfrac{3}{x}$ ,又两曲线的公切线斜率为 -1,则 ${2b} - \dfrac{3}{b} =  - 1$ ,解得 $b = 1$ 或 $b =  - \dfrac{3}{2}$ (舍去).
            
            所以 $1 - 3\ln 1 =  - 1 - m$ ,解得 $m =  - 2$ .
            
            由 $f\left( x\right)  = {x}^{3} + {nx} - {52}$ 知 $f\left( x\right)  = 3{x}^{2} + n$ ，又两曲线的公切线斜率为 -1，则 $3{a}^{2} + n =  - 1$ ，即 $n =  - 3{a}^{2} - 1$ ，故 ${a}^{3} - \left( {3{a}^{2} + 1}\right) a - {52} =  - a + 2$ ,整理得 ${a}^{3} =  - {27}$ ,故 $a =  - 3$ ,
            
            所以 $n =  - 3{a}^{2} - 1 =  - {28}$ ,故 $m - n = {26}$ .
        \end{jieda}
        
        \item 已知函数 $f\left( x\right)  =  - {x}^{3} + x + 1,g\left( x\right)  = {{e}}^{-{2x} + 1}$ .
        \begin{enumerate}
            \item 求曲线 $y = f\left( x\right)$ 在 $x = 1$ 处的切线方程
            \item 若曲线 $y = f\left( x\right)$ 在 $x = 1$ 处的切线与曲线 $y = g\left( x\right)$ 在 $x = t\left( {t \in  {R}}\right)$ 处的切线平行,求 $t$ 的值
        \end{enumerate}
        \begin{jieda}
            \begin{enumerate}
                \item ${f}^{\prime }\left( x\right)  =  - 3{x}^{2} + 1,f\left( 1\right)  = 1,{f}^{\prime }\left( 1\right)  =  - 2$ ,

                所以曲线 $y = f\left( x\right)$ 在 $x = 1$ 处的切线方程为 $y - 1 =  - 2\left( {x - 1}\right)$ ,即 ${2x} + y - 3 = 0$

                \item 由(1)可得:曲线 $y = f\left( x\right)$ 在点 $\left( {1,1}\right)$ 处的切线方程为 $y =  - {2x} + 3$ 
                
                由 ${g}^{\prime }\left( x\right)  =  - 2{{e}}^{-{2x} + 1}$ ,可得曲线 $y = g\left( x\right)$ 在 $x = t\left( {t \in  {R}}\right)$ 处的切线斜率为 ${g}^{\prime }\left( t\right)  =  - 2{{e}}^{-{2t} + 1}$，由题意可得 $- 2{{e}}^{-{2t} + 1} =  - 2$ ，从而 $t = \dfrac{1}{2}$
            \end{enumerate}
        \end{jieda}
        
        \item 若曲线 $y = \sqrt{x}$ 在点 $\left( {4,2}\right)$ 处的切线也是曲线 $y = a{{e}}^{x}$ 的切线,则 $a =$\blkda{$\dfrac{{{e}}^{3}}{4}$}
        \begin{jieda}
            $y = \sqrt{x}$ 的导数 ${y}^{\prime } = \dfrac{1}{2\sqrt{x}}$ ,令 $x = 4$ ,则 ${y}^{\prime } = \dfrac{1}{4}$ ,

            所以曲线 $y = \sqrt{x}$ 在 $\left( {4,2}\right)$ 处的切线方程为 $y - 2 = \dfrac{1}{4}\left( {x - 4}\right)$，即 $x - {4y} + 4 = 0$ .
            
            $y = a{{e}}^{x}$ 的导数 ${y}^{\prime } = a{{e}}^{x}$ ,设直线 $x - {4y} + 4 = 0$ 与曲线 $y = a{{e}}^{x}$ 切于点 $\left( {{x}_{0},a{{e}}^{{x}_{0}}}\right)$ ,
            
            则曲线 $y = a{{e}}^{x}$ 在点 $\left( {{x}_{0},a{{e}}^{{x}_{0}}}\right)$ 处的切线方程为 $y - a{{e}}^{{x}_{0}} = a{{e}}^{{x}_{0}}\left( {x - {x}_{0}}\right)$ ,
            
            即 $x - \dfrac{1}{a{{e}}^{{x}_{0}}}y + 1 - {x}_{0} = 0$ ,所以 $\left\{  \begin{array}{l} \dfrac{1}{a{{e}}^{{x}_{0}}} = 4, \\  1 - {x}_{0} = 4, \end{array}\right.$ 解得 $a = \dfrac{{{e}}^{3}}{4}$ .
        \end{jieda}
        
        \item 已知函数 $f\left( x\right)  = {e}^{x - a},g\left( x\right)  = \ln x - b$ .
        \begin{enumerate}
            \item 当 $a = 1$ 时,求曲线 $y = f\left( x\right)$ 在点 $\left( {1,1}\right)$ 处的切线方程
            \item 若 $a = b + 2$ ，是否存在直线与曲线 $y = f\left( x\right)$ 和 $y = g\left( x\right)$ 都相切？若存在，求出所有这样的直线；若不存在, 请说明理由.
        \end{enumerate}
        
        \begin{jieda}
            \begin{enumerate}
                \item 当 $a = 1$ 时， ${f}^{\prime }\left( x\right)  = {e}^{x - 1},f\left( 1\right)  = 1,{f}^{\prime }\left( 1\right)  = 1$
                
                曲线 $y = f\left( x\right)$ 在点 $\left( {1,1}\right)$ 处的切线方程为: $y - f\left( 1\right)  = {f}^{\prime }\left( 1\right) \left( {x - 1}\right)$ ,
                
                代入整理得: $y = x$ .
                \item 设直线与曲线 $y = f\left( x\right)$ 相切于点 $A\left( {{x}_{1},{y}_{1}}\right)$ ,与曲线 $y = g\left( x\right)$ 相切于点 $B\left( {{x}_{2},{y}_{2}}\right) ,{f}^{\prime }\left( x\right)  = {e}^{x - a}$ , ${g}^{\prime }\left( x\right)  = \dfrac{1}{x}$
                
                曲线 $y = f\left( x\right)$ 在点 $A$ 处的切线为: $y - {e}^{{x}_{1} - a} = {e}^{{x}_{1} - a}\left( {x - {x}_{1}}\right)$
                
                与曲线 $y = g\left( x\right)$ 相切于点 $B$ ,则 $\left\{  \begin{matrix} {e}^{{x}_{1} - a} = \dfrac{1}{{x}_{2}} \\  \ln {x}_{2} - b - {e}^{{x}_{1} - a} = {e}^{{x}_{1} - a}\left( {{x}_{2} - {x}_{1}}\right) \end{matrix}\right.$
                
                由上式得: ${x}_{1} - a = \ln \left(\dfrac{1}{{x}_{2}}\right) =  - \ln {x}_{2}$ ，则 $\ln {x}_{2} = a - {x}_{1}$
                
                将 ${e}^{{x}_{1} - a} = \dfrac{1}{{x}_{2}}\text{ 、 }\ln {x}_{2} = a - {x}_{1}$ 代入下式得: $a - {x}_{1} - b - \dfrac{1}{{x}_{2}} = \dfrac{1}{{x}_{2}}\left( {{x}_{2} - {x}_{1}}\right)$ ，
                
                整理得: $\left( {{x}_{1} - 1}\right) \left( {{x}_{2} - 1}\right)  = 0$

                \begin{dingautolist}{172}
                    \item 当 ${x}_{1} = 1$ 时， $y - {e}^{1 - a} = {e}^{1 - a}\left( {x - 1}\right)$ ，即切线为 $y = {e}^{1 - a}x$
                    \item 当 ${x}_{2} = 1$ 时， $a - {x}_{1} = \ln {x}_{2} = 0,{x}_{1} = a$ ，因此切线为 $y - 1 = x - a$
                \end{dingautolist}
            \end{enumerate}
        \end{jieda}
    \end{enumerate}

    \item 曲线上两点间距离最值问题
    \begin{enumerate}
        \item 已知实数 ${a}$ 、 ${b}$ 、 ${c}$ 、 ${d}$ 满足 $\left| {{b} - \dfrac{{{lna}} }{{a}}}\right|  + \left| {{c} - {d} + 2}\right|  = 0$ ,则 ${\left( {a} - {c}\right) }^{2} + {\left( {b} - {d}\right) }^{2}$ 的最小值为\blkda{$\dfrac{9}{2}$}
        \begin{jieda}
            $\left| {{b} - \dfrac{ {{lna}} }{{a}}}\right|  + \left| {{c} - {d} + 2}\right|  = 0 \Longleftrightarrow \left\{\begin{array}{l}
                 b = \dfrac{lna}{a}  \\
                 d = c+2
            \end{array}\right.$

            ${\left( {a} - {c}\right) }^{2} + {\left( {b} - {d}\right) }^{2}$即$Q(a, b)$与$P(c, d)$两点距离平方，$Q$在$y = \dfrac{lnx}{x}$上，$P$在$y = x+2$上.

            \begin{center}
                \includegraphics[width=0.25\linewidth]{bodyxb//src/quxianshangliangdianjulizuizhi1.png}
            \end{center}

            对于任意的$Q$，过$Q$向直线做垂线，垂足取$P$取得最小值，$PQ$长度即$Q$到直线距离. 计算可知$x = 1$处$y = \dfrac{lnx}{x}$的切线与$y = x+2$平行，此时取得$Q$到直线距离的最小值，此时目标值为$\dfrac{9}{2}$
            
        \end{jieda}

        \item 对于函数 $y = h\left( x\right)$ ,记 ${h}^{\left( 0\right) }\left( x\right)  = h\left( x\right) ,{h}^{\left( 1\right) }\left( x\right)  = {\left( h\left( x\right) \right) }^{\prime },\cdots ,{h}^{\left( n + 1\right) }\left( x\right)  = {\left( {h}^{\left( n\right) }\left( x\right) \right) }^{\prime }\left( {n \in  \mathbf{N}}\right)$ . 如果 $n$ 是满足 ${h}^{\left( n\right) }\left( x\right)  = h\left( x\right)$ 的最小正整数,则称 $n$ 是函数 $y = h\left( x\right)$ 的“最小导周期”.
        \begin{enumerate}
            \item 已知 $f\left( x\right)  = a\sin \left( {x + t}\right)  + b\cos \left( {x + t}\right)$ ,证明：对任意 $a,b,t \in  {R},f\left( x\right)$ 的最小导周期为 4
            \item 设 $m,n \in  {R},g\left( x\right)  = {e}^{mx} + n\cos x$ ,若函数 $y = g\left( x\right)$ 的最小导周期为 2 ,记$M\left( {a,b}\right)  = \sqrt{{\left( a - b\right) }^{2} + {\left( a + 1 + g\left( b\right) \right) }^{2}}$ ,当实数 $a,b$ 变化时,求 $M\left( {a,b}\right)$ 的最小值
            \item 设 $\omega  > 1,h\left( x\right)  = \cos {\omega x}$ ,若函数 $y = h\left( x\right)$ 满足 ${h}^{\left( 2\right) }\left( x\right)  \leq  x$ 对 $x \in  \left( {0, + \infty }\right)$ 恒成立,且存在${x}_{0} \in  \left( {0, + \infty }\right)$ 使得 ${{h}}^{\left( 2\right) }\left( {x}_{0}\right)  = {x}_{0}$ ,试用 $\omega$ 表示 ${x}_{0}$ ,并证明 $\dfrac{\pi }{2\omega } < {x}_{0} < \dfrac{\pi }{\omega }$ .
        \end{enumerate}
        \begin{jieda}
            \begin{enumerate}
                \item 证明: 因为 ${f}^{\left( 1\right) }\left( x\right)  = a\cos \left( {x + t}\right)  - b\sin \left( {x + t}\right)$ ,${f}^{\left( 2\right) }\left( x\right)  =  - a\sin \left( {x + t}\right)  - b\cos \left( {x + t}\right) , {f}^{\left( 3\right) }\left( x\right)  =  - a\cos \left( {x + t}\right)  + b\sin \left( {x + t}\right) , {f}^{\left( 4\right) }\left( x\right)  = a\sin \left( {x + t}\right)  + b\cos \left( {x + t}\right)  = f\left( x\right)$
                
                所以,对任意实数 $a,b,t$ ,都有 ${f}^{\left( 4\right) }\left( x\right)  = f\left( x\right)$ ,即最小导周期为 4 .
                
                \item $g\left( x\right)  = {g}^{\left( 2\right) }\left( x\right)$ ,则 ${{e}}^{mx} + n\cos x = {m}^{2}{{e}}^{mx} - n\cos x$ 对任意实数 $x$ 恒成立
                
                令 $x = 0$ ,则 $1 + n = {m}^{2} - n$ ,即 ${m}^{2} = 1 + {2n}$ 
                
                令 $x = \dfrac{\pi }{2}$ ,则 ${{e}}^{\frac{m\pi }{2}} = {m}^{2}{{e}}^{\frac{m\pi }{2}}$，则 ${m}^{2} = 1$ 

                所以 $m = 1,n = 0$ (舍,由于 ${g}^{\left( 1\right) }\left( x\right)  = {{e}}^{x} = g\left( x\right)$ ,最小导周期不是 2 ) 或 $m =  - 1,n = 0$ .故 $g\left( x\right)  = {{e}}^{-x}$ .

                $M\left( {a,b}\right)  = \sqrt{{\left( a - b\right) }^{2} + {\left( a + 1 + g\left( b\right) \right) }^{2}} = \sqrt{{\left( a - b\right) }^{2} + {\left( -a - 1 - {{e}}^{-b}\right) }^{2}}$ 的最小值可视为点$P\left( {a, - a - 1}\right)$ 与点 $Q\left( {b,{{e}}^{-b}}\right)$ 之间的最短距离,即曲线 $y = {{e}}^{-x}$ 上的点到直线 $y =  - x - 1$距离的最小值.
                \begin{center}
                    \includegraphics[width=0.25\linewidth]{bodyxb//src/quxianshangliangdianjulizuizhi2.png}
                \end{center}
                设切点 $M\left( {{x}_{0},{{e}}^{-{x}_{0}}}\right) ,{y}^{\prime } =  - {{e}}^{-x}$ ,切线斜率 $- {{e}}^{-{x}_{0}} =  - 1$ ,得 ${x}_{0} = 0$ ,切点为(0,1)，点(0,1)到直线 $y =  - x - 1$ 距离 $d = \dfrac{\left| 1 + 0 + 1\right| }{\sqrt{{1}^{2} + {1}^{2}}} = \sqrt{2}$ . 即 $M\left( {a,b}\right)$ 的最小值为 $\sqrt{2}$ .

                \item $h^{(2)}(x) = -\omega^2cos\omega x$，设$\varphi(x) = -\omega^2cos\omega x - x$，则$x = x_0$是$\varphi(x)$的极大值点，$\varphi'(x) = \omega^3sin\omega x - 1$，故有$\left\{\begin{array}{l}
                     sin\omega x_0 = \dfrac{1}{\omega^3} \\
                     cos\omega x_0 = -\dfrac{x_0}{\omega^2}
                \end{array}\right.$，因此有$1 = \dfrac{1}{\omega^6} + \dfrac{x_0^2}{\omega^4}$，因此$x_0 = \dfrac{\sqrt{\omega^6-1}}{\omega}$

                $\left\{\begin{array}{l}
                     sin\omega x_0 > 0 \\
                     cos\omega x_0 < 0
                \end{array}\right.$，故$\omega x_0 \in \left(2k\pi +\dfrac{1}{2}\pi, 2k\pi+\pi\right), k \in \mathbf{Z}$，由$\omega x_0 > 0$可知$k \geq 0$.

                假设$k > 0$，则$\varphi\left(x_0 - \dfrac{2k\pi}{\omega}\right) = -\omega^2cos\omega \left(x_0 - \dfrac{2k\pi}{\omega}\right) - \left(x_0 - \dfrac{2k\pi}{\omega}\right) = -\omega^2cos\omega x_0 - \left(x_0 - \dfrac{2k\pi}{\omega}\right) = \dfrac{2k\pi}{\omega} > 0$，与$\varphi(x) \leq 0$矛盾.

                故$k = 0$，$\omega x_0 \in \left(\dfrac{1}{2}\pi, \pi\right)$，即$\dfrac{\pi }{2\omega } < {x}_{0} < \dfrac{\pi }{\omega }$.
            \end{enumerate}
        \end{jieda}
    \end{enumerate}

    \item 导数与新定义
    \begin{enumerate}
        \item 给出定义:若函数 $f\left( x\right)$ 在 $D$ 上可导，即 $f'\left( x\right)$ 存在，且导函数 $f'\left( x\right)$ 在 $D$ 上也可导,则称 $f\left( x\right)$ 在 $D$ 上存在二阶导数,记 $f''\left( x\right)  = {\left( f'\left( x\right) \right) }^{\prime }$ . 若 $f''\left( x\right)  < 0$ 在 $D$ 上恒成立,则称 $f\left( x\right)$ 在 $D$ 上为凸函数. 以下四个函数在 $\left( {0,\dfrac{{\pi }}{2}}\right)$ 上不是是凸函数的是\dotfill{(\kh{D})}
        \xx{$f\left( x\right)  = \sin x + \cos x$}{$f\left( x\right)  = \ln x - {2x}$}{$f\left( x\right)  =  - {x}^{3} + {2x} - 1$}{$f\left( x\right)  =  - x{{e}}^{-x}$}
        
        \begin{jieda}
            对于 ${A} : f\left( x\right)  = \sin x + \cos x,f'\left( x\right)  = \cos x - \sin x,f''\left( x\right)  =  - \sin x - \cos x$ ,
            
            则 $f\left( x\right)$ 在 $\left( {0,\dfrac{{\pi }}{2}}\right)$ 上恒有 ${f}''\left( x\right)  < 0$ ,故 ${A}$ 错误;
            
            对于 ${B} : f\left( x\right)  = \ln x - {2x},f'\left( x\right)  = \dfrac{1}{x} - 2,{f}''\left( x\right)  =  - \dfrac{1}{{x}^{2}}$ ,
            
            则 $f\left( x\right)$ 在 $\left( {0,\dfrac{\pi }{2}}\right)$ 上恒有 ${f}''\left( x\right)  < 0$ ,故 ${B}$ 错误;
            
            对于 ${C} : f\left( x\right)  =  - {x}^{3} + {2x} - 1,f'\left( x\right)  =  - 3{x}^{2} + 2,{f}''\left( x\right)  =  - {6x}$ ,
            
            则 $f\left( x\right)$ 在 $\left( {0,\dfrac{{\pi }}{2}}\right)$ 上恒有 ${f}''\left( x\right)  < 0$ ,故 ${C}$ 错误;
            
            对于 ${D} : f\left( x\right)  =  - x{{e}}^{-x},f'\left( x\right)  =  - {{e}}^{-x} + x{{e}}^{-x},f''\left( x\right)  = \left( {2 - x}\right) {{e}}^{-x}$ ,
            
            则 $f\left( x\right)$ 在 $\left( {0,\dfrac{{\pi }}{2}}\right)$ 上恒有 ${f}''\left( x\right)  > 0$ ,故 ${D}$ 正确.
        \end{jieda}
        
        \item 曲率是刻画曲线弯曲程度的重要指标，曲线的曲率定义如下: 记 $f'\left( x\right)$ 是 $f\left( x\right)$ 的导函数， $f''\left( x\right)$ 是 $f'\left( x\right)$ 的导函数，那么曲线 $y = f\left( x\right)$ 在点 $\left( {{x}_{0},f\left( {x}_{0}\right) }\right)$ 处的曲率 $K = \dfrac{\left| f''\left( {x}_{0}\right) \right| }{{\left\{  1 + {\left\lbrack  f'\left( {x}_{0}\right) \right\rbrack  }^{2}\right\}  }^{\frac{3}{2}}}$ ， 则曲线 $f\left( x\right)  = \sin x + \cos x$ 在点 $\left( {\dfrac{\pi }{4},\sqrt{2}}\right)$ 处的曲率为\blkda{$\sqrt{2}$}
        \begin{jieda}
            对函数 $f\left( x\right)  = \sin x + \cos x$ 求导,得 $f'\left( x\right)  = \cos x - \sin x$ 
            
            对 $f'\left( x\right)$ 求导，得 $f''\left( x\right)  =  - \sin x - \cos x$ ，所以 $f'\left( \dfrac{\pi }{4}\right)  = \cos \dfrac{\pi }{4} - \sin \dfrac{\pi }{4} = 0$ ， $f''\left( \dfrac{\pi }{4}\right)  =  - \sin \dfrac{\pi }{4} - \cos \dfrac{\pi }{4} =  - \sqrt{2}$
            
            所以曲线 $f\left( x\right)  = \sin x + \cos x$ 在点 $\left( {\dfrac{\pi }{4},\sqrt{2}}\right)$ 处的曲率 $K = \dfrac{\left| f''\left( {x}_{0}\right) \right| }{{\left\{  1 + {\left\lbrack  f'\left( {x}_{0}\right) \right\rbrack  }^{2}\right\}  }^{\frac{3}{2}}} = \dfrac{\left| -\sqrt{2}\right| }{{\left( 1 + {0}^{2}\right) }^{\frac{3}{2}}} = \sqrt{2}$ .
        \end{jieda}

        \item 对于三次函数 $f\left( x\right)  = a{x}^{3} + b{x}^{2} + {cx} + d\left( {a \neq  0}\right)$ ，定义:设 ${f}''\left( x\right)$ 是函数 $y = f\left( x\right)$ 的导函数 $y = f'\left( x\right)$ 的导数，若 ${f}''\left( x\right)  = 0$ 有实数解 ${x}_{0}$ ，则称点 $\left( {{x}_{0},f\left( {x}_{0}\right) }\right)$ 为函数 $y = f\left( x\right)$ 的“拐点”. 现已知 $f\left( x\right)  = {x}^{3} - 3{x}^{2} + {2x} - 2$ . 请解答下列问题:
        \begin{enumerate}
            \item 求函数 $f\left( x\right)$ 的“拐点” $A$ 的坐标
            \item 求证: $f\left( x\right)$ 的图像关于“拐点” $A$ 对称，并求 $f\left( {-{2020}}\right)  + f\left( {-{2019}}\right)  + \cdots f\left( {2019}\right)  + f\left( {2022}\right)$ 的值.
        \end{enumerate}
        \begin{jieda}
            \begin{enumerate}
                \item $\because f'\left( x\right)  = 3{x}^{2} - {6x} + 2,f''\left( x\right)  = {6x} - 6,\therefore$ 令 $f''\left( x\right)  = {6x} - 6 = 0$，得 $x = 1$ .
                
                有 $f\left( 1\right)  = {1}^{3} - 3 + 2 - 2 =  - 2,\therefore$ “拐点” $A$ 为 $\left( {1, - 2}\right)$ .

                \item 证明:设 $P\left( {{x}_{0},{y}_{0}}\right)$ 是 $y = f\left( x\right)$ 图像上任意一点，则 ${y}_{0} = {x}_{0}^{3} - 3{x}_{0}^{2} + 2{x}_{0} - 2$ .
                
                $P\left( {{x}_{0},{y}_{0}}\right)$ 是关于“拐点” $A\left( {1, - 2}\right)$ 的对称点为 $P\left( {2 - {x}_{0}, - 4 - {y}_{0}}\right)$ .
                
                把点 $P$ 坐标代入 $y = f\left( x\right)$ 得左边 $=  - 4 - {y}_{0} =  - {x}_{0}^{3} + 3{x}_{0}^{2} - 2{x}_{0} - 2$ ,
                
                右边 $= {\left( 2 - {x}_{0}\right) }^{3} - 3{\left( 2 - {x}_{0}\right) }^{2} + 2\left( {2 - {x}_{0}}\right)  - 2 =  - {x}_{0}^{3} + 3{x}_{0}^{2} - 2{x}_{0} - 2,\therefore$ 左边 $=$ 右边.
                
                $\therefore$ 点 $P\left( {2 - {x}_{0}, - 4 - {y}_{0}}\right)$ 在 $y = f\left( x\right)$ 的图像上.
                
                $\therefore y = f\left( x\right)$ 关于“拐点” $A$ 对称.
                
                由对称性可得 $f\left( x\right)  + f\left( {2 - x}\right)  =  - 4$
                
                $f\left( {-{2020}}\right)  + f\left( {-{2019}}\right)  + \cdots f\left( {2019}\right)  + f\left( {2022}\right)  = {2021} \times  \left( {-4}\right)  + f\left( 1\right)  =  - {8086}.$
            \end{enumerate}
        \end{jieda}

        

        \item 给出定义: 设 ${f}^{\prime }\left( x\right)$ 是函数 $y = f\left( x\right)$ 的导函数, ${f}^{\prime \prime }\left( x\right)$ 是函数 $y = {f}^{\prime }\left( x\right)$ 的导函数,若方程 ${f}^{\prime \prime }\left( x\right)  = 0$ 有实数解 $x = {x}_{0}$ ,则称 $\left( {{x}_{0},f\left( {x}_{0}\right) }\right)$ 为函数 $y = f\left( x\right)$ 的 “拐点”. 经研究发现所有的三次函数 $f\left( x\right)  = a{x}^{3} + b{x}^{2} + {cx} + d\left( {a \neq  0}\right)$ 都有“拐点”,且该“拐点”也是函数 $y = f\left( x\right)$ 的图像的对称中心. 若函数 $f\left( x\right)  =  - {x}^{3} + 3{x}^{2}$ , 则 $f\left( \dfrac{1}{2021}\right)  + f\left( \dfrac{2}{2021}\right)  + f\left( \dfrac{3}{2021}\right)  + \cdots  + f\left( \dfrac{4040}{2021}\right)  + f\left( \dfrac{4041}{2021}\right)  =$\blkda{8082}
        
        \item 给出以下三个材料:
        \begin{dingautolist}{172}
            \item 若函数 $f\left( x\right)$ 可导,我们通常把导函数 $f'\left( x\right)$ 的导数叫做 $f\left( x\right)$ 的二阶导数，记作 $f''\left( x\right)$ . 类似的，函数 $f\left( x\right)$ 的二阶导数的导数叫做函数 $f\left( x\right)$ 的三阶导数，记作 ${f}^{\prime \prime }\left( x\right)$ ，函数 $f\left( x\right)$ 的三阶导数的导数叫做函数 $f\left( x\right)$ 的四阶导数......,一般地,函数 $f\left( x\right)$ 的 $n - 1$ 阶导数的导数叫做函数 $f\left( x\right)$ 的 $n$ 阶导数,记作 ${f}^{\left( n\right) }\left( x\right)  = {\left\lbrack  {f}^{n - 1}\left( x\right) \right\rbrack  }^{\prime },n \geq  4$
            \item 若 $n \in  {{N}}^{ * }$ ，定义 $n! = n \times  \left( {n - 1}\right)  \times  \left( {n - 2}\right)  \times  \cdots  \times  3 \times  2 \times  1$
            \item 若函数 $f\left( x\right)$ 在包含 ${x}_{0}$ 的某个开区间 $\left( {a,b}\right)$ 上具有任意阶的导数，那么对于任意 $x \in  \left( {a,b}\right)$ 有 $g\left( x\right)  = f\left( {x}_{0}\right)  + \dfrac{f'\left( {x}_{0}\right) }{1!}\left( {x - {x}_{0}}\right)  + \dfrac{f''\left( {x}_{0}\right) }{2!}{\left( x - {x}_{0}\right) }^{2} + \cdots  + \dfrac{{f}^{\left( n\right) }\left( {x}_{0}\right) }{n!}{\left( x - {x}_{0}\right) }^{n} + \cdots$ ,我们将 $g\left( x\right)$ 称为函数 $f\left( x\right)$ 在点 $x = {x}_{0}$ 处的泰勒展开式.

            例如 ${f}_{1}\left( x\right)  = {{e}}^{x}$ 在点 $x = 0$ 处的泰勒展开式为 ${g}_{1}\left( x\right)  = 1 + x + \dfrac{1}{2}{x}^{2} + \cdots  + \dfrac{1}{n!}{x}^{n} + \cdots$
        \end{dingautolist}
        
        根据以上三段材料，完成下面的题目:
        \begin{enumerate}
            \item 求出 $f\left( x\right)  = \cos x$ 在点 $x = 0$ 处的泰勒展开式 $g\left( x\right)$
            \item 用 $f\left( x\right)  = \cos x$ 在点 $x = 0$ 处的泰勒展开式前三项计算 $\cos {0.3}$ 的值，精确到小数点后 4 位.
            \item 现已知 $\dfrac{\sin x}{x} = \left( {1 - \dfrac{x}{\pi }}\right) \left( {1 + \dfrac{x}{\pi }}\right) \left( {1 - \dfrac{x}{2\pi }}\right) \left( {1 + \dfrac{x}{2\pi }}\right) \cdots \left( {1 - \dfrac{x}{n \pi }}\right) \left( {1 + \dfrac{x}{n \pi }}\right) \cdots$ ，试求 $\mathop{\sum }\limits_{{n = 1}}^{\infty }\dfrac{1}{{n}^{2}}$ 的值
        \end{enumerate}
        \begin{jieda}
            \begin{enumerate}
                \item $f\left( x\right)  = \cos x, f'\left( x\right)  =  - \sin x,{f}''\left( x\right)  =  - \cos x,\cdots$ ，
                
                所以 $f\left( 0\right)  = \cos 0 = 1,f'\left( 0\right)  =  - \sin 0 = 0,{f}''\left( x\right)  =  - \cos 0 =  - 1,\cdots$ ,
                
                由 $\cos x = 1 + \dfrac{0}{1!}\left( {x - 0}\right)  + \dfrac{-1}{2!}{\left( x - 0\right) }^{2} + \cdots  + \dfrac{{\left( -1\right) }^{n}}{n!}{\left( x - 0\right) }^{2n} + \cdots$
                
                所以 $\cos x = 1 - \dfrac{{x}^{2}}{2!} + \dfrac{{x}^{4}}{4!} - \dfrac{{x}^{6}}{6!} + \cdots  + \dfrac{{\left( -1\right) }^{n}{x}^{2n}}{\left( {2n}\right) !} + \cdots$

                \item 由(1)可得 $\cos {0.3} = 1 - \dfrac{{0.3}^{2}}{2!} + \dfrac{{0.3}^{4}}{4!} - \dfrac{{0.3}^{6}}{6!} + \cdots  + \dfrac{{\left( -1\right) }^{n} \times  {0.3}^{2n}}{\left( {2n}\right) !} + \cdots  \approx  1 - \dfrac{{0.3}^{2}}{2!} + \dfrac{{0.3}^{4}}{4!}  = 1 - {0.045} + {0.0003375} = {0.9553}$

                \item 因为 $\dfrac{\sin x}{x} = \left( {1 - \dfrac{x}{\pi}}\right) \left( {1 + \dfrac{x}{\pi}}\right) \left( {1 - \dfrac{x}{2\pi}}\right) \left( {1 + \dfrac{x}{2\pi}}\right) \cdots \left( {1 - \dfrac{x}{n\pi}}\right) \left( {1 + \dfrac{x}{n\pi}}\right) \cdots  = \left( {1 - \dfrac{{x}^{2}}{{{\pi}}^{2}}}\right) \left( {1 - \dfrac{{x}^{2}}{4{{\pi}}^{2}}}\right) \cdots \left( {1 - \dfrac{{x}^{2}}{{n}^{2}{{\pi}}^{2}}}\right)$
                
                对 $\cos x = 1 - \dfrac{{x}^{2}}{2!} + \dfrac{{x}^{4}}{4!} - \dfrac{{x}^{6}}{6!} + \cdots  + \dfrac{{\left( -1\right) }^{n}{x}^{2n}}{\left( {2n}\right) !} + \cdots$ ,
                
                两边求导可得: $- \sin x =  - \dfrac{x}{1!} + \dfrac{{x}^{3}}{3!} - \dfrac{{x}^{5}}{5!} + \cdots  + \dfrac{{\left( -1\right) }^{n}{x}^{{2n} - 1}}{\left( {{2n} - 1}\right) !} + \cdots$ ,
                
                所以 $\sin x = \dfrac{x}{1!} - \dfrac{{x}^{3}}{3!} + \dfrac{{x}^{5}}{5!} + \cdots  + \dfrac{{\left( -1\right) }^{n - 1}{x}^{{2n} - 1}}{\left( {{2n} - 1}\right) !} + \cdots$ ,
                
                所以 $\dfrac{\sin x}{x} = 1 - \dfrac{{x}^{2}}{3!} + \dfrac{{x}^{4}}{5!} + \cdots  + \dfrac{{\left( -1\right) }^{n - 1}{x}^{{2n} - 2}}{\left( {{2n} - 1}\right) !}$
                
                比较两式中 ${x}^{2}$ 的系数，可得：$- \dfrac{1}{3!} =  - \dfrac{1}{{\pi }^{2}}\left( {\dfrac{1}{{1}^{2}} + \dfrac{1}{{2}^{2}} + \dfrac{1}{{3}^{2}}\cdots  + \dfrac{1}{{n}^{2}} + \cdots }\right)$
                
                所以 $\mathop{\sum }\limits_{{n = 1}}^{\infty }\dfrac{1}{{n}^{2}} = \dfrac{1}{{1}^{2}} + \dfrac{1}{{2}^{2}} + \dfrac{1}{{3}^{2}}\cdots  + \dfrac{1}{{n}^{2}} + \cdots  = \dfrac{{\pi}^{2}}{6}$ .
            \end{enumerate}
        \end{jieda}
    \end{enumerate}

    \item 利用导数判断单调性、求单调区间
    \begin{enumerate}
        \item 函数 $f\left( x\right)  = \ln x + \dfrac{{{e}}^{x} + 1}{x}$ 的单调增区间为\blkda{$\left( {1, + \infty }\right)$}
        \begin{jieda}
            由函数 $f\left( x\right)  = \ln x + \dfrac{{{e}}^{x} + 1}{x}$ ,可得其定义域为 $\left( {0, + \infty }\right)$，且 ${f}^{\prime }\left( x\right)  = \dfrac{1}{x} + \dfrac{x{{e}}^{x} - \left( {{{e}}^{x} + 1}\right) }{{x}^{2}} = \dfrac{\left( {x - 1}\right) \left( {{{e}}^{x} + 1}\right) }{{x}^{2}}$，令 ${f}^{\prime }\left( x\right)  > 0$ ,解得 $x > 1$ ,所以函数 $f\left( x\right)$ 的单调增区间为 $\left( {1, + \infty }\right)$ .
        \end{jieda}
        
        \item 已知函数 $f\left( x\right)  = \ln \left( {\cos x}\right)  + \dfrac{1}{2}\sin {2x} - x$ ，则\dotfill{(\kh{B})}
        \xx{$f\left( x\right)$ 在区间 $\left( {0,\dfrac{\pi }{2}}\right)$ 内单调递增}{$f\left( x\right)$ 在区间 $\left( {0,\dfrac{\pi }{2}}\right)$ 内单调递减}{$f\left( x\right)$ 在区间 $\left( {0,\dfrac{\pi }{4}}\right)$ 内单调递增，在区间 $\left( {\dfrac{\pi }{4},\dfrac{\pi }{2}}\right)$ 内单调递减}{$f\left( x\right)$ 在区间 $\left( {0,\dfrac{\pi }{4}}\right)$ 内单调递减，在区间 $\left( {\dfrac{\pi }{4},\dfrac{\pi }{2}}\right)$ 内单调递增}
        
        \begin{jieda}
            由题意 $\cos x > 0,{f}^{\prime }\left( x\right)  =  - \dfrac{\sin x}{\cos x} + \cos {2x} - 1 =  - \left( {\tan x + 1}\right)  + \dfrac{1 - {\tan }^{2}x}{1 + {\tan }^{2}x} = \dfrac{-\tan x{\left( \tan x + 1\right) }^{2}}{{\tan }^{2}x + 1}$
            
            当 $x \in  \left( {0,\dfrac{\pi}{2}}\right)$ 时， $\tan x > 0$ ，则 ${f}^{\prime }\left( x\right)  < 0$ ，所以 $f\left( x\right)$ 在区间 $\left( {0,\dfrac{\pi }{2}}\right)$ 上单调递减.
        \end{jieda}
        

        \item 已知函数 $f\left( x\right)  = {{e}}^{x} + 2{f}^{\prime }\left( 0\right) x - \cos x$ .
        \begin{enumerate}
            \item 求 $f\left( x\right)$ 的解析式
            \item 判断 $f\left( x\right)$ 在 $\left( {-\infty ,0}\right)$ 上的单调性.
        \end{enumerate}
        \begin{jieda}
            \begin{enumerate}
                \item 因为 $f\left( x\right)  = {{e}}^{x} + 2{f}^{\prime }\left( 0\right) x - \cos x$ ，所以 ${f}^{\prime }\left( x\right)  = {{e}}^{x} + 2{f}^{\prime }\left( 0\right)  + \sin x$ ，
                
                则 ${f}^{\prime }\left( 0\right)  = {{e}}^{0} + 2{f}^{\prime }\left( 0\right)  + \sin 0 = 1 + 2{f}^{\prime }\left( 0\right)$ ，所以 ${f}^{\prime }\left( 0\right)  =  - 1$ ，
                
                所以 $f\left( x\right)  = {{e}}^{x} - {2x} - \cos x$ .
                \item ${f}^{\prime }\left( x\right)  = {{e}}^{x} - 2 + \sin x$ ,
                
                当 $x \leq  0$ 时， $0 < {{e}}^{x} \leq  1,\sin x \leq  1$ ，
                
                所以 ${f}^{\prime }\left( x\right)  = {{e}}^{x} - 2 + \sin x \leq  0$ 恒成立,
                
                所以 $f\left( x\right)$ 在 $( - \infty ,0\rbrack$ 上的单调递减.
            \end{enumerate}
        \end{jieda}
        
        \item 已知函数 $f\left( x\right)  = {x}^{2} - \left( {{2a} + 1}\right) x + a\ln x + a$ .
        \begin{enumerate}
            \item 当 $a = 1$ 时，求函数 $f\left( x\right)$ 的单调区间
            \item 当 $a \in  {R}$ 时,求 $f\left( x\right)$ 的单调区间
        \end{enumerate}
        \begin{jieda}
            \begin{enumerate}
                \item 当 $a = 1$ 时， $f\left( x\right)  = {x}^{2} - {3x} + \ln x + 1$ ，定义域为 $\left( {0, + \infty }\right)$ ，
                
                ${f}^{\prime }\left( x\right)  = {2x} - 3 + \dfrac{1}{x} = \dfrac{\left( {{2x} - 1}\right) \left( {x - 1}\right) }{x},$
                
                令 ${f}^{\prime }\left( x\right)  > 0$ ,得 $x \in  \left( {0,\dfrac{1}{2}}\right)  \cup  \left( {1, + \infty }\right)$ ,令 ${f}^{\prime }\left( x\right)  < 0$ ,得 $x \in  \left( {\dfrac{1}{2},1}\right)$ ,
                
                所以 $f\left( x\right)$ 的单调递增区间为 $\left( {0,\dfrac{1}{2}}\right) ,\left( {1, + \infty }\right)$ ,单调递减区间为 $\left( {\dfrac{1}{2},1}\right)$ .

                \item $f\left( x\right)  = {x}^{2} - \left( {{2a} + 1}\right) x + a\ln x + a$ ，定义域为 $\left( {0, + \infty }\right)$ ，
                
                ${f}^{\prime }\left( x\right)  = {2x} - \left( {{2a} + 1}\right)  + \dfrac{a}{x} = \dfrac{\left( {{2x} - 1}\right) \left( {x - a}\right) }{x}$ ,令 ${f}^{\prime }\left( x\right)  = 0$ ,得 $x = \dfrac{1}{2}$ 或 $x = a$ .

                \begin{dingautolist}{172}
                    \item 当 $a \leq  0$ 时，当 $x \in  \left( {0,\dfrac{1}{2}}\right)$ 时， ${f}^{\prime }\left( x\right)  < 0,f\left( x\right)$ 单调递减，
                    
                    当 $x \in  \left( {\dfrac{1}{2}, + \infty }\right)$ 时， ${f}^{\prime }\left( x\right)  > 0,f\left( x\right)$ 单调递增
                    \item 当 $0 < a < \dfrac{1}{2}$ 时，当 $x \in  \left( {0,a}\right)$ 和 $x \in  \left( {\dfrac{1}{2}, + \infty }\right)$ 时， ${f}^{\prime }\left( x\right)  > 0,f\left( x\right)$ 单调递增，
                    
                    当 $x \in  \left( {a,\dfrac{1}{2}}\right)$ 时， ${f}^{\prime }\left( x\right)  < 0,f\left( x\right)$ 单调递减

                    \item 当 $a = \dfrac{1}{2}$ 时， ${f}^{\prime }\left( x\right)  \geq  0$ 对 $\forall x \in  \left( {0, + \infty }\right)$ 恒成立，所以 $f\left( x\right)$ 在 $\left( {0, + \infty }\right)$ 单调递增

                    \item 当 $a > \dfrac{1}{2}$ 时，当 $x \in  \left( {0,\dfrac{1}{2}}\right)$ 和 $x \in  \left( {a, + \infty }\right)$ 时， ${f}^{\prime }\left( x\right)  > 0,f\left( x\right)$ 单调递增，
                
                    当 $x \in  \left( {\dfrac{1}{2},a}\right)$ 时， ${f}^{\prime }\left( x\right)  < 0,f\left( x\right)$ 单调递减
                \end{dingautolist}
            \end{enumerate}
        \end{jieda}
        
        \item 设函数 $y = f\left( x\right)$ ,其中 $f\left( x\right)  = {kx} + \sin x,k \geq  1,k \in  \mathbf{R}$ ,定义域为 $D = \left( {-\infty , + \infty }\right)$ . 对于 $t \in  D$ ,定义${S}_{f\left( t\right) } = \{ x \mid  f\left( x\right)  \leq  f\left( t\right) \}$ ,则 ${S}_{f\left( {2025}\right) } =$\blkda{$(-\infty, 2025]$}
        \begin{jieda}
            $f'(x) = k + cosx \geq 0$，因此$f(x)$严增，故${S}_{f\left( t\right) } = \{ x \mid  f\left( x\right)  \leq  f\left( t\right) \} = (0, t]$.
        \end{jieda}
    
    \end{enumerate}


    \item 根据单调性解不等式
    \begin{enumerate}
        \item 定义在 $(-\infty, 0) \cup \left( {0, + \infty }\right)$的函数 $f\left( x\right)$ 的导函数为 ${f}^{\prime }\left( x\right)$ ,若 $x{f}^{\prime }\left( x\right)  + f\left( x\right) = e^x$ ,且 $f\left( 2\right)  = 0$ ,则不等式 $\left( {x - 2}\right) f\left( x\right)  < 0$ 的解集为\blkda{$\left(-\infty, 0\right)$}
        \begin{jieda}
            $x{f}^{\prime }\left( x\right)  + f\left( x\right) = e^x \Longleftrightarrow \left[xf(x)\right]' = e^x$，则$f(x) = \dfrac{e^x+c}{x}$，$f(2) = 0$则$f(x) = \dfrac{e^x-e^2}{x}$

            \begin{dingautolist}{172}
                \item 当$x > 2$时，$f(x) > 0, x-2 > 0$，故不满足
                \item 当$x \in (0, 2)$时，$f(x) < 0, x-2 < 0$，故不满足
                \item 当$x < 0$时，$f(x) > 0, x-2 < 0$，故满足.
            \end{dingautolist}
        \end{jieda}
        \item 定义在 $\left( {0, + \infty }\right)$ 上的函数 $f\left( x\right)$ 的导函数为 ${f}^{\prime }\left( x\right)$ ,若 $x{f}^{\prime }\left( x\right)  - f\left( x\right)  < 0$ ,且 $f\left( 3\right)  = 0$ ,则不等式 $\left( {x - 2}\right) f\left( x\right)  < 0$ 的解集为\blkda[3]{$\left( {0,2}\right)  \cup  \left( {3, + \infty }\right)$}
        \begin{jieda}
            依题意,令 $g\left( x\right)  = \dfrac{f\left( x\right) }{x}$ ,求导得 ${g}^{\prime }\left( x\right)  = \dfrac{x{f}^{\prime }\left( x\right)  - f\left( x\right) }{{x}^{2}} < 0$ ,则 $g\left( x\right)$ 在 $\left( {0, + \infty }\right)$ 上单调递减, 由 $f\left( 3\right)  = 0$ ,得 $g\left( 3\right)  = 0$ ,不等式 $\left( {x - 2}\right) f\left( x\right)  < 0 \Longleftrightarrow  \left( {x - 2}\right)  \cdot  \dfrac{f\left( x\right) }{x} < 0 \Longleftrightarrow  \left( {x - 2}\right) g\left( x\right)  < 0$ , 则 $\left\{  \begin{array}{l} x - 2 < 0 \\  g\left( x\right)  > 0 \end{array}\right.$ 或 $\left\{  \begin{array}{l} x - 2 > 0 \\  g\left( x\right)  < 0 \end{array}\right.$ ,即 $\left\{  \begin{matrix} x < 2 \\  0 < x < 3 \end{matrix}\right.$ 或 $\left\{  \begin{array}{l} x > 2 \\  x > 3 \end{array}\right.$ ,解得 $0 < x < 2$ 或 $x > 3$ , 所以不等式 $\left( {x - 2}\right) f\left( x\right)  < 0$ 的解集为 $\left( {0,2}\right)  \cup  \left( {3, + \infty }\right)$ .
        \end{jieda}
        
        \item 已知定义在 $\left\lbrack  {-3,3}\right\rbrack$ 上的函数 $f\left( x\right)  = {{e}}^{x} - {{e}}^{-x} - {2x} + 1$ ，若 $f\left( {m}^{2}\right)  +  f\left( {m - 2}\right)  \leq  2$ ，则 $m$ 的取值范围是\blkda{$\left\lbrack  {-1,1}\right\rbrack$}

        \begin{jieda}
            记 $g\left( x\right)  = {{e}}^{x} - {{e}}^{-x} - {2x},x \in  \left\lbrack  {-3,3}\right\rbrack$ ,则 $g\left( {-x}\right)  = {{e}}^{-x} - {{e}}^{x} + {2x} =  - g\left( x\right)$，故 $g\left( x\right)$ 为 $\left\lbrack  {-3,3}\right\rbrack$ 的奇函数
            
            因为${g}^{\prime }\left( x\right)  = {{e}}^{x} + {{e}}^{-x} - 2 \geq  2\sqrt{{{e}}^{x}{{e}}^{-x}} - 2 = 0$，所以$g\left( x\right)$ 为 $\left\lbrack  {-3,3}\right\rbrack$ 上的单调递增函数
            
            因为 $f\left( x\right)  = g\left( x\right)  + 1$，由 $f\left( {m}^{2}\right)  + f\left( {m - 2}\right)  \leq  2$ 可得 $g\left( {m}^{2}\right)  + g\left( {m - 2}\right)  + 2 \leq  2$ ,进而 $g\left( {m}^{2}\right)  \leq  g\left( {2 - m}\right)$ ,
            
            故 $- 3 \leq  {m}^{2} \leq  2 - m \leq  3$ ，解得 $- 1 \leq  m \leq  1$
        \end{jieda}
        
        \item 已知定义在 $\mathbf{R}$ 上的函数 $f\left( x\right) ,f\left( {x - 1}\right)$ 关于直线 $x = 1$ 对称,当 $x \in  \left( {0, + \infty }\right)$ 时, ${f}^{\prime }\left( x\right)  > 0$ ,设 $a = f\left( {\sqrt{3}\ln 2}\right) ,b = f\left( {\ln 3}\right) ,c = f\left( {-\dfrac{1}{2}}\right)$ ,则 $a,b,c$ 的大小关系为\dotfill{(\kh{C})}
        \xx{$c < a < b$}{$b < a < c$}{$a > b > c$}{$a < b < c$}

        \begin{jieda}
            因为 $f\left( {x - 1}\right)$ 关于直线 $x = 1$ 对称,则 $f\left( x\right)$ 是偶函数,
            
            当 $x \in  \left( {0, + \infty }\right)$ 时, ${f}^{\prime }\left( x\right)  > 0$ ,所以函数 $f\left( x\right)$ 在 $\left( {0, + \infty }\right)$ 上单调递增.
            
            因为 $\ln 3 > \ln {e} = 1,\left| {-\dfrac{1}{2}}\right|  = \dfrac{1}{2} < 1$ ,
            
            再比较 $\sqrt{3}\ln 2$ 和 $\ln 3$ 的大小，构造函数 $g\left( x\right)  = \dfrac{\ln x}{x}$ ，
            
            则 ${g}^{\prime }\left( x\right)  = \dfrac{1 - \ln x}{{x}^{2}}$ ,令 ${g}^{\prime }\left( x\right)  = \dfrac{1 - \ln x}{{x}^{2}} > 0$ ,得 $0 < x < {e}$ ,此时 $g\left( x\right)  = \dfrac{\ln x}{x}$ 单调递增,
            
            所以 $g\left( 2\right)  = \dfrac{\ln 2}{2} > g\left( \sqrt{3}\right)  = \dfrac{\ln \sqrt{3}}{\sqrt{3}}$ ,也即 $\sqrt{3}\ln 2 > \ln 3 > \ln {e} = 1$ ,
            
            结合函数 $f\left( x\right)$ 在 $\left( {0, + \infty }\right)$ 上单调递增及偶函数可知 $a > b > c$ .
        \end{jieda}
        
        \item 已知 $a = \pi \ln 3,b = 3\ln \pi ,c = \dfrac{3\pi }{{e}}$ ，则\dotfill{(\kh{C})}
        \xx{$a < b < c$}{$c < a < b$}{$b < a < c$}{$c < b < a$}
        \begin{jieda}
            因为 $a = \pi \ln 3,b = 3\ln \pi ,c = \dfrac{3\pi }{{e}}$ ,即 $\dfrac{a}{3\pi } = \dfrac{\ln 3}{3},\dfrac{b}{3\pi } = \dfrac{\ln \pi }{\pi },\dfrac{c}{3\pi } = \dfrac{1}{{e}} = \dfrac{\ln {e}}{{e}}$ ,
            
            根据特点设函数 $f\left( x\right)  = \dfrac{\ln x}{x},x > 0$ ,则 ${f}^{\prime }\left( x\right)  = \dfrac{1 - \ln x}{{x}^{2}}$ ,
            
            当 $x \in  \left( {{e}, + \infty }\right)$ 时， ${f}^{\prime }\left( x\right)  < 0$ ，可知 $f\left( x\right)$ 在 $\left( {{e}, + \infty }\right)$ 内单调递减，
            
            且 ${e} < 3 < \pi$ ，则 $f\left( {e}\right)  > f\left( 3\right)  > f\left( \pi \right)$ ，即 $\dfrac{c}{3\pi } > \dfrac{a}{3\pi } > \dfrac{b}{3\pi }$ ，所以 $b < a < c$ .
        \end{jieda}
    \end{enumerate}

    \item 利用导数求函数的极值
    \begin{enumerate}
        \item 已知函数 $f\left( x\right)  = \dfrac{{2x} - 1}{x - 1}{{e}}^{x}$ ,则\dotfill{(\kh{C})}
        \xx{$f\left( x\right)$ 有极大值，无极小值}{$f\left( x\right)$ 无极大值，有极小值}{$f\left( x\right)$ 既有极大值，也有极小值}{$f\left( x\right)$ 既无极大值，也无极小值}

        \begin{jieda}
            $f\left( x\right)$ 定义域为 $\left( {-\infty ,1}\right)  \cup  \left( {1, + \infty }\right)$ ,且 ${f}^{\prime }\left( x\right)  = \dfrac{x\left( {{2x} - 3}\right) }{{\left( x - 1\right) }^{2}}{{e}}^{x}$ ,
            
            由 ${f}^{\prime }\left( x\right)  > 0$ 可得: $x \in  \left( {-\infty ,0}\right)  \cup  \left( {\dfrac{3}{2}, + \infty }\right)$ ,所以 $f\left( x\right)$ 在 $\left( {-\infty ,0}\right) ,\left( {\dfrac{3}{2}, + \infty }\right)$ 单调递增;
            
            由 ${f}^{\prime }\left( x\right)  < 0$ ,可得: $x \in  \left( {0,1}\right)  \cup  \left( {1,\dfrac{3}{2}}\right)$ ,所以 $f\left( x\right)$ 在 $\left( {0,1}\right) ,\left( {1,\dfrac{3}{2}}\right)$ 单调递减.
            
            所以 $f\left( x\right)$ 有极大值 $f\left( 0\right)$ ,有极小值 $f\left( \dfrac{3}{2}\right)$ .
        \end{jieda}
        
        \item 函数 $f\left( x\right)  = \dfrac{{x}^{2}}{2} - {3x} + 2\ln x$ 的极大值与极小值之和为\blkda{$2\ln 2 - \dfrac{13}{2}$}

        \begin{jieda}
            由题意知: ${f}^{\prime }\left( x\right)  = x - 3 + \dfrac{2}{x} = \dfrac{{x}^{2} - {3x} + 2}{x} = \dfrac{\left( {x - 1}\right) \left( {x - 2}\right) }{x}\left( {x > 0}\right)$
            
            当 $x \in  \left( {1,2}\right)$ 时， ${f}^{\prime }\left( x\right)  < 0,f\left( x\right)$ 单调递减；当 $x \in  \left( {0,1}\right) ,\left( {2, + \infty }\right)$ 时，
            
            ${f}^{\prime }\left( x\right)  > 0,f\left( x\right)$ 单调递增,所以 $f\left( x\right)$ 的极大值为 $f\left( 1\right)  =  - \dfrac{5}{2}$ ,
            
            极小值为 $f\left( 2\right)  = 2\ln 2 - 4$ ,故 $f\left( 1\right)  + f\left( 2\right)  = 2\ln 2 - \dfrac{13}{2}$ .
        \end{jieda}
        
        \item 已知函数 $f\left( x\right)  = \left( {{x}^{2} + {ax} + b}\right) {{e}}^{x}$ 的图象在点 $\left( {0,f\left( 0\right) }\right)$ 处的切线方程为 ${2x} + y - 1 = 0$ .
        \begin{enumerate}
            \item 求$a$，$b$的值
            \item 求 $f\left( x\right)$ 的单调区间与极值
        \end{enumerate}
        \begin{jieda}
            \begin{enumerate}
                \item 由题意，易知 $f\left( 0\right)  = 1$ ，得 $b = 1$ .
                
                ${f}^{\prime }\left( x\right)  = \left( {{2x} + a}\right) {{e}}^{x} + \left( {{x}^{2} + {ax} + 1}\right) {{e}}^{x} = \left\lbrack  {{x}^{2} + \left( {a + 2}\right) x + a + 1}\right\rbrack  {{e}}^{x},$
                
                由 ${f}^{\prime }\left( 0\right)  = a + 1 =  - 2$ ，解得 $a =  - 3$ .

                \item 由(1)知 ${f}^{\prime }\left( x\right)  = \left( {{x}^{2} - x - 2}\right) {{e}}^{x} = \left( {x + 1}\right) \left( {x - 2}\right) {{e}}^{x}$ ，易知 ${{e}}^{x} > 0$ ，
                
                当 $x$ 变化时， $f\left( x\right)$ ， ${f}^{\prime }\left( x\right)$ 的变化情况如下表所示.
                
                \begin{center}
                \begin{tabular}{|c|c|c|c|c|c|}
                \hline
                $x$ & $(-0, - 1)$ & $-1$ & $(-1,2)$ & $2$ & $(2, +\infty)$ \\
                \cline{1-6}
                ${f}^{\prime }\left( x\right)$ & + & 0 & - & 0 & + \\
                \cline{1-6}
                $f\left( x\right)$ & 单调递增 & $\dfrac{5}{e}$ & 单调递减 & $- {{e}}^{2}$ & 单调递增 \\
                \cline{1-6}
                \hline
                \end{tabular}
                \end{center}
                
                因此，函数 $f\left( x\right)$ 在 $\left( {-\infty , - 1}\right)$ 和 $\left( {2, + \infty }\right)$ 上单调递增，在 $\left( {-1,2}\right)$ 上单调递减.
                
                当 $x =  - 1$ 时, $f\left( x\right)$ 有极大值,且极大值为 $f\left( {-1}\right)  = \dfrac{5}{{e}}$ ;
                
                当 $x = 2$ 时, $f\left( x\right)$ 有极小值,且极小值为 $f\left( 2\right)  =  - {{e}}^{2}$ .
            \end{enumerate}
        \end{jieda}
    \end{enumerate}

    \item 根据极值 (点) 求参数
    \begin{enumerate}
        \item 已知 $f\left( x\right)  = {x}^{3} + {3a}{x}^{2} + {bx} + {a}^{2}$ 在 $x =  - 1$ 处有极值 0，则 $a - b =$\blkda{$-7$}
        \begin{jieda}
            根据题意， ${f}^{\prime }\left( x\right)  = 3{x}^{2} + {6ax} + b$ ,
            
            $\because$ 函数 $f\left( x\right)$ 在 $x =  - 1$ 处有极值 0,
            
            $\therefore {f}^{\prime }\left( {-1}\right)  = 3 - {6a} + b = 0$ 且 $f\left( {-1}\right)  =  - 1 + {3a} - b + {a}^{2} = 0$ ,
            
            $\therefore a = 1,b = 3$ 或 $a = 2,b = 9$ ,
            
            $a = 1,b = 3$ 时 ${f}^{\prime }\left( x\right)  = 3{x}^{2} + {6x} + 3 \geq  0$ 恒成立,此时函数无极值点,
            
            当 $a = 2,b = 9$ 时， ${f}^{\prime }\left( x\right)  = 3{x}^{2} + {12x} + 9 = 3\left( {x + 1}\right) \left( {x + 3}\right)$ ,
            
            此时 $x =  - 1$ 是函数的极值,满足条件,
            
            $\therefore a = 2,b = 9,\therefore a - b =  - 7$ .
        \end{jieda}

        
        \item 若函数 $f\left( x\right)  = \left( {{x}^{2} + {ax} + 3}\right) {{e}}^{x}$ 在 $\left( {0, + \infty }\right)$ 上无极值点,则实数 $a$ 的取值范围是\blkda{$a \geq   - 2\sqrt{2}$}
        \begin{jieda}
            因为 ${f}^{\prime }\left( x\right)  = \left\lbrack  {{x}^{2} + \left( {a + 2}\right) x + a + 3}\right\rbrack  {{e}}^{x}$ ,令 $g\left( x\right)  = {x}^{2} + \left( {a + 2}\right) x + a + 3$，对称轴为$x = -(a+2)$
            \begin{dingautolist}{172}
                \item $-(a+2) \leq 0$，此时应有$g(0) = a+3 \geq 0$
                \item $-(a+2) > 0$，此时应有$\Delta \leq 0$
            \end{dingautolist}
            综上，解得$a \geq   - 2\sqrt{2}$
        \end{jieda}

        
        \item 若函数 $f\left( x\right)  = a\ln x + \dfrac{3}{x} - x$ 既有极大值也有极小值，则实数 $a$ 的取值范围为\blkda{$\left( {2\sqrt{3}, + \infty }\right)$}

        \begin{jieda}
            由题意可知: $f\left( x\right)$ 的定义域为 $\left( {0, + \infty }\right)$ ,且 ${f}^{\prime }\left( x\right)  = \dfrac{a}{x} - \dfrac{3}{{x}^{2}} - 1 =  - \dfrac{{x}^{2} - {ax} + 3}{{x}^{2}}$
            
            若函数 $f\left( x\right)$ 既有极大值也有极小值,则 ${x}^{2} - {ax} + 3 = 0$ 有 2 个不相等的正根 $x_1,x_2$，则 $\left\{  \begin{matrix} \Delta  = {a}^{2} - {12} > 0 \\  {x}_{1} + {x}_{2} = a > 0 \\  {x}_{1}{x}_{2} = 3 > 0 \end{matrix}\right.$ ,解得 $a > 2\sqrt{3}$
        \end{jieda}
        
        \item 已知函数 $f\left( x\right)  = \left( {{x}^{2} + 2}\right)  - k\ln \left( {x + 1}\right)$ ( $k$ 为常数).
        \begin{enumerate}
            \item 当 $k = 1$ 时，求 $f\left( x\right)$ 在 $x = 0$ 处的切线方程
            \item 若函数 $f\left( x\right)$ 在区间 $\left\lbrack  {0,1}\right\rbrack$ 上存在极值，求实数 $k$ 的取值范围
        \end{enumerate}
        
        \begin{jieda}
            \begin{enumerate}
                \item 当 $k = 1$ 时， $f\left( x\right)  = \left( {{x}^{2} + 2}\right)  - \ln \left( {x + 1}\right) ,x \in  \left( {-1, + \infty }\right)$ ,
                
                所以 $f\left( 0\right)  = 2$ ,
                
                所以 ${f}^{\prime }\left( x\right)  = {2x} - \dfrac{1}{x + 1}$
                
                所以 $f\left( x\right)$ 在 $x = 0$ 处的切线的斜率为 $k = {f}^{\prime }\left( 0\right)  = 2 \times  0 - \dfrac{1}{0 + 1} =  - 1$ ,
                
                所以 $f\left( x\right)$ 在 $x = 0$ 处的切线方程为 $y - 2 = \left( {-1}\right) \left( {x - 0}\right)$ ,即 $x + y - 2 = 0$ .

                \item 因为 $f\left( x\right)  = \left( {{x}^{2} + 2}\right)  - k\ln \left( {x + 1}\right) ,x \in  \left( {-1, + \infty }\right)$ ,
                
                所以 ${f}^{\prime }\left( x\right)  = {2x} - \dfrac{k}{x + 1}$ ,
                
                因为函数 $f\left( x\right)$ 在区间 $\left( {0,1}\right)$ 上存在极值,
                
                所以 $\exists {x}_{0} \in  \left( {0,1}\right)$ ,使得 ${f}^{\prime }\left( {x}_{0}\right)  = 0,{x}_{0}$ 两侧的导数异号,
                
                所以 $2{x}_{0} - \dfrac{k}{{x}_{0} + 1} = 0$ ,即 $k = 2{x}_{0}^{2} + 2{x}_{0},{x}_{0} \in  \left( {0,1}\right)$ ,
                
                令 $g\left( {x}_{0}\right)  = 2{x}_{0}^{2} + 2{x}_{0},{x}_{0} \in  \left( {0,1}\right)$ ,
                
                由二次函数的性质知,对称轴为 ${x}_{0} =  - \dfrac{1}{2}$ ,开口向上,
                
                所以 $g\left( {x}_{0}\right)$ 在 $\left( {0,1}\right)$ 上单调递增,
                
                所以 $g\left( 0\right)  < g\left( {x}_{0}\right)  < g\left( 1\right)$ ,即 $0 < g\left( {x}_{0}\right)  < 4$ ,
                
                所以实数 $k$ 的取值范围为 $\left( {0,4}\right)$ .
            \end{enumerate}
        \end{jieda}
        
        \item 已知函数 $f\left( x\right)  = \dfrac{1}{2}{x}^{2} - \left( {a + 1}\right) x + a\ln x$ .
        \begin{enumerate}
            \item 讨论 $f\left( x\right)$ 的单调性
            \item 若函数 $f\left( x\right)$ 有极小值，且极小值大于 $\left( {1 - \dfrac{1}{2}{{e}}^{2}}\right) a$ ，求实数 $a$ 的取值范围
        \end{enumerate}
        \begin{jieda}
            \begin{enumerate}
                \item 函数 $f\left( x\right)  = \dfrac{1}{2}{x}^{2} - \left( {a + 1}\right) x + a\ln x$ 的定义域为 $\left( {0, + \infty }\right)$ ,
                
                ${f}^{\prime }\left( x\right)  = \dfrac{a}{x} + x - \left( {a + 1}\right)  = \dfrac{{x}^{2} - \left( {a + 1}\right) x + a}{x} = \dfrac{\left( {x - 1}\right) \left( {x - a}\right) }{x}$

                \begin{dingautolist}{172}
                    \item 当 $a > 1$ 时,令 ${f}^{\prime }\left( x\right)  > 0$ ,得 $x > a$ 或 $0 < x < 1$ ,令 ${f}^{\prime }\left( x\right)  < 0$ ,得 $1 < x < a$ ,
                    
                    所以 $f\left( x\right)$ 的递增区间是 $\left( {a, + \infty }\right)$ 和 $\left( {0,1}\right)$ ,递减区间是 $\left( {1,a}\right)$

                    \item 当 $a = 1$ 时， ${f}^{\prime }\left( x\right)  \geq  0$ 恒成立，所以 $f\left( x\right)$ 的递增区间是 $\left( {0, + \infty }\right)$

                    \item 当 $0 < a < 1$ 时,令 ${f}^{\prime }\left( x\right)  > 0$ ,得 $0 < x < a$ 或 $x > 1$ ,令 ${f}^{\prime }\left( x\right)  < 0$ 得 $a < x < 1$ ,
                    
                    所以 $f\left( x\right)$ 的递增区间是 $\left( {0,a}\right)$ 和 $\left( {1, + \infty }\right)$ ,递减区间是 $\left( {a,1}\right)$
                    \item 当 $a \leq  0$ 时,令 ${f}^{\prime }\left( x\right)  > 0$ 得 $x > 1$ ,令 ${f}^{\prime }\left( x\right)  < 0$ 得 $0 < x < 1$ ,
                
                    所以 $f\left( x\right)$ 的递增区间是 $\left( {1, + \infty }\right)$ ,递减区间是 $\left( {0,1}\right)$ 
                \end{dingautolist}
                
                \item 由(1)知:若函数 $f\left( x\right)$ 有极小值，则 $a \neq  1$ ，

                \begin{dingautolist}{172}
                    \item 当 $a \leq  0$ 或 $0 < a < 1$ 时， $f\left( x\right)$ 在 $x = 1$ 取得极小值，
                    
                    又 $f\left( 1\right)  = \dfrac{1}{2} \times  {1}^{2} - \left( {a + 1}\right)  \times  1 + a\ln 1 =  - \dfrac{{2a} + 1}{2}$ ,
                    
                    所以 $- \dfrac{{2a} + 1}{2} > \left( {1 - \dfrac{1}{2}{{e}}^{2}}\right) a$ ,解得 $a > \dfrac{1}{{{e}}^{2} - 4}$ ,则 $\dfrac{1}{{{e}}^{2} - 4} < a < 1$ 

                    \item 当 $a > 1$ 时, $f\left( x\right)$ 在 $x = a$ 取得极小值,
                    
                    又 $f\left( a\right)  = \dfrac{1}{2} \times  {a}^{2} - \left( {a + 1}\right)  \times  a + a\ln a =  - \dfrac{{a}^{2}}{2} - a + a\ln a$ ,
                    
                    所以 $- \dfrac{{a}^{2}}{2} - a + a\ln a > \left( {1 - \dfrac{1}{2}{{e}}^{2}}\right) a$ ,即 $\ln a - \dfrac{a}{2} - 2 + \dfrac{1}{2}{{e}}^{2} > 0$ ,
                    
                    令 $h\left( a\right)  = \ln a - \dfrac{1}{2}a - 2 + \dfrac{1}{2}{{e}}^{2}\left( {a > 1}\right)$ ,则 ${h}^{\prime }\left( a\right)  = \dfrac{2 - a}{2a},h\left( a\right)  > 0$ ,
                    
                    令 ${h}^{\prime }\left( a\right)  > 0$ ,得 $1 < a < 2$ ; 令 ${h}^{\prime }\left( a\right)  < 0$ ,得 $a > 2$ ,
                    
                    所以函数 $h\left( a\right)$ 在 $\left( {1,2}\right)$ 上单调递增, $\left( {2, + \infty }\right)$ 上单调递减,
                    
                    且 $h\left( 1\right)  = \ln 1 - \dfrac{1}{2} \times  1 - 2 + \dfrac{1}{2}{{e}}^{2} = \dfrac{{{e}}^{2} - 5}{2} > 0,h\left( {{e}}^{2}\right)  = {\operatorname{lne}}^{2} - \dfrac{1}{2} \times  {{e}}^{2} - 2 + \dfrac{1}{2}{{e}}^{2} = 0$ ,
                    
                    所以由 $h\left( a\right)  > 0$ ,可得 $1 < a < {{e}}^{2}$ 
                \end{dingautolist}
                
                综上， $\dfrac{1}{{{e}}^{2} - 4} < a < 1$ 或 $1 < a < {{e}}^{2}$ ，即 $a$ 的取值范围是 $\left( {\dfrac{1}{{{e}}^{2} - 4},1}\right)  \cup  \left( {1,{{e}}^{2}}\right)$ .
            \end{enumerate}
        \end{jieda}
        
        \item 已知 $x_1,{x}_{2}$ 是函数 $f\left( x\right)  = \dfrac{1}{2}a{x}^{2} - {2x} + \ln x$ 的两个极值点,若不等式 $m > f\left( {x}_{1}\right)  + f\left( {x}_{2}\right)  + {x}_{1}{x}_{2}$ 恒成立，则实数 $m$ 的取值范围是\blkda{$\lbrack  - 2, + \infty )$}
        \begin{jieda}
            $f\left( x\right)  = \dfrac{1}{2}a{x}^{2} - {2x} + \ln x,x > 0$ ，则 ${f}^{\prime }\left( x\right)  = {ax} - 2 + \dfrac{1}{x} = \dfrac{a{x}^{2} - {2x} + 1}{x}$
            
            令 ${f}^{\prime }\left( x\right)  = 0$ 得， $a{x}^{2} - {2x} + 1 = 0$ ，
            
            由题意知 ${x}_{1},{x}_{2}$ 是方程 $a{x}^{2} - {2x} + 1 = 0$ 的两正根，
            
            则 $\left\{  \begin{matrix} \Delta  = 4 - {4a} > 0 \\  {x}_{1} + {x}_{2} = \dfrac{2}{a} > 0 \\  {x}_{1}{x}_{2} = \dfrac{1}{a} > 0 \end{matrix}\right.$ ，解得 $0 < a < 1$ ，且 ${x}_{1} + {x}_{2} = 2{x}_{1}{x}_{2}$ .
            
            由 $f\left( {x}_{1}\right)  + f\left( {x}_{2}\right)  + {x}_{1}{x}_{2} = \dfrac{1}{2}a{x}_{1}^{2} - 2{x}_{1} + \ln {x}_{1} + \dfrac{1}{2}a{x}_{2}^{2} - 2{x}_{2} + \ln {x}_{2} + {x}_{1}{x}_{2}$
            
            $= \dfrac{1}{2}\left( {2{x}_{1} - 1 + 2{x}_{2} - 1}\right)  - 2\left( {{x}_{1} + {x}_{2}}\right)  + \ln \left( {{x}_{1}{x}_{2}}\right)  + {x}_{1}{x}_{2}$
            
            $=  - \left( {{x}_{1} + {x}_{2}}\right)  + \ln {x}_{1}{x}_{2} + {x}_{1}{x}_{2} - 1$
            
            $=  - {x}_{1}{x}_{2} + \ln {x}_{1}{x}_{2} - 1 =  - \dfrac{1}{a} + \ln \dfrac{1}{a} - 1$ ,
            
            令 $t = \dfrac{1}{a},t > 1$ ,则 $g\left( t\right)  = \ln t - t - 1,t > 1$ ,
            
            由 ${g}^{\prime }\left( t\right)  = \dfrac{1}{t} - 1 < 0$ ,故 $g\left( t\right)$ 在 $\left( {1, + \infty }\right)$ 单调递减,故 $g\left( t\right)  < g\left( 1\right)  =  - 2$ ,
            
            要使 $m > f\left( {x}_{1}\right)  + f\left( {x}_{2}\right)  + {x}_{1}{x}_{2}$ 恒成立,即 $m > g\left( t\right)$ 恒成立,
            
            则 $m \geq   - 2$ ，则实数 $m$ 的取值范围是 $\lbrack  - 2, + \infty )$ .
        \end{jieda}
    \end{enumerate}

    \item 利用导数求函数的最值
    \begin{enumerate}
        \item 函数 $f\left( x\right)  = \left( {{x}^{2} - {4x} + 1}\right) {{e}}^{x}$ 在区间 $\left\lbrack  {-3,3}\right\rbrack$ 内的最大值和最小值分别为\blkda[3]{$\dfrac{6}{{e}}, - 2{{e}}^{3}$}
        \begin{jieda}
            $f\left( x\right)$ 的定义域为 $R,{f}^{\prime }\left( x\right)  = \left( {x - 3}\right) \left( {x + 1}\right) {{e}}^{x}$ ,
            
            当 $x \in  \left( {-1,3}\right)$ 时， ${f}^{\prime }\left( x\right)  < 0,f\left( x\right)$ 单调递减，
            
            当 $x \in  \left( {-3, - 1}\right)$ 时, ${f}^{\prime }\left( x\right)  > 0,f\left( x\right)$ 单调递增,
            
            所以 $f\left( x\right)$ 的最大值为 $f\left( {-1}\right)  = \dfrac{6}{{e}}$ ,
            
            又 $f\left( 3\right)  =  - 2{{e}}^{3},f\left( {-3}\right)  = {22}{{e}}^{-3}$ ,
            
            所以 $f\left( x\right)$ 的最大值为 $\dfrac{6}{{e}}$ ,最小值为 $- 2{{e}}^{3}$ .
        \end{jieda}
        \item 已知函数 $f\left( x\right)  = {{e}}^{x} - {e}x$ ，则 $f\left( x\right)$ \dotfill{(\kh{C})}
        \xx{有最小值 1 , 无最大值}{有最大值 1 , 无最小值}{有最小值 0 ，无最大值}{有最大值 0 , 无最小值}

        \begin{jieda}
            因为 $f\left( x\right)  = {{e}}^{x} - {e}x$ ,所以 ${f}^{\prime }\left( x\right)  = {{e}}^{x} - {e}$ .
            
            当 $x \in  \left( {-\infty ,1}\right)$ 时， ${f}^{\prime }\left( x\right)  < 0,f\left( x\right)$ 单调递减，
            
            当 $x \in  \left( {1, + \infty }\right)$ 时， ${f}^{\prime }\left( x\right)  > 0,f\left( x\right)$ 单调递增，
            
            故 $f\left( x\right)$ 的最小值为 $f\left( 1\right)  = 0$ ，无最大值.
        \end{jieda}
        
        \item 设$f(x) = -\dfrac{1}{3}x^3 + \dfrac{1}{2}x^2+2ax$
        \begin{enumerate}
            \item 若$y = f(x)$在$\left(\dfrac{2}{3}, +\infty\right)$上存在单调增区间，求$a$的取值范围
            \item 当$0 < a < 2$时，$y = f(x)$在区间$[1, 4]$上的最小值为$-\dfrac{16}{3}$，求$y = f(x)$在该区间上的最大值.
        \end{enumerate}
        \begin{jieda}
            \begin{enumerate}
                \item $f'(x) = -x^2+x+2a$，对称轴为$x = \dfrac{1}{2}$，故$y = f(x)$在$\left(\dfrac{2}{3}, +\infty\right)$上存在单调增区间即$f'\left(\dfrac{2}{3}\right) > 0 $，故$a \in \left(-\dfrac{1}{9}, +\infty\right)$
                \item $f'(x) = -x^2+x+2a$，对称轴为$x = \dfrac{1}{2}$,因为$a \in (0, 2)$，因此$f'(1) > 0, f'(4) < 0$，故$f(x)$在$[1, 4]$内先增后减，故最小值为$f(1)$或$f(4)$ \\
                $f(1) = 2a+\dfrac{1}{6}, f(4) = 8a-\dfrac{40}{3}$，$f(4)-f(1) = 6a-\dfrac{27}{2} < 0$，因此$f(x)$在$[1, 4]$上最小值为$f(4)$，故$8a-\dfrac{40}{3} = -\dfrac{16}{3}$，因此$a = 1$ \\
                $f'(x) = -x^2+x+2$，令$f'(x) = 0$，在$[1, 4]$内的解为$x = 2$，$f'(x)$在$[1, 2)$内为正，在$(2, 4]$内为负，因此$f(x)$最大值为$f(2) = \dfrac{10}{3}$
            \end{enumerate}
        \end{jieda}
        
        \item 已知函数 $g\left( x\right)  = {x}^{3} + {2m}{x}^{2} - {mx} + n$ 的图象在点 $( - 1,g( -$ 1))处的切线与直线 $x + {8y} - 2 = 0$ 垂直.
        \begin{enumerate}
            \item 求 $m$ 的值
            \item 已知 $g\left( x\right)$ 在区间 $\left\lbrack  {-1,2}\right\rbrack$ 上的最小值为 -5，求 $g\left( x\right)$ 在区间 $\left\lbrack  {-1,2}\right\rbrack$ 上的最大值.
        \end{enumerate}
        \begin{jieda}
            \begin{enumerate}
                \item 由已知,得 ${g}^{\prime }\left( x\right)  = 3{x}^{2} + {4mx} - m$ ,
                
                由题知 ${g}^{\prime }\left( {-1}\right)  = 3 - {4m} - m = 8$ ，解得 $m =  - 1$ .
                
                \item 由 (1) 可知, $g\left( x\right)  = {x}^{3} - 2{x}^{2} + x + n,{g}^{\prime }\left( x\right)  = 3{x}^{2} - {4x} + 1 = 3\left( {x - \dfrac{1}{3}}\right) \left( {x - 1}\right)$ ,
                
                $x,{g}^{\prime }\left( x\right) ,g\left( x\right)$ 的变化情况如表所示:
                
                \begin{center}
                \begin{tabular}{|c|c|c|c|c|c|c|c|}
                \hline
                $x$ & -1 & $\left( {-1,\dfrac{1}{3}}\right)$ & $\dfrac{1}{3}$ & $\left( \dfrac{1}{3}, 1\right)$ & 1 & $(1,2)$ & 2 \\
                \cline{1-8}
                ${g}^{\prime }\left( x\right)$ & \phantom{x} & + & 0 & - & 0 & + & \phantom{x} \\
                \cline{1-8}
                \hline
                $g\left( x\right)$ & $n -4$ & 单调递增 & 极大值 $n + \dfrac{4}{27}$  & 单调递减 & 极小值 $n$ & 单调递增 & $n + 2$  \\
                \cline{1-8}
                \hline
                \end{tabular}
                \end{center}
                
                $\because n - 4 < n,\therefore g{\left( x\right) }_{\min } = n - 4 =  - 5,\therefore n =  - 1$ ,
                
                $\because n + \dfrac{4}{27} < n + 2,\therefore g{\left( x\right) }_{\max } = n + 2 = 1$ .
                
                即 $g\left( x\right)$ 在区间 $\left\lbrack  {-1,2}\right\rbrack$ 上的最大值为 1 .
            \end{enumerate}
        \end{jieda}
        
        \item 已知函数 $f\left( x\right)  = {{e}}^{2x} + {{e}}^{x} - x$ .
        \begin{enumerate}
            \item 求曲线 $y = f\left( x\right)$ 在点 $\left( {0,f\left( 0\right) }\right)$ 处的切线方程
            \item 当 $x \in  \left\lbrack  {-1,0}\right\rbrack$ 时，求函数 $f\left( x\right)$ 的最大值与最小值
        \end{enumerate}
        \begin{jieda}
            \begin{enumerate}
                \item 因为 $f\left( x\right)  = {{e}}^{2x} + {{e}}^{x} - x$ ,则 ${f}^{\prime }\left( x\right)  = 2{{e}}^{2x} + {{e}}^{x} - 1$ ,
                
                可得 $f\left( 0\right)  = 2,{f}^{\prime }\left( 0\right)  = 2$ ,
                
                即切点坐标为 $\left( {0,2}\right)$ ,切线斜率为 $k = 2$ ,
                
                所以切线方程为 $y = {2x} + 2$ .
                \item 由 (1) 可得 ${f}^{\prime }\left( x\right)  = 2{{e}}^{2x} + {{e}}^{x} - 1 = \left( {{{e}}^{x} + 1}\right) \left( {2{{e}}^{x} - 1}\right)$ ,
                
                且 $x \in  \left\lbrack  {-1,0}\right\rbrack$ ,则 ${{e}}^{x} + 1 > 0$ ,
                
                令 ${f}^{\prime }\left( x\right)  > 0$ ，则 $2{{e}}^{x} - 1 > 0$ ，解得 $- \ln 2 < x \leq  0$ ；
                
                令 ${f}^{\prime }\left( x\right)  < 0$ ，则 $2{{e}}^{x} - 1 < 0$ ，解得 $- 1 \leq  x <  - \ln 2$ ；
                
                可知 $f\left( x\right)$ 在 $\lbrack  - 1, - \ln 2)$ 内单调递减,在 $( - \ln 2,0\rbrack$ 内单调递增,
                
                又因为 $f\left( {-1}\right)  = \dfrac{1 + {e} - {{e}}^{2}}{{{e}}^{2}},f\left( 0\right)  = 2,f\left( {-\ln 2}\right)  = \dfrac{3}{4} + \ln 2$ ,且 $\dfrac{1 + {e} - {{e}}^{2}}{{{e}}^{2}} < 2$ ,
                
                所以函数 $f\left( x\right)$ 的最大值为 2,最小值 $\dfrac{3}{4} + \ln 2$ .
            \end{enumerate}
        \end{jieda} 
    \end{enumerate}
    \item 已知函数最值求参数
    \begin{enumerate}
        \item 已知定义在 $\left( {a,2 - {3a}}\right)$的函数 $f\left( x\right)  = x + \dfrac{4}{x} + 3\ln x$ 有最小值，则实数 $a$ 的取值范围是\blkda{$\left\lbrack  {0,\dfrac{1}{3}}\right)$}
        \begin{jieda}
            函数 $f\left( x\right)  = x + \dfrac{4}{x} + 3\ln x$ 的定义域为 $\left( {0, + \infty }\right)$
            
            ${f}^{\prime }\left( x\right)  = 1 - \dfrac{4}{{x}^{2}} + \dfrac{3}{x} = \dfrac{{x}^{2} + {3x} - 4}{{x}^{2}} = \dfrac{\left( {x + 4}\right) \left( {x - 1}\right) }{{x}^{2}},$
            
            令 ${f}^{\prime }\left( x\right)  = 0$ 可得 $x = 1$ 或 $x =  - 4$ (舍),
            
            当 $0 < x < 1$ 时 ${f}^{\prime }\left( x\right)  < 0$ ,当 $x > 1$ 时 ${f}^{\prime }\left( x\right)  > 0$ ,
            
            所以 $f\left( x\right)$ 在 $\left( {0,1}\right)$ 上单调递减,在 $\left( {1, + \infty }\right)$ 上单调递增,所以 $f\left( x\right)$ 在 $x = 1$ 处取得极小值,即最小值，又因为函数 $f\left( x\right)$ 在 $\left( {a,2 - {3a}}\right)$ 内有最小值，故 $0 \leq  a < 1 < 2 - {3a}$ ，解得 $0 \leq  a < \dfrac{1}{3}$
        \end{jieda}
        
        \item 已知 $f\left( x\right)  = \dfrac{1}{3}{x}^{3} - x$ 在区间 $\left( {m,6 - {m}^{2}}\right)$ 上有最小值，则实数 $m$ 的取值范围是\blkda{$\left\lbrack  {-2,1}\right)$}
        \begin{jieda}
            由函数 $f\left( x\right)  = \dfrac{1}{3}{x}^{3} - x$ ,可得 ${f}^{\prime }\left( x\right)  = {x}^{2} - 1 = \left( {x + 1}\right) \left( {x - 1}\right)$

            \begin{dingautolist}{172}
                \item 当 $x <  - 1$ 时， ${f}^{\prime }\left( x\right)  > 0,f\left( x\right)$ 单调递增
                
                \item 当 $- 1 < x < 1$ 时, ${f}^{\prime }\left( x\right)  < 0,f\left( x\right)$ 单调递减
                
                \item 当 $x > 1$ 时, ${f}^{\prime }\left( x\right)  > 0,f\left( x\right)$ 单调递增
            \end{dingautolist}
            
            要使得函数 $y = f\left( x\right)$ 在区间 $\left( {m,6 - {m}^{2}}\right)$ 上有最小值，则满足 $\left\{  \begin{matrix} m < 1 < 6 - {m}^{2} \\  f\left( m\right)  \geq  f\left( 1\right)  \end{matrix}\right.$ ,即 $\left\{  \begin{array}{l}  - \sqrt{5} < m < 1 \\  \dfrac{1}{3}{m}^{3} - m \geq   - \dfrac{2}{3} \end{array}\right.$
            
            因为 $\dfrac{1}{3}{m}^{3} - m \geq   - \dfrac{2}{3}$ ,可得 ${m}^{3} - {3m} + 2 \geq  0$ ,即 ${\left( m - 1\right) }^{2}\left( {x + 2}\right)  \geq  0$ ,解得 $m \geq   - 2$，所以 $- 2 \leq  m < 1$ ，即实数 $m$ 的取值为 $\lbrack  - 2,1)$ .
        \end{jieda}
        
        \item 已知函数 $f\left( x\right)  = \dfrac{x}{a} + \ln x$ ，其中 $a$ 为常数， ${e}$ 为自然对数的底数.
        \begin{enumerate}
            \item 当 $a =  - 1$ 时，求 $f\left( x\right)$ 的单调区间
            \item 若 $f\left( x\right)$ 在区间 $(0,{e}\rbrack$ 上的最大值为 2，求 $a$ 的值
        \end{enumerate}

        \begin{jieda}
            \begin{enumerate}
                \item 函数 $f\left( x\right)$ 的定义域为 $\left( {0, + \infty }\right)$ ,
                
                当 $a =  - 1$ 时， $f\left( x\right)  = \ln x - x,{f}^{\prime }\left( x\right)  = \dfrac{1 - x}{x}\left( {x > 0}\right)$ ,
                
                令 ${f}^{\prime }\left( x\right)  > 0$ ,得 $0 < x < 1$ ; 令 ${f}^{\prime }\left( x\right)  < 0$ ,得 $x > 1$ ,
                
                $\therefore$ 函数 $f\left( x\right)$ 的单调递增区间为 $\left( {0,1}\right)$ ,单调递减区间为 $\left( {1, + \infty }\right)$ .
                
                \item ${f}^{\prime }\left( x\right)  = \dfrac{a + x}{ax}$

                \begin{dingautolist}{172}
                    \item 当 $a > 0$ 时， $\because x > 0$ ， $\therefore {f}^{\prime }\left( x\right)  > 0$ ，
                    
                    $\therefore$ 函数 $f\left( x\right)$ 在 $\left( {0,{e}}\right)$ 上单调递增，
                    
                    $\therefore f{\left( x\right) }_{\max } = f\left( {e}\right)  = 2,\therefore \dfrac{{e}}{a} + 1 = 2,$
                    
                    $\therefore a = {e}$ ,符合题意
                    
                    \item 当 $a < 0$ 且 $- a < {e}$ ，即 $- {e} < a < 0$ 时，令 ${f}^{\prime }\left( x\right)  = 0$ ，得 $x =  - a$ ，
                    
                    当 $x$ 变化时， ${f}^{\prime }\left( x\right) ,f\left( x\right)$ 的变化情况如下表.
                    
                    \begin{center}
                    \begin{tabular}{|c|c|c|c|}
                    \hline
                    $x$ & $\left( {0, - a}\right)$ & $- a$ & $（-a,e)$ \\
                    \cline{1-4}
                    ${f}^{\prime }(x)$ & + & 0 & - \\
                    \cline{1-4}
                    $f\left( x\right)$ & 单调递增 & 极大值 & 单调递减 \\
                    \cline{1-4}
                    \hline
                    \end{tabular}
                    \end{center}
                    
                    $\therefore f{\left( x\right) }_{\max } = f\left( {-a}\right)  = 2,\therefore  - 1 + \ln \left( {-a}\right)  = 2$ ,
                    
                    $\therefore a =  - {{e}}^{3}$ ,不符合题意,舍去;
                
                    \item 当 $- a \geq  {e}$ ，即 $a \leq   - {e}$ 时，在 $(0,{e}\rbrack$ 上， ${f}^{\prime }\left( x\right)  \geq  0$ ，
                
                    $\therefore f\left( x\right)$ 在 $(0,{e}\rbrack$ 上单调递增，故 $f\left( x\right)$ 在 $( 0,{e}\rbrack$ 上的最大值为 $f\left( {e}\right)  = \dfrac{{e}}{a} + 1 = 2$，不符合题意,舍去.
                
                \end{dingautolist}
                
                综上可得 $a = {e}$ .
            \end{enumerate}
        \end{jieda}
        
        \item 设函数 $f\left( x\right)  = \ln x - {ax}\left( {a \in  R}\right)$ .
        \begin{enumerate}
            \item 讨论函数 $f\left( x\right)$ 的单调性
            \item 当函数 $f\left( x\right)$ 有最大值，且最大值小于 $a - 2$ 时，求 $a$ 的取值范围
        \end{enumerate}
        \begin{jieda}
            \begin{enumerate}
                \item 由 $f\left( x\right)  = \ln x - {ax}$ ，知 ${f}^{\prime }\left( x\right)  = \dfrac{1}{x} - a$ ，定义域为 $\left( {0, + \infty }\right)$ ，
                
                当 $a \leq  0$ 时， ${f}^{\prime }\left( x\right)  > 0$ 恒成立，所以 $f\left( x\right)$ 在 $\left( {0, + \infty }\right)$ 上单调递增；
                
                当 $a > 0$ 时,令 ${f}^{\prime }\left( x\right)  > 0$ ,则 $0 < x < \dfrac{1}{a},f\left( x\right)$ 在 $\left( {0,\dfrac{1}{a}}\right)$ 上单调递增;
                
                令 ${f}^{\prime }\left( x\right)  < 0$ ，则 $x > \dfrac{1}{a},f\left( x\right)$ 在 $\left( {\dfrac{1}{a}, + \infty }\right)$ 上单调递减；
                
                综上所述，当 $a \leq  0$ 时， $f\left( x\right)$ 在 $\left( {0, + \infty }\right)$ 上单调递增；
                
                当 $a > 0$ 时， $f\left( x\right)$ 在 $\left( {0,\dfrac{1}{a}}\right)$ 上单调递增，在 $\left( {\dfrac{1}{a}, + \infty }\right)$ 上单调递减.

                \item 由(1)知，若 $f\left( x\right)$ 有最大值，则 $a > 0$ ，且 $f{\left( x\right) }_{\max } = f\left( \dfrac{1}{a}\right)  =  - \ln a - 1$ ，
                
                因为 $f\left( x\right)$ 的最大值小于 $a - 2$ ,
                
                所以 $- \ln a - 1 < a - 2$ ,即 $a + \ln a - 1 > 0$ ,
                
                设 $g\left( a\right)  = a + \ln a - 1$ ，问题转化为解不等式 $g\left( a\right)  > 0$ ，
                
                因为 ${g}^{\prime }\left( a\right)  = 1 + \dfrac{1}{a} > 0$ 恒成立,所以 $g\left( a\right)$ 在 $a \in  \left( {0, + \infty }\right)$ 上单调递增,
                
                又 $g\left( 1\right)  = 0$ ，所以 $g\left( a\right)  > 0 = g\left( 1\right)$ ，所以 $a > 1$
            \end{enumerate}
        \end{jieda}
    \end{enumerate}

    \item 由函数的单调性求参数
    \begin{enumerate}
        \item 已知函数 $f\left( x\right)  = e^x - ax^2$在$(0, +\infty)$单调递增,则实数 $a$ 的取值范围是\blkda{$\left(-\infty, \dfrac{e}{2}\right]$}
        \begin{jieda}
            ${f}^{\prime }\left( x\right)  = e^x - 2ax \geq  0$在$(0, +\infty)$恒成立
            \begin{dingautolist}{172}
                \item 直曲分离：$y = 2ax$恒在$y = e^x$下方，相切为边界，$y = e^x$过原点的切线为$y = ex$，故$2a \in (0, e]$，因此$a \in \left(-\infty, \dfrac{e}{2}\right]$
                \item 参变分离：$2a \leq \dfrac{e^x}{x}$，设$g(x) = \dfrac{e^x}{x}$，$g'(x) = \dfrac{e^x(x-1)}{x^2}$，故$g(x)$在$(0, 1)$严减，在$(1, +\infty)$严增，故$2a \leq e$，因此$a \in \left(-\infty, \dfrac{e}{2}\right]$
                \item 含参求导：$f''(x) = e^x - 2a$
                
                若$a \leq 0$，则$f'(x)$严增，$f'(x) \geq f'(0) = 1$满足要求

                若$a > 0$，则$f'(x)$在$(0, ln(2a))$严减，在$(ln(2a), +\infty)$严增，故$f'(x) \geq f'(ln(2a)) = 2a - 2aln(2a) \geq 0$，故$a \leq \dfrac{e}{2}$

                综上，$a \in \left(-\infty, \dfrac{e}{2}\right]$
            \end{dingautolist}
        \end{jieda}
        \item 已知函数 $f\left( x\right)  = \dfrac{1}{2}{x}^{2} + {kx} + \ln x$ 在区间 $\left( {3, + \infty }\right)$ 上单调递增,则实数 $k$ 的取值范围是\blkda{${k} \geq - \dfrac{10}{3}$}
        \begin{jieda}
            由于函数 $f\left( x\right)  = \dfrac{1}{2}{x}^{2} + {kx} + \ln x$ 在区间 $\left( {3, + \infty }\right)$ 上单调递增,
            所以 ${f}^{\prime }\left( x\right)  = x + k + \dfrac{1}{x} \geq  0$ 在 $\left( {3, + \infty }\right)$ 上恒成立,即 $x + \dfrac{1}{x} \geq   - k$ 恒成立,
            由 ${y} = {x} + \dfrac{1}{{x}}$ 在 $\left( {3, + \infty }\right)$ 上单调递增,则 ${x} + \dfrac{1}{{x}} > 3 + \dfrac{1}{3} = \dfrac{10}{3}$ ,
            $\therefore  - {k} \leq  \dfrac{10}{3} \Longrightarrow  {k} \geq   - \dfrac{10}{3}$ .
            所以实数 ${k}$ 的取值范围为 $\left\lbrack  {-\dfrac{10}{3}, + \infty }\right)$ .
        \end{jieda}
        \item 若函数 $f\left( x\right)  = \ln x + {x}^{2} - {ax}$ 在其定义域内单调递增，则实数 $a$ 的取值范围是\blkda{$( - \infty ,2\sqrt{2}\rbrack$}
        \begin{jieda}
            $f\left( x\right)$ 的定义域为 $\left( {0, + \infty }\right) ,{f}^{\prime }\left( x\right)  = \dfrac{1}{x} + {2x} - a$ ,
            
            因为函数 $f\left( x\right)  = \ln x + {x}^{2} - {ax}$ 在其定义域内单调递增,
            
            所以 $\dfrac{1}{x} + {2x} - a \geq  0$ 在 $\left( {0, + \infty }\right)$ 上恒成立,即 $\dfrac{1}{x} + {2x} \geq  a$ 在 $\left( {0, + \infty }\right)$ 上恒成立,
            
            因为 $\dfrac{1}{x} + {2x} \geq  2\sqrt{\dfrac{1}{x} \cdot  {2x}} = 2\sqrt{2}$ ，当且仅当 $x = \dfrac{\sqrt{2}}{2}$ 时，等号成立，
            
            所以 ${\left( \dfrac{1}{x} + 2x\right) }_{\min } = 2\sqrt{2}$ ,所以 $a \leq  2\sqrt{2}$ .
        \end{jieda}
        
        \item 已知函数 $F\left( x\right)  = \left\{  \begin{matrix}  - \dfrac{1}{3}{x}^{3} + a{x}^{2} - a - 4\left( {x \geq  0}\right) \\  {ax} - \sin x\left( {x < 0}\right)  \end{matrix}\right.$ 在 ${R}$ 上单调, 则实数 $a$ 的取值范围为\blkda{$\left\lbrack  {-4, - 1}\right\rbrack$}
        \begin{jieda}
            设 $f\left( x\right)  =  - \dfrac{1}{3}{x}^{3} + a{x}^{2} - a - 4$ ,则 $f\left( x\right)$ 在 $\lbrack 0, + \infty )$ 上单调递减.
            
            因为 ${f}^{\prime }\left( x\right)  =  - {x}^{2} + {2ax}$ ,由 ${f}^{\prime }\left( x\right)  \leq  0$ 在 $\lbrack 0, + \infty )$ 上恒成立,得:
            
            若 $x = 0$ ,则 $0 \leq  0$ ,所以 $a \in  {R}$ ;
            
            若 $x > 0$ ，则 $- {x}^{2} + {2ax} \leq  0 \Longrightarrow  {2a} \leq  x$ ，所以 $a \leq  0$ .
            
            设 $g\left( x\right)  = {ax} - \sin x$ ,则 $g\left( x\right)$ 在 $\left( {-\infty ,0}\right)$ 上单调递减.
            
            由 ${g}^{\prime }\left( x\right)  = a - \cos x \leq  0$ 在 $\left( {-\infty ,0}\right)$ 上恒成立,所以 $a \leq  \cos x,x \in  \left( {-\infty ,0}\right)$ ,所以 $a \leq   - 1$ .
            
            且 $- a - 4 \leq  0 \Longrightarrow  a \geq   - 4$ .
        \end{jieda}
        
        \item 已知函数 $f\left( x\right)  = \left( {{x}^{2} - {2x} + a}\right) {{e}}^{x}$ .
        \begin{enumerate}
            \item $f\left( x\right)$ 在 $\left\lbrack  {1,5}\right\rbrack$ 上是增函数，求 $a$ 的取值范围
            \item 讨论函数 $f\left( x\right)$ 的单调性
        \end{enumerate}
        \begin{jieda}
            \begin{enumerate}
                \item 因为 $f\left( x\right)  = \left( {{x}^{2} - {2x} + a}\right) {{e}}^{x}$ ,所以 $f\left( x\right)$ 的定义域为 ${R}$，则 ${f}^{\prime }\left( x\right)  = \left( {{2x} - 2}\right) {{e}}^{x} + \left( {{x}^{2} - {2x} + a}\right) {{e}}^{x} = \left( {{x}^{2} + a - 2}\right) {{e}}^{x}$
                
                因为 $f\left( x\right)$ 在 $\left\lbrack  {1,5}\right\rbrack$ 上是增函数,即 ${f}^{\prime }\left( x\right)  \geq  0$ 在 $\left\lbrack  {1,5}\right\rbrack$ 上恒成立，则 $\left( {{x}^{2} + a - 2}\right) {{e}}^{x} \geq  0$ 在 $\left\lbrack  {1,5}\right\rbrack$ 上恒成立,
                
                因为 ${{e}}^{x} > 0$ 在 $\left\lbrack  {1,5}\right\rbrack$ 上恒成立,所以 ${x}^{2} + a - 2 \geq  0$ 在 $\left\lbrack  {1,5}\right\rbrack$ 上恒成立，即 $a - 2 \geq   - {x}^{2}$ 在 $\left\lbrack  {1,5}\right\rbrack$ 上恒成立,即 $a - 2 \geq  {\left( -{x}^{2}\right) }_{\max }$
                
                因为 $1 \leq  x \leq  5$ ,所以 $1 \leq  {x}^{2} \leq  {25}$ ,则 $- {25} \leq   - {x}^{2} \leq   - 1$ ,
                
                所以 $a - 2 \geq   - 1$ ，则 $a \geq  1$ .

                \item 由 (1) 得 ${f}^{\prime }\left( x\right)  = \left( {{x}^{2} + a - 2}\right) {{e}}^{x}$
                
                当 $a \geq  2$ 时， ${f}^{\prime }\left( x\right)  \geq  0$ ，则 $f\left( x\right)$ 在 $R$ 上是增函数；
                
                当 $a < 2$ 时， ${f}^{\prime }\left( x\right)  = \left\lbrack  {{x}^{2} - \left( {2 - a}\right) }\right\rbrack  {{e}}^{x} = \left( {x + \sqrt{2 - a}}\right) \left( {x - \sqrt{2 - a}}\right) {{e}}^{x}$ ,
                
                所以 ${f}^{\prime }\left( x\right)  = 0 \Longleftrightarrow  x =  \pm  \sqrt{2 - a}$ ;
                
                ${f}^{\prime }\left( x\right)  > 0 \Longleftrightarrow  x <  - \sqrt{2 - a}$ 或 $x > \sqrt{2 - a}$ ;
                
                ${f}^{\prime }\left( x\right)  < 0 \Longleftrightarrow   - \sqrt{2 - a} < x < \sqrt{2 - a},$
                
                所以 $f\left( x\right)$ 在 $\left( {-\sqrt{2 - a},\sqrt{2 - a}}\right)$ 上是减函数,在 $\left( {-\infty , - \sqrt{2 - a}}\right)$ 和 $\left( {\sqrt{2 - a}, + \infty }\right)$ 上是增函数.
            \end{enumerate}
        \end{jieda}
        \item 已知函数 $f\left( x\right)  = \dfrac{1}{3}{x}^{3} + a{x}^{2} - {4x}$ .
        \begin{enumerate}
            \item 当 $a =  - 1$ 时，求曲线 $y = f\left( x\right)$ 在点 $\left( {3,f\left( 3\right) }\right)$ 处的切线方程
            \item 试问是否存在实数 $a$ ，使得 $f\left( x\right)$ 在 $\left\lbrack  {1,a}\right\rbrack$ 上单调递增？若存在，求 $a$ 的取值范围；若不存在，请说明理由.
        \end{enumerate}
        \begin{jieda}
            \begin{enumerate}
                \item 当 $a =  - 1$ 时， $f\left( x\right)  = \dfrac{1}{3}{x}^{3} - {x}^{2} - {4x}$ ，则 ${f}^{\prime }\left( x\right)  = {x}^{2} - {2x} - 4$ ，所以 ${f}^{\prime }\left( 3\right)  =  - 1$ ，因为 $f\left( 3\right)  =  - {12}$ ，
                
                所以曲线 $y = f\left( x\right)$ 在点 $\left( {3,f\left( 3\right) }\right)$ 处的切线方程为 $y + {12} =  - \left( {x - 3}\right)$，即 $y =  - x - 9\left( {\text{ 或 }x + y + 9 = 0}\right)$ .
                
                
                \item 假设存在实数 $a$ ，使得 $f\left( x\right)$ 在 $\left\lbrack  {1,a}\right\rbrack$ 上单调递增，则 ${f}^{\prime }\left( x\right)  = {x}^{2} + {2ax} - 4 \geq  0$ 对 $x \in  \left\lbrack  {1,a}\right\rbrack$ 恒成立，即 $a \geq  \dfrac{2}{x} - \dfrac{x}{2}$ 对 $x \in  \left\lbrack  {1,a}\right\rbrack$ 恒成立.
                
                当 $x \in  \left\lbrack  {1,a}\right\rbrack$ 时， $y = \dfrac{2}{x} - \dfrac{x}{2}$ 为减函数，则 ${\left( \dfrac{2}{x} - \dfrac{x}{2}\right) }_{\max } = \dfrac{2}{1} - \dfrac{1}{2} = \dfrac{3}{2}$ .
                
                所以 $a \geq  \dfrac{3}{2}$ ，又 $a > 1$ ，所以 $a$ 的取值范围为 $\left\lbrack  {\dfrac{3}{2}, + \infty }\right)$ .
            \end{enumerate}
        \end{jieda}
    \end{enumerate}

    \item 恒成立问题求参数
    \begin{enumerate}
        
        \item 已知函数 $f\left( x\right)  = e^x - 2ax \geq  0$在$(0, +\infty)$恒成立,则实数 $a$ 的取值范围是\blkda{$\left(-\infty, \dfrac{e}{2}\right]$}
        \begin{jieda}
            \begin{dingautolist}{172}
                \item 直曲分离：$y = 2ax$恒在$y = e^x$下方，相切为边界，$y = e^x$过原点的切线为$y = ex$，故$2a \in (0, e]$，因此$a \in \left(-\infty, \dfrac{e}{2}\right]$
                \item 参变分离：$2a \leq \dfrac{e^x}{x}$，设$g(x) = \dfrac{e^x}{x}$，$g'(x) = \dfrac{e^x(x-1)}{x^2}$，故$g(x)$在$(0, 1)$严减，在$(1, +\infty)$严增，故$2a \leq e$，因此$a \in \left(-\infty, \dfrac{e}{2}\right]$
                \item 含参求导：$f''(x) = e^x - 2a$
                
                若$a \leq 0$，则$f'(x)$严增，$f'(x) \geq f'(0) = 1$满足要求

                若$a > 0$，则$f'(x)$在$(0, ln(2a))$严减，在$(ln(2a), +\infty)$严增，故$f'(x) \geq f'(ln(2a)) = 2a - 2aln(2a) \geq 0$，故$a \leq \dfrac{e}{2}$

                综上，$a \in \left(-\infty, \dfrac{e}{2}\right]$
            \end{dingautolist}
        \end{jieda}
        
        \item 已知函数 $f\left( x\right)  = {2a}\left( {x - 1}\right) {{e}}^{x} - {x}^{2}\left( {a > 0}\right)$ .
        \begin{enumerate}
            \item 当 $0 < a < 1$ 时，讨论 $f\left( x\right)$ 的单调性
            \item 若不等式 $f\left( x\right)  > x - {4a}{{e}}^{x}$ 对 $x \in  \left( {-1, + \infty }\right)$ 恒成立，求 $a$ 的取值范围
        \end{enumerate}

        \begin{jieda}
            \begin{enumerate}
                \item 由题意得 ${f}^{\prime }\left( x\right)  = {2a}{{e}}^{x} + {2a}\left( {x - 1}\right) {{e}}^{x} - {2x} = {2ax}{{e}}^{x} - {2x} = {2x}\left( {a{{e}}^{x} - 1}\right)$ ,
                
                因为 $0 < a < 1$ 时,令 ${f}^{\prime }\left( x\right)  = 0$ ,得 $x = \ln \dfrac{1}{a} =  - \ln a > 0$ 或 $x = 0$ ,
                
                当 $x < 0$ 时， ${f}^{\prime }\left( x\right)  > 0,f\left( x\right)$ 在 $\left( {-\infty ,0}\right)$ 上单调递增，
                
                当 $0 < x <  - \ln a$ 时， ${f}^{\prime }\left( x\right)  < 0,f\left( x\right)$ 在 $\left( {0, - \ln a}\right)$ 上单调递减，
                
                当 $x >  - \ln a$ 时, ${f}^{\prime }\left( x\right)  > 0,f\left( x\right)$ 在 $\left( {-\ln a, + \infty }\right)$ 上单调递增,
                
                故当 $0 < a < 1$ 时， $f\left( x\right)$ 在 $\left( {-\infty ,0}\right)$ 上单调递增，在 $\left( {0, - \ln a}\right)$ 上单调递减，在 $\left( {-\ln a, + \infty }\right)$ 上单调递增.

                \item 由 $f\left( x\right)  > x - {4a}{{e}}^{x}$ ，得到 $\left( {x + 1}\right) \left( {{2a}{{e}}^{x} - x}\right)  > 0$ ，
                
                因为 $x >  - 1$ ，所以 $x + 1 > 0$ ，则 ${2a} > {\left( \dfrac{x}{{{e}}^{x}}\right) }_{\max }$ ，
                
                令 $g\left( x\right)  = \dfrac{x}{{{e}}^{x}}\left( {x >  - 1}\right)$ ，则 ${g}^{\prime }\left( x\right)  = \dfrac{1 - x}{{{e}}^{x}}$ ，
                
                当 $- 1 < x < 1$ 时， ${g}^{\prime }\left( x\right)  > 0$ ，即 $g\left( x\right)$ 在区间 $\left( {-1,1}\right)$ 上单调递增，
                
                当 $x > 1$ 时， ${g}^{\prime }\left( x\right)  < 0$ ，即 $g\left( x\right)$ 在区间 $\left( {1, + \infty }\right)$ 上单调递减，所以 $g\left( x\right)  \leq  g\left( 1\right)  = \dfrac{1}{{e}}$ ，
                
                得到 ${2a} > \dfrac{1}{{e}}$ ,所以 $a > \dfrac{1}{2{e}}$ ,故 $a$ 的取值范围为 $\left( {\dfrac{1}{2{e}}, + \infty }\right)$ .
            \end{enumerate}
        \end{jieda}
        \item 已知函数 $f\left( x\right)  = x - \ln x - 2$ .
        \begin{enumerate}
            \item 求曲线 $y = f\left( x\right)$ 在点 $\left( {1,f\left( 1\right) }\right)$ 处的切线方程
            \item 函数 $f\left( x\right)$ 在区间 $\left( {k,k + 1}\right) \left( {k \in  \mathbf{N}}\right)$ 上有零点,求 $k$ 的值
            \item 记函数 $g\left( x\right)  = \dfrac{1}{2}{x}^{2} - {bx} - 2 - f\left( x\right)$ ,设 ${x}_{1},{x}_{2}\left( {{x}_{1} < {x}_{2}}\right)$ 是函数 $g\left( x\right)$ 的两个极值点,若 $b \geq  \dfrac{3}{2}$ ,且$g\left( {x}_{1}\right)  - g\left( {x}_{2}\right)  \geq  k$ 恒成立,求实数 $k$ 的最大值.
        \end{enumerate}
        \begin{jieda}
            \begin{enumerate}
                \item 因为 $f\left( x\right)  = x - \ln x - 2$ ,所以 ${f}^{\prime }\left( x\right)  = 1 - \dfrac{1}{x}$ ,则切线斜率为 ${f}^{\prime }\left( 1\right)  = 0$ ,又 $f\left( 1\right)  =  - 1$ ,切点为(1, - 1),所以切线方程为 $y =  - 1$ ;
                \item $\because {f}^{\prime }\left( x\right)  = \dfrac{x - 1}{x},x \in  \left( {0, + \infty }\right)$ ,当 $0 < x < 1$ 时, ${f}^{\prime }\left( x\right)  < 0$ ,函数 $f\left( x\right)$ 单调递减；当 $x > 1$ 时, ${f}^{\prime }\left( x\right)  > 0$ ,函数 $f\left( x\right)$ 单调递增,所以 $f\left( x\right)$ 的极小值为 $f\left( 1\right)  =  - 1 < 0,f\left( {{e}}^{-2}\right)  = {{e}}^{-2} - \ln {{e}}^{-2} - 2 = {{e}}^{-2} > 0$ ,
                $\therefore f\left( x\right)$ 在区间(0,1)上存在一个零点 ${x}_{3}$ ,此时 $k = 0$ ;
                又 $f\left( 3\right)  = 3 - \ln 3 - 2 = 1 - \ln 3 < 0$ , $f\left( 4\right)  = 4 - \ln 4 - 2 = 2 - 2\ln 2 = 2\left( {1 - \ln 2}\right)  > 0$ ,
                $\therefore f\left( x\right)$ 在区间(3,4)上存在一个零点 ${x}_{4}$ ,此时 $k = 3$ ,
                综上, $k$ 的值为 0 或 3 ;
                \item 函数 $g\left( x\right)  = \dfrac{1}{2}{x}^{2} - {bx} - 2 - f\left( x\right)  = \ln x + \dfrac{1}{2}{x}^{2} - \left( {b + 1}\right) x,x \in  \left( {0, + \infty }\right)$ ,
                所以 ${g}^{\prime }\left( x\right)  = \dfrac{1}{x} + x - \left( {b + 1}\right)  = \dfrac{{x}^{2} - \left( {b + 1}\right) x + 1}{x}$ ,
                由 ${g}^{\prime }\left( x\right)  = 0$ 得 ${x}^{2} - \left( {b + 1}\right) x + 1 = 0$ ,依题意方程 ${x}^{2} - \left( {b + 1}\right) x + 1 = 0$ 有两不相等的正实根 ${x}_{1}\text{、}{x}_{2}$，根据韦达定理有$\left\{\begin{array}{l}
                     {x}_{1} + {x}_{2} = b + 1  \\
                     {x}_{1}{x}_{2} = 1
                \end{array} \right.$ ,
                又 $b \geq  \dfrac{3}{2},{x}_{1} + \dfrac{1}{{x}_{1}} = b + 1 \geq  \dfrac{5}{2},0 < {x}_{1} < {x}_{2} = \dfrac{1}{{x}_{1}}$ ,解得 $0 < {x}_{1} \leq  \dfrac{1}{2}$ ,
                $\therefore g\left( {x}_{1}\right)  - g\left( {x}_{2}\right)  = \ln \dfrac{{x}_{1}}{{x}_{2}} + \dfrac{1}{2}\left( {{x}_{1}^{2} - {x}_{2}^{2}}\right)  - \left( {b + 1}\right) \left( {{x}_{1} - {x}_{2}}\right)  = 2\ln {x}_{1} - \dfrac{1}{2}\left( {{x}_{1}^{2} - \dfrac{1}{{x}_{1}^{2}}}\right)$ ,
                构造函数 $F\left( x\right)  = 2\ln x - \dfrac{1}{2}\left( {{x}^{2} - \dfrac{1}{{x}^{2}}}\right) ,x \in  \left( {0,\dfrac{1}{2}}\right\rbrack$ ,
                所以 ${F}^{\prime }\left( x\right)  = \dfrac{2}{x} - x - \dfrac{1}{{x}^{3}} = \dfrac{-{\left( {x}^{2} - 1\right) }^{2}}{{x}^{3}} < 0,\therefore F\left( x\right)$ 在 $\left( {0,\dfrac{1}{2}}\right\rbrack$ 上单调递减,
                所以当 $x = \dfrac{1}{2}$ 时, $F{\left( x\right) }_{\min } = F\left( \dfrac{1}{2}\right)  = \dfrac{15}{8} - 2\ln 2$ ,
                因为 $g\left( {x}_{1}\right)  - g\left( {x}_{2}\right)  \geq  k$ 恒成立,所以 $k \leq  \dfrac{15}{8} - 2\ln 2$ ,则 $k$ 的最大值为 $\dfrac{15}{8} - 2\ln 2$ .
            \end{enumerate}
        \end{jieda}
    \end{enumerate}

    \item 利用导数证明不等式/求解不等式恒成立
    \begin{enumerate}
        \item 已知函数 $f\left( x\right)  = {{e}}^{x} + {x}^{2} - \left( {{ax} + 1}\right) ,x \in  {R}$ 的图象在 $x = 0$ 处的切线方程为 $y =  - x$ .
        \begin{enumerate}
            \item 求$a$的值
            \item 证明: $f\left( x\right)  > {x}^{2} - 1$
        \end{enumerate}
        \begin{jieda}
            \begin{enumerate}
                \item ${f}^{\prime }\left( x\right)  = {{e}}^{x} + {2x} - a$ ，则 ${f}^{\prime }\left( 0\right)  = 1 - a,f\left( 0\right)  = 0$，则过切点 $\left( {0,0}\right)$ 的切线方程为 $y - 0 = \left( {1 - a}\right) \left( {x - 0}\right)$，即 $y = \left( {1 - a}\right) x$ ,所以 $1 - a =  - 1$ ,即 $a = 2$ 
                
                \item 要证: $f\left( x\right)  > {x}^{2} - 1$ ,即证: ${{e}}^{x} - {2x} > 0$ ,
                
                设 $h\left( x\right)  = {{e}}^{x} - {2x},{h}^{\prime }\left( x\right)  = {{e}}^{x} - 2$ ,
                
                令 ${h}^{\prime }\left( x\right)  = 0$ ，解得 $x = \ln 2$ ，
                
                当 $x \in  \left( {-\infty ,\ln 2}\right)$ 时， $h\left( x\right)$ 单调递减， $x \in  \left( {\ln 2, + \infty }\right)$ 时， $h\left( x\right)$ 单调递增，
                
                所以 $h\left( x\right)  \geq  h\left( {\ln 2}\right)  = {{e}}^{\ln 2} - 2\ln 2 = 2 - \ln 4 = {\operatorname{lne}}^{2} - \ln 4 > 0$ ,
                
                所以 $f\left( x\right)  > {x}^{2} - 1$ .
            \end{enumerate}
        \end{jieda}
        
        \item 已知函数 $f\left( x\right)  =  - \dfrac{1}{2}{{e}}^{2x} + 4{{e}}^{x} - {ax},f\left( x\right)$ 有两个极值点 ${x}_{1},{x}_{2}$ .
        \begin{enumerate}
            \item 求 $a$ 的取值范围
            \item 证明: $a - \left( {a - 1}\right) \ln a - 2 < 0$
        \end{enumerate}
        \begin{jieda}
            \begin{enumerate}
                \item 由题得 ${f}^{\prime }\left( x\right)  =  - {{e}}^{2x} + 4{{e}}^{x} - a = 0$ 有两个根
                
                令 $t = {{e}}^{x}$ ,则 $- {t}^{2} + {4t} - a = 0$ 有两个正根,
                
                令 ${t}_{1} = {{e}}^{{x}_{1}},{t}_{2} = {{e}}^{{x}_{2}}$ ,则 ${t}_{1},{t}_{2}$ 是方程 ${t}^{2} - {4t} + a = 0$ 的两个正根,
                
                则 $\left\{  \begin{matrix} \Delta  = {\left( -4\right) }^{2} - {4a} = {16} - {4a} > 0 \\  {t}_{1}{t}_{2} = a > 0 \\  {t}_{1} + {t}_{2} = 4 \end{matrix}\right.$ ,解得 $0 < a < 4$ .

                经检验，此时$f(x)$有两个极值点.

                \item 证明:令 $g\left( x\right)  = x - \left( {x - 1}\right) \ln x - 2\left( {0 < x < 4}\right)$ ，则 ${g}^{\prime }\left( x\right)  = 1 - \left( {\ln x + \dfrac{x - 1}{x}}\right)  = \dfrac{1}{x} - \ln x$，$g''\left( x\right)  =  - \dfrac{1}{{x}^{2}} - \dfrac{1}{x} < 0$，则 ${g}^{\prime }\left( x\right)$ 在 $\left( {0,4}\right)$ 上单调递减.
                
                又 ${g}^{\prime }\left( 1\right) = 1,{g}^{\prime }\left( 2\right)  = \dfrac{1}{2} - \ln 2 < 0$，故存在 ${x}_{0} \in  \left( {1,2}\right)$ ,使 ${g}^{\prime }\left( {x}_{0}\right)  = \dfrac{1}{{x}_{0}} - \ln {x}_{0} = 0$ ,即 $\dfrac{1}{{x}_{0}} = \ln {x}_{0}$ ,则当 $x \in  \left( {0,{x}_{0}}\right)$ 时, ${g}^{\prime }\left( x\right)  > 0$ ;
                
                当 $x \in  \left( {{x}_{0},4}\right)$ 时， ${g}^{\prime }\left( x\right)  < 0$ ，故 $g\left( x\right)$ 在 $\left( {0,{x}_{0}}\right)$ 上单调递增， $g\left( x\right)$ 在 $\left( {{x}_{0},4}\right)$ 上单调递减，
                
                则 $g\left( x\right)  \leq  g\left( {x}_{0}\right)  = {x}_{0} - \left( {{x}_{0} - 1}\right) \ln {x}_{0} - 2 = {x}_{0} - \left( {{x}_{0} - 1}\right)  \times  \dfrac{1}{{x}_{0}} - 2 = {x}_{0} + \dfrac{1}{{x}_{0}} - 3$ .
                
                又 ${x}_{0} \in  \left( {1,2}\right)$ ，则 ${x}_{0} + \dfrac{1}{{x}_{0}} \in  \left( {2,\dfrac{5}{2}}\right)$ ，故 $g\left( {x}_{0}\right)  = {x}_{0} + \dfrac{1}{{x}_{0}} - 3 < 0$ ，即 $g\left( x\right)  < 0$，所以 $a - \left( {a - 1}\right) \ln a - 2 < 0$ .
            \end{enumerate}
        \end{jieda}
        
        \item 已知函数 $f\left( x\right)  = \ln x - x + a$ .
        \begin{enumerate}
            \item 若 $x > 0$ 时 $f\left( x\right)  < 0$ ，求 $a$ 的取值范围
            \item 若 $0 < a \leq  1$ ，证明:当 $x \geq  1$ 时， $f\left( x\right)  + x \leq  \left( {x - 1}\right) {{e}}^{x - a} + 1$ 
        \end{enumerate}
        \begin{jieda}
            \begin{enumerate}
                \item $\ln x - x + a < 0 \Longrightarrow  a < x - \ln x$ ,
                
                令 $g\left( x\right)  = x - \ln x,x > 0$ ,
                
                则 ${g}^{\prime }\left( x\right)  = 1 - \dfrac{1}{x} = \dfrac{x - 1}{x}$ ,令 ${g}^{\prime }\left( x\right)  > 0$ 得 $x > 1$ ,令 ${g}^{\prime }\left( x\right)  < 0$ 得 $0 < x < 1$ ,
                
                故 $g\left( x\right)  = x - \ln x$ 在 $\left( {0,1}\right)$ 上单调递减,在 $\left( {1, + \infty }\right)$ 上单调递增,
                
                故 $g\left( x\right)  = x - \ln x$ 在 $x = 1$ 处取得极小值,也是最小值,最小值为 $g\left( 1\right)  = 1$ ,
                
                故 $a < 1$

                \item $f\left( x\right)  + x \leq  \left( {x - 1}\right) {{e}}^{x - a} + 1 \Longleftrightarrow  \ln x + a \leq  \left( {x - 1}\right) {{e}}^{x - a} + 1$ ,
                
                令 $h\left( x\right)  = \ln x + a - 1 - \left( {x - 1}\right) {{e}}^{x - a},x \geq  1$，则 ${h}^{\prime }\left( x\right)  = \dfrac{1}{x} - x{{e}}^{x - a}$
                
                $h''\left( x\right)  =  - \dfrac{1}{{x}^{2}} - \left( {x + 1}\right) {{e}}^{x - a} < 0$ 恒成立，故 ${h}^{\prime }\left( x\right)  = \dfrac{1}{x} - x{{e}}^{x - a}$ 在 $\lbrack 1, + \infty )$ 上单调递减，
                
                又 $0 < a \leq  1$ ，故 ${h}^{\prime }\left( 1\right)  = 1 - {{e}}^{1 - a} \leq  0$，故 ${h}^{\prime }\left( x\right)  = \dfrac{1}{x} - x{{e}}^{x - a} \leq  0$ 在 $\lbrack 1, + \infty )$ 上恒成立，故 $h\left( x\right)  = \ln x + a - 1 - \left( {x - 1}\right) {{e}}^{x - a}$ 在 $\lbrack 1, + \infty )$ 上单调递减，
                
                又 $h\left( 1\right)  = a - 1 \leq  0$，故 $\ln x + a - 1 - \left( {x - 1}\right) {{e}}^{x - a} \leq  0$ ,结论得证.
            \end{enumerate}
        \end{jieda}
        \item 已知函数 $f\left( x\right)  = \ln x - a{x}^{2}$ .
        \begin{enumerate}
            \item 当 $a = \dfrac{1}{2}$ 时,求函数的单调区间;
            \item 若 $f\left( x\right)  \leq  0$ 恒成立,求 $a$ 的取值范围；
            \item 求证： $2\ln \left( {n + 1}\right)  < 3 + \dfrac{5}{4} + \dfrac{7}{9} + \cdots  + \dfrac{{2n} + 1}{{n}^{2}}$ .
        \end{enumerate}
        \begin{jieda}
            \begin{enumerate}
                \item 当 $a = \dfrac{1}{2}$ 时, $f\left( x\right)  = \ln x - \dfrac{1}{2}{x}^{2}$ ,定义域为 $\left( {0, + \infty }\right)$ ,${f}^{\prime }\left( x\right)  = \dfrac{1}{x} - x = \dfrac{1 - {x}^{2}}{x},$令 ${f}^{\prime }\left( x\right)  > 0$ 可得 $0 < x < 1$ ,增区间为(0,1); 令 ${f}^{\prime }\left( x\right)  < 0$ 可得 $x > 1$ ,减区间为 $\left( {1, + \infty }\right)$ 
                \item 由 $f\left( x\right)  \leq  0$ 恒成立,可得 $a \geq  \dfrac{\ln x}{{x}^{2}}$ 恒成立,
                令 $g\left( x\right)  = \dfrac{\ln x}{{x}^{2}}$ ,则 ${g}^{\prime }\left( x\right)  = \dfrac{1 - 2\ln x}{{x}^{3}}$ ,令 ${g}^{\prime }\left( x\right)  > 0$ 得 $0 < x < \sqrt{{e}},g\left( x\right)$ 的增区间为 $\left( {0,\sqrt{{e}}}\right)$ ,
                令 ${g}^{\prime }\left( x\right)  < 0$ 得 $x > \sqrt{{e}}$ , $g\left( x\right)$ 的减区间为 $\left( {\sqrt{{e}}, + \infty }\right)$ ,所以 $g\left( x\right)$ 的最大值为 $g\left( \sqrt{{e}}\right)  = \dfrac{1}{2{e}}$ ,所以 $a \geq  \dfrac{1}{2{e}}$ ,故 $a$
                的取值范围是 $\left\lbrack  {\dfrac{1}{2{e}}, + \infty }\right)$ .
                \item 证明：设 $h\left( x\right)  = 2\ln \left( {1 + x}\right)  - {2x} - {x}^{2},x > 0$ ,
                ${h}^{\prime }\left( x\right)  = \dfrac{2}{x + 1} - 2 - {2x} =  - \dfrac{2{x}^{2} + {4x}}{x + 1}$ ,当 $x > 0$ 时, ${h}^{\prime }\left( x\right)  < 0,h\left( x\right)$ 为减函数,
                所以 $h\left( x\right)  < h\left( 0\right)  = 0$ ,即 $2\ln \left( {1 + x}\right)  < {2x} + {x}^{2}$ ,
                令 $x = \dfrac{1}{n}$ ,则 $2\ln \left( {1 + \dfrac{1}{n}}\right)  < \dfrac{2}{n} + \dfrac{1}{{n}^{2}} = \dfrac{{2n} + 1}{{n}^{2}}$ ,
                所以 $2\ln \left( {1 + \dfrac{1}{1}}\right)  < \dfrac{2 \times  1 + 1}{{1}^{2}},2\ln \left( {1 + \dfrac{1}{2}}\right)  < \dfrac{2 \times  2 + 1}{{2}^{2}},2\ln \left( {1 + \dfrac{1}{3}}\right)  < \dfrac{2 \times  3 + 1}{{3}^{2}}$，$2\ln \left( {1 + \dfrac{1}{n}}\right)  < \dfrac{2 \times  n + 1}{{n}^{2}}$
                以上各式相加可得 $2\ln \left( {n + 1}\right)  < 3 + \dfrac{5}{4} + \dfrac{7}{9} + \cdots  + \dfrac{{2n} + 1}{{n}^{2}}$.

                \ps 根据切线不等式$ln(1+x) \leq x$可以证得$2\ln \left( {n + 1}\right) <2 + \dfrac{4}{4} + \dfrac{6}{9} + \cdots  + \dfrac{{2n}}{{n}^{2}} < 3 + \dfrac{5}{4} + \dfrac{7}{9} + \cdots  + \dfrac{{2n} + 1}{{n}^{2}}$.
            \end{enumerate}
        \end{jieda}
        \item 已知函数 $f\left( x\right)  = x{e}^{ax} - {e}^{x},a$ 是实数.
        \begin{enumerate}
            \item 当 $a = 1$ 时,求 $f\left( x\right)$ 的极值
            \item 求最大的实数 $a$ ,使得 $x > 0$ 时恒有 $f\left( x\right)  <  - 1$
            \item 对于正整数 $n$ ,证明: $\dfrac{1}{\sqrt{{1}^{2} + 1}} + \dfrac{1}{\sqrt{{2}^{2} + 2}} + \cdots  + \dfrac{1}{\sqrt{{n}^{2} + n}} > \ln \left( {n + 1}\right)$ .
        \end{enumerate}
        \begin{jieda}
            \begin{enumerate}
                \item $f\left( x\right) = (x-1)e^x$，$f'(x) = xe^x$，故$x = 0$是$f(x)$的唯一驻点，$x < 0$，$f'(x) < 0$，$f(x)$严减；$x > 0$，$f'(x) > 0$，$f(x)$严增；故$f(0) = -1$是$f(x)$极大值，无极小值.
                \item 设$g(x) = f(x)+1$，$f\left( x\right)  <  - 1$即$g(x) < 0$，$g(0) = 0$

                $g'(x) = (1+ax)e^{ax} - e^x$，$g'(0) = 0$

                $g''(x) = (a^2x+2a)e^{ax} - e^x$，$g''(0) = 2a-1$

                若$g''(0) > 0$，即$a > \dfrac{1}{2}$，则$\exists x_0 > 0$使$g'(x)$在$(0, x_0)$严增，故$g'(x) > 0$在$(0, x_0)$恒成立，使$g(x)$在$(0, x_0)$严增，故$g(x) > 0$在$(0, x_0)$恒成立，矛盾

                当$a \leq \dfrac{1}{2}$时，$g(x) \leq h(x) = x{e}^{\frac{x}{2}} - {e}^{x} + 1$，故$h'(x) = \left(1+\dfrac{x}{2} - \sqrt{e^x}\right)\sqrt{e^x}$，设$m(x) = 1+\dfrac{x}{2} - \sqrt{e^x}$，$m'(x) = \dfrac{1}{2} - \dfrac{1}{2}\sqrt{e^x} < 0$，因此$h'(x) < 0$，故$h(x)$在$(0, +\infty)$严减，故$h(x) < h(0) = 0$

                故$a$最大可取$\dfrac{1}{2}$

                \item $\dfrac{1}{\sqrt{{1}^{2} + 1}} + \dfrac{1}{\sqrt{{2}^{2} + 2}} + \cdots  + \dfrac{1}{\sqrt{{n}^{2} + n}} > \ln \left( {n + 1}\right) \Longleftrightarrow \dfrac{1}{\sqrt{{1}^{2} + 1}} + \dfrac{1}{\sqrt{{2}^{2} + 2}} + \cdots  + \dfrac{1}{\sqrt{{n}^{2} + n}} > (\ln \left( {n + 1}\right) - lnn) + (lnn - ln(n-1)) + ... + (ln2 - ln1)$

                故只需证$\dfrac{1}{\sqrt{{n}^{2} + n}} \geq \ln \left( {n + 1}\right) - lnn = ln\left(1+\dfrac{1}{n}\right)$对$n \in \mathbf{N^*}$恒成立，且等号不恒取等即可.

                根据上一问可知$x{e}^{\frac{x}{2}} - {e}^{x} + 1 < 0$对$x > 0$恒成立，换元$t = {e}^{x}$，即$lnt \cdot \sqrt{t} - t + 1 < 0$对$t > 1$恒成立.

                变形有$lnt < \sqrt{t} - \dfrac{1}{\sqrt{t}}$，令$t = 1+\dfrac{1}{n}$，即有$\dfrac{1}{\sqrt{{n}^{2} + n}} \geq  ln\left(1+\dfrac{1}{n}\right)$，得证.
            \end{enumerate}
        \end{jieda}
        \item 已知函数 $f\left( x\right)  = {e}^{x} - {e}^{-x} - {2x}$ .
        \begin{enumerate}
            \item 讨论 $f\left( x\right)$ 的单调性
            \item 设 $g\left( x\right)  = f\left( {2x}\right)  - {4bf}\left( x\right)$ ,当 $x > 0$ 时, $g\left( x\right)  > 0$ ,求 $b$ 的最大值
            \item 已知 ${1.4142} < \sqrt{2} < {1.4143}$ ,估计 $\ln 2$ 的近似值(精确到 0.001)
        \end{enumerate}
        \begin{jieda}
            \begin{enumerate}
                \item 因为 ${f}^{\prime }\left( x\right)  = {e}^{x} + \dfrac{1}{{e}^{x}} - 2 \geq  0$ ,当且仅当 $x = 0$ 时等号成立,所以函数 $f\left( x\right)$ 在 $\mathbf{R}$ 上是增函数
                \item 易知 $g(0) = 0$
                
                $g'\left( x\right)  = 2f'\left( {2x}\right)  - {4bf}'\left( x\right) = 2\left\lbrack  {{e}^{2x} + {e}^{-{2x}} - {2b}\left( {{e}^{x} + {e}^{-x}}\right)  + \left( {{4b} - 2}\right) }\right\rbrack$

                令$t = {{e}^{x} + {e}^{-x}}$，$t \geq 2$，因此$g'(x) = h(t) = 2[(t^2-2)-2bt+(4b-2)] = 2[t^2-2bt+(4b-4)]$

                令$h(t) = 0$，$t = 2$或$2b-2$，因此$g'(0) = h(2) = 0$
                \begin{dingautolist}{172}
                    \item 若对称轴$t = b$在$t = 2$及其左侧，即$b \leq 2$，此时$g'(x) \geq 0$恒成立，$g(x) > g(0) = 0$满足要求
                    \item 若$b > 2$，则在$t \in (2, 2b-2)$，即$x \in \left(0, ln\left(b-1+\sqrt{b^2-2b}\right)\right)$内$g'(x) < 0$，故在该区间上$g(x) < g(0) = 0$，不满足要求.
                \end{dingautolist}
                

                \item 考虑$g\left( {\ln \sqrt{2}}\right)  = \dfrac{3}{2} - 2\sqrt{2}b + 2\left( {{2b} - 1}\right) \ln 2$

                当$b \leq 2$时，$g(x)$在$(0, +\infty)$严增，故$g\left( {\ln \sqrt{2}}\right) > 0$，代入$b = 2$，可得$g\left( {\ln \sqrt{2}}\right)  = \dfrac{3}{2} - 4\sqrt{2} + 6\ln 2 > 0$，故$\ln 2 > \dfrac{8\sqrt{2} - 3}{12} > {0.6928}$

                当$b > 2$时，$g(x)$在$(0, +\infty)$先减后增，考虑令$g\left( {\ln \sqrt{2}}\right) < 0$

                $g(x)$在$\left(0, \ln \left( {b - 1 + \sqrt{{b}^{2} - {2b}}}\right)\right)$严减，故当$\ln \left( {b - 1 + \sqrt{{b}^{2} - {2b}}}\right) \geq ln\sqrt{2}$时必有$g\left( {\ln \sqrt{2}}\right) < 0$，即${b - 1 + \sqrt{{b}^{2} - {2b}}} \geq \sqrt{2}$，解得$b \geq \dfrac{3\sqrt{2}}{4} + 1$

                代入$b = \dfrac{3\sqrt{2}}{4} + 1$, $g\left( {\ln \sqrt{2}}\right)  =  - \dfrac{3}{2} - 2\sqrt{2} + \left( {3\sqrt{2} + 2}\right) \ln 2 < 0$ ,$\ln 2 < \dfrac{{18} + \sqrt{2}}{28} < {0.6934}$

                所以 $\ln 2$ 的近似值为 $0.693$.
            \end{enumerate}
        \end{jieda}
        \item 已知集合 $T = \{ \left( {x,y}\right)  \mid  x > 1,y \in  \mathbf{R}\} ,{M}_{a} = \left\{  \left( {x,y}\right)  \left|y = {\left( a\ln x + \dfrac{1}{a}\right) }^{2} \right. \right\}$. 如果存在$\left( {{x}_{0},{y}_{0}}\right)  \in  {M}_{{a}_{0}}$ ,对属于 $T$ 且不属于任意 ${M}_{a}\left( {a \neq  0}\right)$ 的所有元素$(x,y)$,都有 $y - x \leq  {y}_{0} - {x}_{0}$ 成立,则 ${y}_{0} - {x}_{0}$ 的取值范围是\blkda[3]{$\left[8ln2 - 4, +\infty\right)$}.
        \begin{jieda}
            当固定$x$的时候$y = \left( a\ln x + \dfrac{1}{a}\right)^2 \geq (2\sqrt{lnx})^2 = 4lnx$，故所有${M}_{a}\left( {a \neq  0}\right)$的并集为$y = 4lnx$及其上方部分.

            故属于 $T$ 且不属于任意 ${M}_{a}\left( {a \neq  0}\right)$ 的所有元素$(x,y)$构成集合$\{(x, y)| y < 4lnx, x > 1\}$

            设$z = y - x$，即$y = x+z$，为斜率为1的动直线，其纵截距为$z$. 数形结合，寻找边界$y = 4lnx$的斜率为1的切线为$y = x + 4ln4 - 4$，故$z < 4ln4 - 4 = 8ln2 - 4$

            因此${y}_{0} - {x}_{0} \in \left[8ln2 - 4, +\infty\right)$
        \end{jieda}
        \item 已知关于 $x$ 的不等式 $\left( {\ln x - {kx}}\right) \left\lbrack  {{x}^{2} - \left( {k + 3}\right) x + 4}\right\rbrack   \leq  0$ 对任意 $x \in  \left( {0, + \infty }\right)$ 均成立，则实数 $k$ 的取值范围为\blkda[2]{$\left\lbrack  {\dfrac{1}{{e}},1}\right\rbrack$}.
        \begin{jieda}
            $\left( {\ln x - {kx}}\right) \left\lbrack  {{x}^{2} - \left( {k + 3}\right) x + 4}\right\rbrack   \leq  0$ 对任意 $x \in  (0$ , $+ \infty )$ 均成立
            
            即$\left( {\dfrac{\ln x}{x} - k}\right) \left\lbrack  {x + \dfrac{4}{x} - \left( {k + 3}\right) }\right\rbrack   \leq  0$在$(0, +\infty)$恒成立
            
            即$ \left( {k - \dfrac{\ln x}{x}}\right)  \left\lbrack  {k - \left( {x + \dfrac{4}{x} - 3}\right) }\right\rbrack   \leq  0$在$(0, +\infty)$恒成立
            
            即$\dfrac{\ln x}{x} \leq  k \leq  x + \dfrac{4}{x} - 3$在$(0, +\infty)$恒成立
            
            ${y}_{1} = \dfrac{\ln x}{x} \leq  \dfrac{1}{{e}},{y}_{2} = x + \dfrac{4}{x} - 3 \geq  2\sqrt{x \times  \dfrac{4}{x}} - 3 = 1$ , 所以实数 $k$ 的取值范围为 $k \in  \left\lbrack  {\dfrac{1}{{e}},1}\right\rbrack$ .

            \begin{center}
            \includegraphics[width=0.2\textwidth]{bodyxb/src/12_423_1184_221_274_0.jpg}
            \end{center}
        \end{jieda}
        \item 设 $k \in  \mathbf{R}$ ,函数 $y = f\left( x\right)$ 和 $y = g\left( x\right)$ 的表达式分别为$f\left( x\right)  = e^{x} - {kx},g\left( x\right)  = {x}^{2} + \left( {k - 2}\right) x + \dfrac{1}{4}$ ,若对任意的实数 ${x}_{0}$ ,皆有 $f\left( {x}_{0}\right)  \cdot  g\left( {x}_{0}\right)  \geq  0$ 成立,则 $k$ 的取值范围是\blkda{$[1,e]$}.
        \begin{jieda}
            $f\left( {x}_{0}\right)  \cdot  g\left( {x}_{0}\right)  \geq  0$即$\left[e^{x_0} - {kx_0}\right] \cdot \left[x_0^{2} + \left( {k - 2}\right) x_0 + \dfrac{1}{4}\right] \geq 0$

            \begin{dingautolist}{172}
                \item 易知$x_0 = 0$时不等式成立
                \item 当$x_0 \neq 0$时，即$\left[\dfrac{e^{x_0}}{x_0} - k\right] \cdot \left[x_0 + \left( {k - 2}\right) + \dfrac{1}{4x_0}\right] \geq 0$
                
                即$\left[k - \dfrac{e^{x_0}}{x_0}\right] \cdot \left[k - \left(2 - \left(x_0 + \dfrac{1}{4x_0}\right)\right)\right] \leq 0$

                \begin{center}
                \includegraphics[width=0.2\textwidth]{bodyxb/src/2026putuoyimo12.png}
                \end{center}

                当$x_0 < 0$，$\dfrac{e^x_0}{x_0} \in (-\infty, 0)$，$2 - \left(x_0 + \dfrac{1}{4x_0}\right) \in [3, +\infty)$

                当$x_0 > 0$，$\dfrac{e^x_0}{x_0} \in [e, +\infty)$，$2 - \left(x_0 + \dfrac{1}{4x_0}\right) \in (-\infty, 1]$

                故$k \in [1, e]$
            \end{dingautolist}
        \end{jieda}
    \end{enumerate}

    
    \item 洛必达法则
    \begin{enumerate}
        \item 若对任意$x \geq 1$，不等式$x\left(lnx-ax\right)+a \leq 0$恒成立，求实数$a$的取值范围.
        \begin{jieda}
            不等式化为$xlnx-ax^2+a\leq 0$
            \begin{enumerate}
                \item 方法1 分类讨论
                
                设$f(x) = xlnx-ax^2+a$，则$f'(x) = lnx+1-2ax$，$f''(x) = \dfrac{1}{x}-2a$

                \begin{dingautolist}{172}
                    \item 若$a \leq 0$，则$f''(x) > 0$恒成立，故$f'(x)$严增，故$f'(x) \geq f'(1) = 1-2a > 0$，故$f(x)$严增，故$f(x) \geq f(1)$，不符合要求
                    \item 若$a \geq \dfrac{1}{2}$，则$f''(x) \leq 0$恒成立，故$f'(x)$严减，故$f'(x) \leq f'(1) = 1-2a \leq 0$，故$f(x)$严减，故$f(x) \leq f(1)$，满足要求
                    \item 若$a \in \left(0, \dfrac{1}{2}\right)$，则$f''(x)$在$\left(1, \dfrac{1}{2a}\right)$大于0，故$f'(x)$在$\left(1, \dfrac{1}{2a}\right)$严增，故$f'(x)$在$\left(1, \dfrac{1}{2a}\right)$大于0，故$f(x)$在$\left(1, \dfrac{1}{2a}\right)$严增，故$f(x)$在$\left(1, \dfrac{1}{2a}\right)$大于0，不符合要求.
                \end{dingautolist}
                综上，$a \in \left[\dfrac{1}{2}, +\infty\right)$
            
                \item 方法2 参变分离

                当$x = 1$时恒成立，参变分离有$a \geq \dfrac{xlnx}{x^2-1}(x > 1)$，设$f(x) = \dfrac{xlnx}{x^2-1}$，则$f'(x) = \dfrac{(lnx+1)(x^2-1)-2x^2lnx}{(x^2-1)^2} = \dfrac{x^2-1-(1+x^2)lnx}{(x^2-1)^2}$

                令$g(x) = x^2-1-(1+x^2)lnx$，$g'(x) = 2x-2xlnx-\dfrac{1+x^2}{x} = x-2xlnx-\dfrac{1}{x}$，$g''(x) = 1+\dfrac{1}{x^2}-2(lnx+1)$，$g'''(x) = -\dfrac{2}{x^3}-\dfrac{2}{x} < 0$，故$g''(x) < 0$，故$g'(x) < 0$，故$g(x) < 0$，故$f'(x) < 0$，根据洛必达法则可知$f(x) < \dfrac{1}{2}$，故$a \geq \dfrac{1}{2}$
            \end{enumerate}
        \end{jieda}
        \item 已知函数$f(x) = \left(\dfrac{1}{x}+a\right)ln(x+1)$，若函数$f(x)$在$(0, +\infty)$上存在极值点，求实数$a$的取值范围.
        \begin{jieda}
            $f'(x) = \left(-\dfrac{1}{x^2}\right)ln(x+1)+\left(\dfrac{1}{x}+a\right)\dfrac{1}{x+1}$，故$f'(x)$在$(0, +\infty)$上存在变号零点

            $f'(x) = 0 \Longleftrightarrow (x+1)ln(x+1) = ax^2+x$
            \begin{enumerate}
                \item 方法1 分类讨论

                设$g(x) = (x+1)ln(x+1) - ax^2-x$，故$g'(x) = ln(x+1) - 2ax$，$g''(x) = \dfrac{1}{x+1}-2a$

                \begin{dingautolist}{172}
                    \item 若$a \leq 0$，则$g''(x) > 0$在$(0, +\infty)$恒成立，即$g'(x)$在$(0, +\infty)$严增，又因为$g'(0) = 0$，所以$g'(x) > 0$在$(0, +\infty)$恒成立，即$g(x)$在$(0, +\infty)$严增，又因为$g(0) = 0$，所以$g(x) > 0$在$(0, +\infty)$恒成立，不符合要求

                    \item 若$a \geq \dfrac{1}{2}$，则$g''(x) < 0$在$(0, +\infty)$恒成立，即$g'(x)$在$(0, +\infty)$严减，又因为$g'(0) = 0$，所以$g'(x) < 0$在$(0, +\infty)$恒成立，即$g(x)$在$(0, +\infty)$严减，又因为$g(0) = 0$，所以$g(x) < 0$在$(0, +\infty)$恒成立，不符合要求

                    \item 若$a \in \left(0, \dfrac{1}{2}\right)$，则$g''(x)$在$\left(0, \dfrac{1}{2a}-1\right)$大于$0$，在$\left(\dfrac{1}{2a}-1, +\infty\right)$小于$0$.
                    故$g'(x)$在$\left(0, \dfrac{1}{2a}-1\right)$严增，在$\left(\dfrac{1}{2a}-1, +\infty\right)$严减，因为$g'(0) = 0$，由一次函数和对数函数性质可知$x \rightarrow +\infty$时$g'(x) < 0$，故$g'(x)$在$(0, +\infty)$存在唯一变号零点$x_0$

                    又因为$g(0) = 0$，所以$g(x)$在$(0, x_0)$严增，在$(x_0, +\infty)$严减，由二次函数和对数函数性质可知$x \rightarrow +\infty$时$g(x) < 0$，故$g(x)$在$(0, +\infty)$存在唯一变号零点$x_1$


                    故$f'(x)$在$(0, +\infty)$上存在变号零点
                \end{dingautolist}
                因此$a \in \left(0, \dfrac{1}{2}\right)$满足要求

                \item 方法2 参变分离

                将方程变形为$a = \dfrac{(x+1)ln(x+1)-x}{x^2}$，设$g(x) = \dfrac{(x+1)ln(x+1)-x}{x^2}$，则$g'(x) = \dfrac{x^2ln(x+1)-2x[(x+1)ln(x+1)-x]}{x^4} = \dfrac{(-x^2-2x)ln(x+1)+2x^2}{x^4} = \dfrac{(-x-2)ln(x+1)+2x}{x^3}$

                设$h(x) = (-x-2)ln(x+1)+2x$，则$h'(x) = -ln(x+1)-\dfrac{x+2}{x+1}+2 = \dfrac{x}{x+1}-ln(x+1)$，$h''(x) = \dfrac{1}{(x+1)^2}-\dfrac{1}{x+1} < 0$，故$h'(x) < h'(0) = 0$，故$h(x) < h(0) = 0$，故$g'(x) < 0$，故$g(x)$在$(0, + \infty)$严减

                根据洛必达法则可得$x \rightarrow 0$时$g(x) \rightarrow \dfrac{1}{2}$，因此$a \in \left(0, \dfrac{1}{2}\right)$
            \end{enumerate}
        \end{jieda}
    \end{enumerate}

    \item 双变量恒成立问题
    \begin{enumerate}
        \item 已知函数 $f\left( x\right)  = \dfrac{2 - x}{{e}^{x}}$ ,若对任意 ${x}_{1},{x}_{2} \in  \lbrack a, + \infty )$ ,都有 $f\left( {x}_{1}\right)  - f\left( {x}_{2}\right)  \geq   - \dfrac{1}{{e}^{3}}$ 成立,则实数 $a$ 的最小值为\blkda{2}.
        \begin{jieda}
            $f\left( {x}_{1}\right)  - f\left( {x}_{2}\right)  \geq   - \dfrac{1}{{e}^{3}}$即$f_{min} - f_{max} \geq -\dfrac{1}{e^3}$
        
            $f'(x) = \dfrac{x-3}{e^x}$，故$f(x)$在$(-\infty, 3)$严减，在$(3, +\infty)$严增，当$x \rightarrow +\infty$时$f(x) \rightarrow 0$，$f(3) = -\dfrac{1}{e^3}$

            结合图象可知$a = 2$.
        \end{jieda}
        \item 已知函数 $f\left( x\right)  = {2k}{x}^{2} - 4\ln x,g\left( x\right)  = \ln \dfrac{x}{k}$ ,其中 $x \in  (0,{e}\rbrack$ , $k > 0$ .
        \begin{enumerate}
            \item 若 $y = f\left( x\right)  + \dfrac{32}{9}x$ 在 $x = 1$ 处取得极值,求 $k$ 的值
            \item 讨论函数 $f\left( x\right)$ 的单调性
            \item 若对任意 ${x}_{1},{x}_{2} \in  (0,{e}\rbrack$ ，当 $k > 1$ 时，不等式 $f\left( {x}_{1}\right)  > g\left( {x}_{2}\right)  + 4$ 恒成立，求 $k$ 的取值范围
        \end{enumerate}
        \begin{jieda}
            \begin{enumerate}
                \item 令 $t\left( x\right)  = f\left( x\right)  + \dfrac{32}{9}x$ ,由题意 ${f}^{\prime }\left( x\right)  = {4kx} - \dfrac{4}{x} = \dfrac{{4k}{x}^{2} - 4}{x},x \in  \left( {0,{e}}\right\rbrack  .{t}^{\prime }\left( x\right)  = \dfrac{{4k}{x}^{2} - 4}{x} + \dfrac{32}{9}$
                
                由已知得 ${t}^{\prime }\left( 1\right)  = \dfrac{{4k} - 4}{1} + \dfrac{32}{9} = 0$ ，解得 $k = \dfrac{1}{9}$ ，
                
                此时 ${t}^{\prime }\left( x\right)  = \dfrac{\dfrac{4}{9}{x}^{2} - 4}{x} + \dfrac{32}{9} = \dfrac{4\left( {x - 1}\right) \left( {x + 9}\right) }{9x}$ ,
                
                易知在区间 $\left( {0,1}\right)$ 上 $f\left( x\right)$ 单调递增,在 $\left( {1,{e}}\right)$ 上 $f\left( x\right)$ 单调递减,
                
                则函数 $f\left( x\right)$ 在 $x = 3$ 处取得极小值,因此 $k = \dfrac{1}{9}$
                
                
                \item 由题意 $f\left( x\right)  = \dfrac{{4k}\left( {{x}^{2} - \dfrac{1}{k}}\right) }{x} = \dfrac{{4k}\left( {x + \dfrac{\sqrt{k}}{k}}\right) \left( {x - \dfrac{\sqrt{k}}{k}}\right) }{x}$ ，其中 $x \in  (0,{e}\rbrack ,k > 0$

                \begin{dingautolist}{172}
                    \item 当 $\dfrac{\sqrt{k}}{k} < {e}$ ，即 $k > \dfrac{1}{{{e}}^{2}},f\left( x\right)$ 在 $\left( {0,\dfrac{\sqrt{k}}{k}}\right)$ 上单调递减，在 $\left( {\dfrac{\sqrt{k}}{k},{e}}\right\rbrack$ 上单调递增.
                    \item 当 $\dfrac{\sqrt{k}}{k} \geq  {e}$ ,即 $0 < k \leq  \dfrac{1}{{{e}}^{2}}$ ,则 $f\left( x\right)$ 在 $(0,{e}\rbrack$ 上单调递减
                \end{dingautolist}
                
                综上,当 $0 < k \leq  \dfrac{1}{{{e}}^{2}}$ 时, $f\left( x\right)$ 的单调递减区间为 $(0,{e}\rbrack$
                
                当 $k > \dfrac{1}{{{e}}^{2}}$ 时, $f\left( x\right)$ 的单调递减区间为 $\left( {0,\dfrac{\sqrt{k}}{k}}\right)$ ,单调递增区间为 $\left( {\dfrac{\sqrt{k}}{k},{e}}\right)$ .
                
                \item 当 $k > 1 > \dfrac{1}{{{e}}^{2}}$ 时，由(2)可知当 $x = \dfrac{\sqrt{k}}{k}$ 时，函数 $f\left( x\right)$ 取得最小值，即 $f{\left( x\right) }_{\min } = f\left( \dfrac{\sqrt{k}}{k}\right)  = 2 + 2\ln k$ ,
                
                由 $g\left( x\right)  = \ln \dfrac{x}{k} = \ln x - \ln k$ ,可得 $g\left( x\right)$ 在 $(0,{e}\rbrack$ 上单调递增,
                
                即当 $x = {e}$ 时， $g{\left( x\right) }_{\max } = g\left( {e}\right)  = 1 - \ln k$ ，
                
                对任意 ${x}_{1},{x}_{2} \in  (0,{e}\rbrack$ ,当 $k > 1$ 时,不等式 $f\left( {x}_{1}\right)  > g\left( {x}_{2}\right)  + 4$ 恒成立,
                
                则必有 $f{\left( x\right) }_{\min } > g{\left( x\right) }_{\max } + 4$ ，即 $2 + 2\ln k > 1 - \ln k + 4$ ，解得 $k > {e}$ ，
                
                所以 $k$ 的取值范围是 $\left( {{e}, + \infty }\right)$ .
            \end{enumerate}
        \end{jieda}
        
        \item 已知函数$f(x) = \dfrac{1}{2}ax^2-(2a+1)x + 2lnx$（$a \in \mathbf{R}$）
        \begin{enumerate}
            \item 求$y = f(x)$单调区间
            \item 设$g(x) = x^2-2x$，若对任意的$x_1 \in (0, 2]$，均存在$x_2 \in (0, 2]$，使得$f(x_1) < g(x_2)$，求$a$的取值范围.
        \end{enumerate}
        \begin{jieda}
            \begin{enumerate}
                \item $f'(x) = ax-(2a+1) + \dfrac{2}{x} = \dfrac{(ax-1)(x-2)}{x} (x > 0)$
                \begin{dingautolist}{172}
                    \item 若$a < 0$，则$f(x)$在$(0, 2)$内严增，在$(2, +\infty)$严减
                    \item 若$a = 0$，则$f(x)$在$(0, 2)$内严增，在$(2, +\infty)$严减
                    \item 若$a \in \left(0, \dfrac{1}{2}\right)$，则$f(x)$在$(0, 2)$严增，在$\left(2, \dfrac{1}{a}\right)$严减，在$\left(\dfrac{1}{a}, +\infty\right)$严增.
                    \item 若$a = \dfrac{1}{2}$，则$f(x)$在$(0, +\infty)$严增
                    \item 若$a \in \left(\dfrac{1}{2}, +\infty\right)$，则$f(x)$在$\left(0, \dfrac{1}{a}\right)$严增，在$\left(\dfrac{1}{a}, 2\right)$严减，在$\left(2, +\infty\right)$严增.
                \end{dingautolist}
                \item 由题意可知$g(x)$在$(0, 2]$上的最大值大于$f(x)$在$(0, 2]$上的最大值. $g(x)$在$(0, 2]$上的最大值为$0$. 
                \begin{dingautolist}{172}
                    \item 若$a \leq \dfrac{1}{2}$，则$f(x) \leq f(2) = -2-2a+2ln2$，当$a \in \left(ln2-1, \dfrac{1}{2}\right]$时满足要求
                    \item 若$a \in \left(\dfrac{1}{2}, +\infty\right)$，则$f(x) \leq f\left(\dfrac{1}{a}\right) = -\dfrac{1}{2a} - 2-2lna$ \\
                    设$h(a) = -\dfrac{1}{2a} - 2-2lna$，$h'(a) = \dfrac{1}{2a^2} - \dfrac{2}{a}$，$h'(a) < 0$在$a \in \left(\dfrac{1}{2}, +\infty\right)$恒成立，故$h(a) < h\left(\dfrac{1}{2}\right) = -3+2ln2 < 0$，因此满足要求
                \end{dingautolist}
                综上，$a \in (ln2-1, +\infty)$
            \end{enumerate}
        \end{jieda}
    \end{enumerate}
    
    \item 利用导数研究不等式存在性问题
    \begin{enumerate}
        \item 已知函数 $f\left( x\right)  = x\ln x$ ,若存在 $x \in  \left( {0, + \infty }\right)$ ,使得 $f\left( x\right)  \leq  \dfrac{-{x}^{2} + {mx} - 3}{2}$ 成立，则实数 $m$ 的最小值是\blkda{4}
        \begin{jieda}
            由 $f\left( x\right)  \leq  \dfrac{-{x}^{2} + {mx} - 3}{2} \Longrightarrow  m \geq  \dfrac{{2x}\ln x + {x}^{2} + 3}{x}\left( {x > 0}\right)$ 能成立，问题转化为 $m \geq  {\left( 2\ln x + x + \dfrac{3}{x}\right) }_{\min }$
            
            令 $g\left( x\right)  = 2\ln x + x + \dfrac{3}{x} \Longrightarrow  {g}^{\prime }\left( x\right)  = \dfrac{2}{x} + 1 - \dfrac{3}{{x}^{2}} = \dfrac{\left( {x + 3}\right) \left( {x - 1}\right) }{{x}^{2}}$ ,
            
            由 $x > 1 \Longrightarrow  {g}^{\prime }\left( x\right)  > 0$ ; 由 $0 < x < 1 \Longrightarrow  {g}^{\prime }\left( x\right)  < 0$ ,
            
            $\therefore g\left( x\right)$ 在 $\left( {0,1}\right)$ 上单调递减,在 $\left( {1, + \infty }\right)$ 上单调递增,
            
            $\therefore g{\left( x\right) }_{\min } = g\left( 1\right)  = 4$ ,则 $m \geq  4$ 
        \end{jieda}
        
        \item 已知 $a >  - 1$ ，存在唯一的整数 ${x}_{0}$ ，使得 $\left( {{x}_{0} - 1}\right) {{e}}^{2{x}_{0}} + a{x}_{0} - {2a} < 0$ 成立，则 $a$ 的取值范围是\blkda[3]{$\left( {-\dfrac{1}{2}, - \dfrac{2}{3{{e}}^{2}}}\right\rbrack$}

        \begin{jieda}
            设 $f\left( x\right)  = \left( {x - 1}\right) {{e}}^{2x},y =  - a\left( {x - 2}\right)$
            
            由题意可知函数 $f\left( x\right)$ 在直线 $y =  - a\left( {x - 2}\right)$ 下方的图象有且只有一个点的横坐标为整数,
            
            因为 $f\left( x\right)  = \left( {x - 1}\right) {{e}}^{2x}$ ,所以 ${f}^{\prime }\left( x\right)  = \left( {{2x} - 1}\right) {{e}}^{2x}$，由 ${f}^{\prime }\left( x\right)  > 0$ ，解得 $x > \dfrac{1}{2}$ ，由 ${f}^{\prime }\left( x\right)  < 0$ ，解得 $x < \dfrac{1}{2}$ ，则 $f\left( x\right)$ 在 $\left( {\dfrac{1}{2}, + \infty }\right)$ 上单调递增，在 $\left( {-\infty ,\dfrac{1}{2}}\right)$ 上单调递减；
            
            又 $f\left( 0\right)  =  - 1,f\left( {-1}\right)  =  - 2{{e}}^{-2},f\left( 1\right)  = 0$ ,即 $f\left( x\right)$ 过点 $B\left( {0, - 1}\right) ,C\left( {-1, - 2{{e}}^{-2}}\right)$，且当 $x < 1$ 时 $f\left( x\right)  < 0$ ,当 $x > 1$ 时 $f\left( x\right)  > 0$ 
            
            如图，作出 $f\left( x\right)$ 的大致图象如下所示:
            
            \begin{center}
            \includegraphics[width=0.3\textwidth]{bodyxb/src/bo_d5rb33bef24c73bob2k0_17_141_939_397_232_0.jpg}
            \end{center}
            
            
            因为直线 $y =  - a\left( {x - 2}\right)$ 过定点 $A\left( {2,0}\right)$ ,且当 $x = 1$ 时 $y = a$ ,
            
            所以 ${k}_{AC} \leq   - a < {k}_{AB}$ ,即 $\dfrac{2}{3{{e}}^{2}} \leq   - a < \dfrac{1}{2}$ ,故 $- \dfrac{1}{2} < a \leq   - \dfrac{2}{3{{e}}^{2}}$ ,
            
            即 $a$ 的取值范围是 $\left( {-\dfrac{1}{2}, - \dfrac{2}{3{{e}}^{2}}}\right\rbrack$ .
        \end{jieda}
        \item 已知函数 $f\left( x\right)  = {ax} - \ln x,a \in  \mathbf{R}$
        \begin{enumerate}
            \item 讨论 $f\left( x\right)$ 在区间 $\left\lbrack  {1,2}\right\rbrack$ 上的单调性
            \item 若 $x \in  \left(0,e\right]$ ，不等式 $f\left( x\right)  \leq  3$ 有解，求 $a$ 的取值范围
        \end{enumerate}
        \begin{jieda}
            \begin{enumerate}
                \item 函数 $f\left( x\right)  = {ax} - \ln x$ ，求导得 ${f}^{\prime }\left( x\right)  = a - \dfrac{1}{x}$ ，由 $x \in  \left\lbrack  {1,2}\right\rbrack$ ，得 $\dfrac{1}{x} \in  \left\lbrack  {\dfrac{1}{2},1}\right\rbrack$ ，
                
                当 $a \leq  \dfrac{1}{2}$ 时， ${f}^{\prime }\left( x\right)  \leq  0$ ，当且仅当 $a = \dfrac{1}{2}$ ， $x = 2$ 时取等号，函数 $f\left( x\right)$ 在 $\left\lbrack  {1,2}\right\rbrack$ 上单调递减；
                
                当 $a \geq  1$ 时， ${f}^{\prime }\left( x\right)  \geq  0$ ，当且仅当 $a = 1,x = 1$ 时取等号，函数 $f\left( x\right)$ 在 $\left\lbrack  {1,2}\right\rbrack$ 上单调递增；
                
                当 $\dfrac{1}{2} < a < 1$ 时,由 ${f}^{\prime }\left( x\right)  < 0$ ,得 $1 \leq  x < \dfrac{1}{a}$ ; 由 ${f}^{\prime }\left( x\right)  > 0$ ,得 $\dfrac{1}{a} < x \leq  2$ ,
                
                因此函数 $f\left( x\right)$ 在 $\left\lbrack  {1,\dfrac{1}{a}}\right)$ 上单调递减，在 $\left( {\dfrac{1}{a},2}\right\rbrack$ 上单调递增，
                
                所以当 $a \leq  \dfrac{1}{2}$ 时，函数 $f\left( x\right)$ 在 $\left\lbrack  {1,2}\right\rbrack$ 上单调递减；
                
                当 $\dfrac{1}{2} < a < 1$ 时，函数 $f\left( x\right)$ 在 $\left\lbrack  {1,\dfrac{1}{a}}\right)$ 上单调递减，在 $\left( {\dfrac{1}{a},2}\right\rbrack$ 上单调递增；
                
                当 $a \geq  1$ 时，函数 $f\left( x\right)$ 在 $\left\lbrack  {1,2}\right\rbrack$ 上单调递增.

                \item 依题意，不等式 ${ax} - \ln x \leq  3$ 在 $x \in  \left(0,e\right]$时有解，即$a \leq  \dfrac{3}{x} + \dfrac{\ln x}{x}$ 在 $x \in  \left(0,e\right]$时有解.
                
                令 $g\left( x\right)  = \dfrac{3}{x} + \dfrac{\ln x}{x},x \in  \left( 0,e\right\rbrack$,求导得${g}^{\prime }\left( x\right)  =  - \dfrac{3}{{x}^{2}} + \dfrac{1 - \ln x}{{x}^{2}} =  - \dfrac{2 + \ln x}{{x}^{2}}$ ,
                
                由 ${g}^{\prime }\left( x\right)  > 0$ ,得 $0 < x < \dfrac{1}{{{e}}^{2}}$ ; 由 ${g}^{\prime }\left( x\right)  < 0,\dfrac{1}{{{e}}^{2}} < x \leq  {e}$ ,
                
                函数 $g\left( x\right)$ 在 $\left( {0,\dfrac{1}{{{e}}^{2}}}\right)$ 上单调递增,在 $\left( {\dfrac{1}{{{e}}^{2}},}\right. {e}J$ 上单调递减,
                
                当 $x = \dfrac{1}{{{e}}^{2}}$ 时，函数 $g\left( x\right)$ 取得最大值 $g\left( \dfrac{1}{{{e}}^{2}}\right)  = {{e}}^{2}$ ，因此 $a \leq  {{e}}^{2}$ ，
                
                所以实数 $a$ 的取值范围是 $\left( {-\infty ,{{e}}^{2}}\right\rbrack$ .
            \end{enumerate}
        \end{jieda}

    \end{enumerate}
    
    \item 同构-脱壳
    \begin{enumerate}
        \item 已知 $a > 0$ ，设函数 $f\left( x\right)  = {{e}}^{2x} + \left( {2 - a}\right) x - \ln x - \ln a$ ，若 $f\left( x\right)  \geq  0$ 在 $\left( {0, + \infty }\right)$ 上恒成立，则 $a$ 的取值范围是\blkda{$(0,2{e}\rbrack$}
        \begin{jieda}
            由题意可知: $f\left( x\right)  = {{e}}^{2x} + \left( {2 - a}\right) x - \ln x - \ln a \geq  0$ ,整理可得 ${{e}}^{2x} + \ln \left( {{e}}^{2x}\right)  \geq  {ax} + \ln \left( {ax}\right)$
            
            设 $g\left( x\right)  = x + \ln x,x > 0$ ,则 ${g}^{\prime }\left( x\right)  = 1 + \dfrac{1}{x} > 0$ ,可知 $g\left( x\right)$ 在 $\left( {0, + \infty }\right)$ 内单调递增,
            
            由题意可知: $g\left( {{e}}^{2x}\right)  \geq  g\left( {ax}\right)$ ,则 ${{e}}^{2x} \geq  {ax}$ 对任意 $x \in  \left( {0, + \infty }\right)$ 内恒成立，可得 $a \leq  \dfrac{{{e}}^{2x}}{x}$ 对任意 $x \in  \left( {0, + \infty }\right)$ 内恒成立
            
            设函数 $h\left( x\right)  = \dfrac{{{e}}^{2x}}{x},x > 0$ ,则 ${h}^{\prime }\left( x\right)  = \dfrac{\left( {{2x} - 1}\right) {{e}}^{2x}}{{x}^{2}}$ ,
            
            令 ${h}^{\prime }\left( x\right)  > 0$ ,解得 $x > \dfrac{1}{2}$ ,令 ${h}^{\prime }\left( x\right)  < 0$ ,解得 $0 < x < \dfrac{1}{2}$ ;
            
            可知 $h\left( x\right)$ 在 $\left( {0,\dfrac{1}{2}}\right)$ 内单调递减，在 $\left( {\dfrac{1}{2}, + \infty }\right)$ 内单调递增，所以 $h\left( x\right)$ 的最小值为 $h\left( \dfrac{1}{2}\right)  = 2{e}$ ,可得 $0 < a \leq  2{e}$，所以 $a$ 的取值范围为 $(0,2{e}\rbrack$ .
        \end{jieda}
        
        \item 对于 $x \in  \left( {0, + \infty }\right)$ ,不等式 ${{e}}^{x} - \ln \left( {mx}\right)  + \left( {1 - m}\right) x \geq  0$ 恒成立，则实数 $m$ 的取值范围为\blkda{$0 < m \leq  {e}$}

        \begin{jieda}
            已知 $x \in  \left( {0, + \infty }\right)$ ,由 ${{e}}^{x} - \ln \left( {mx}\right)  + \left( {1 - m}\right) x \geq  0$ 得, ${{e}}^{x} + x \geq  {{e}}^{\ln \left( {mx}\right) } + \ln \left( {mx}\right)$ ,
            
            构造函数 $f\left( x\right)  = {{e}}^{x} + x$ ,则 $f\left( x\right)$ 是 $\mathbf{R}$ 上的增函数,则由 $f\left( x\right)  \geq  f\left( {\ln \left( {mx}\right) }\right)$ 得: $x \geq  \ln \left( {mx}\right)$，即 $m \leq  \dfrac{{{e}}^{x}}{x}$
            
            令 $g\left( x\right)  = \dfrac{{{e}}^{x}}{x},x \in  \left( {0, + \infty }\right) ,{g}^{\prime }\left( x\right)  = \dfrac{\left( {x - 1}\right) {{e}}^{x}}{{x}^{2}}$，当 $x \in  \left( {0,1}\right)$ , ${g}^{\prime }\left( x\right)  < 0$ ,则 $g\left( x\right)$ 单调递减，当 $x \in  \left( {1, + \infty }\right)$ , ${g}^{\prime }\left( x\right)  > 0$ ,则 $g\left( x\right)$ 单调递增
            
            故$g{\left( x\right) }_{\min } = g\left( 1\right)  = {e}$ ,则 $m \leq  {e}$ ,又 $m > 0$ ,则 $0 < m \leq  {e}$ .
        \end{jieda}
        \item 已知函数 $f\left( x\right)  = \dfrac{\ln x}{x} - x + 1$ .若对于任意的 $\left( {0, + \infty }\right) ,f\left( x\right)  +x+ \dfrac{1}{x} \leq  a{e}^{x}$ 恒成立,求实数 $a$ 的最小值.
        \begin{jieda}
            $\dfrac{\ln x}{x} - x + x + 1 + \dfrac{1}{x} \leq ae^x$在$\left( {0, + \infty }\right)$ 恒成立

            即$a \geq \dfrac{lnx+x+1}{xe^x}$， 令$t = xe^x$，则$lnt = x + lnx$，故有$a \geq \dfrac{lnt+1}{t}$在$t \in (0, +\infty)$恒成立.

            \ps 若为填空题，此时直曲分离，数形结合可出答案

            设$g(t) = \dfrac{lnt+1}{t}$，$g'(t) = \dfrac{-lnt}{t^2}$，故$t \in (0, 1)$时$g'(t) > 0$，$g(t)$严增，$t \in (1, +\infty)$时$g'(t) < 0$，$g(t)$严减，因此$a \geq g(1) = 1$，所以$a$ 的最小值为$1$
        \end{jieda}
    \end{enumerate}

    
    \item 逆向构造
    \begin{enumerate}
        \item 已知偶函数 $f\left( x\right)$ 的定义域为 $\left( {-\dfrac{\pi }{2},\dfrac{\pi }{2}}\right)$ ,其导函数为 ${f}^{\prime }\left( x\right)$ ,当 $0 < x < \dfrac{\pi }{2}$ 时,有 ${f}^{\prime }\left( x\right) \cos x + f\left( x\right) \sin x < 0$ 成立，则关于 $x$ 的不等式 $f\left( x\right)  < {2f}\left( \dfrac{\pi }{3}\right) \cos x$ 的解集为\blkda[4]{$\left( {-\dfrac{\pi }{2}, - \dfrac{\pi }{3}}\right)  \cup  \left( {\dfrac{\pi }{3},\dfrac{\pi }{2}}\right)$}

        \begin{jieda}
            设$F(x) = \dfrac{f(x)}{cosx}$，则$F(x)$为偶函数
            
            $F'(x) = \dfrac{f'(x)cosx+f(x)sinx}{cos^2x}$，故$F(x)$在$0 < x < \dfrac{\pi }{2}$严格减

            $f\left( x\right)  < {2f}\left( \dfrac{\pi }{3}\right) \cos x \Longleftrightarrow F(x) < F\left(\dfrac{\pi}{3}\right)$，故根据$F(x)$性质可知解集为$\left( {-\dfrac{\pi }{2}, - \dfrac{\pi }{3}}\right)  \cup  \left( {\dfrac{\pi }{3},\dfrac{\pi }{2}}\right)$
        \end{jieda}

        \item 定义在 $\left( {-\dfrac{\pi }{2},0}\right)  \cup  \left( {0,\dfrac{\pi }{2}}\right)$ 上的奇函数 $f\left( x\right)$ 的导函数为 ${f}^{\prime }\left( x\right)$ ,且当 $x \in  \left( {0,\dfrac{\pi }{2}}\right)$ 时, ${f}^{\prime }\left( x\right) \tan x - f\left( x\right)  > 0$ ，则不等式 $f\left( x\right)  < {2f}\left( \dfrac{\pi }{6}\right) \sin x$ 的解集为\blkda[4]{$\left( {-\dfrac{\pi }{2}, - \dfrac{\pi }{6}}\right)  \cup  \left( {0,\dfrac{\pi }{6}}\right)$}.
        \begin{jieda}
            $x \in  \left( {0,\dfrac{\pi }{2}}\right)$ 时, ${f}^{\prime }\left( x\right) \tan x - f\left( x\right)  > 0 \Longleftrightarrow f'(x)sinx-f(x)cosx > 0$

            设$F(x) = \dfrac{f(x)}{sinx}$，$F'(x) = \dfrac{f'(x)sinx-f(x)cosx}{sin^2x} > 0$，则$F(x)$为偶函数，在$\left( {-\dfrac{\pi }{2},0}\right)$严减，在$\left( {0,\dfrac{\pi }{2}}\right)$严增

            当$x \in \left(0, \dfrac{\pi}{2}\right)$时，$f\left( x\right)  < {2f}\left( \dfrac{\pi }{6}\right) \sin x \Longleftrightarrow F(x) < F\left(\dfrac{\pi}{6}\right)$，当$x \in \left(-\dfrac{\pi}{2}, 0\right)$时，$f\left( x\right)  < {2f}\left( \dfrac{\pi }{6}\right) \sin x \Longleftrightarrow F(x) > F\left(\dfrac{\pi}{6}\right)$，故原不等式解集为$\left( {-\dfrac{\pi }{2}, - \dfrac{\pi }{6}}\right)  \cup  \left( {0,\dfrac{\pi }{6}}\right)$
        \end{jieda}

        \item 已知偶函数 $f\left( x\right)$ 的定义域是 $\left( {-\infty ,0}\right)  \cup  \left( {0, + \infty }\right)$ ,其导函数为 $f\left( x\right)$ ,对定义域内的任意 $x$ ,都有 ${2f}\left( x\right)  + x{f}^{\prime }\left( x\right)  > 2$ 成立,则不等式 ${x}^{2}f\left( x\right)  - {4f}\left( 2\right)  < {x}^{2} - 4$ 的解集为\blkda[4]{$\left( {-2,0}\right)  \cup  \left( {0,2}\right)$}.
        \begin{jieda}
            设$F(x) = x^2f(x)-x^2$，则$F'(x) = 2xf(x)+x^2f'(x)-2x = x[{2f}\left( x\right)  + x{f}^{\prime }\left( x\right)-2]$，故$F(x)$为偶函数，在$(-\infty, 0)$严减，在$(0, +\infty)$严增

            ${x}^{2}f\left( x\right)  - {4f}\left( 2\right)  < {x}^{2} - 4 \Longleftrightarrow F(x) < F(2)$，故解集为$\left( {-2,0}\right)  \cup  \left( {0,2}\right)$
        \end{jieda}
        
        \item 已知定义在 ${R}$ 上的偶函数 $f\left( x\right)$ ,其导函数为 ${f}^{\prime }\left( x\right)$ ,若 $x{f}^{\prime }\left( x\right)  - {2f}\left( x\right)  > 0,f\left( {-1}\right)  = \dfrac{1}{2}$ ,则不等式 ${2f}\left( x\right)  < {x}^{2}$ 的解集是\blkda{$(-1, 1)$}.
        \begin{jieda}
            设$F(x) = \dfrac{f(x)}{x^2}$，则$F'(x) = \dfrac{f'(x)x^2-2xf(x)}{x^2} = \dfrac{xf'(x)-2f(x)}{x}$，故$F(x)$为偶函数，在$(-\infty, 0)$严减，在$(0, +\infty)$严增，$F(-1) = f(-1) = \dfrac{1}{2}$

            当$x = 0$时，不等式满足，当$x \neq 0$时${2f}\left( x\right)  < {x}^{2} \Longleftrightarrow F(x) < \dfrac{1}{2}$，故解集为$(-1, 1)$

            \ps $x{f}^{\prime }\left( x\right)  - {2f}\left( x\right)  > 0$中有商的导数结构，因为$f(x)$系数为2，考虑构造$y = \dfrac{f(x)}{x^2}$
        \end{jieda}
        
        
        \item 函数 $f\left( x\right)$ 的导函数 ${f}^{\prime }\left( x\right)$ 满足 ${2f}\left( x\right)  + {f}^{\prime }\left( x\right)  > 2$ ,且 $f\left( 1\right)  = {2025}$ ,则不等式 $f\left( x\right)  > 1 + \dfrac{2024}{{{e}}^{{2x} - 2}}$ 的解集是\blkda{$\left( {1, + \infty }\right)$}
        \begin{jieda}
            设$F(x) = f(x)e^{2x-2}-e^{2x-2}$，则$F'(x) = [f'(x)+2f(x)-2]e^{2x-2}$，故$F'(x) > 0$，$F(1) = 2024$

            $f\left( x\right)  > 1 + \dfrac{2024}{{{e}}^{{2x} - 2}} \Longleftrightarrow F(x) > 2024$，故解集为$\left( {1, + \infty }\right)$

            \ps 尝试构造过的函数：$y = f(x)e^{2x}, y = f(x)e^{2x-2}$
        \end{jieda}
        \item 已知函数 $f\left( x\right)$ 及其导函数 ${f}^{\prime }\left( x\right)$ 的定义域为 ${R},f\left( 1\right)  =  - 1$ ,且 ${f}^{\prime }\left( x\right)  < {3f}\left( x\right)  + 6$ ,则不等式 $f\left( {\ln x}\right)  + 2 > \dfrac{{x}^{3}}{{{e}}^{3}}$ 的解集为\blkda{$(0, e)$}
        \begin{jieda}
            令$t = lnx$，$F(x) = f(x)+2$，$f\left( {\ln x}\right)  + 2 > \dfrac{{x}^{3}}{{{e}}^{3}} \Longleftrightarrow F(t) > \dfrac{e^{3t}}{e^3}$，$F(1) = 1$

            ${f}^{\prime }\left( x\right)  < {3f}\left( x\right)  + 6 \Longleftrightarrow F'(x) < 3F(x)$，$\left(\dfrac{F(x)}{e^{3x}}\right)' = \dfrac{F'(x) - 3F(x)}{e^{3x}} < 0$

            $F(t) > \dfrac{e^{3t}}{e^3} \Longleftrightarrow \dfrac{F(t)}{e^{3t}} > \dfrac{F(1)}{e^3}$，故$t < 1$，即$lnx < 1$，解得$x \in (0, e)$
        \end{jieda}
        \item 已知连续函数 $f\left( x\right)$ 及其导函数 ${f}^{\prime }\left( x\right)$ 的定义域均为 ${R}$ ,且 $\left( {x - 2}\right) \left\lbrack  {{f}^{\prime }\left( x\right)  - f\left( x\right) }\right\rbrack  > 0 (x \neq 2)$ , $f\left( {4 - x}\right)  = f\left( x\right) {{e}}^{4 - {2x}}$ ,则不等式 ${{e}}^{3}f\left( {\ln x}\right)  < {xf}\left( 3\right)$ 的解集是\blkda{$\left( {{e},{{e}}^{3}}\right)$}
        \begin{jieda}
            % 设$F(x) = (x-2)f(x)$，则$F'(x) > 0$

            $f\left( {4 - x}\right)  = f\left( x\right) {{e}}^{4 - {2x}} \Longleftrightarrow f(4-x)e^x = f(x)e^{4-x} \Longleftrightarrow \dfrac{f(4-x)}{e^{4-x}} = \dfrac{f(x)}{e^x}$

            设$F(x) = f(x)e^{-x}$，则$F(4-x) = F(x)$，即$F(x)$关于$x = 2$对称. 

            $F'(x) = [f'(x)-f(x)]e^{-x}$

            由$\left( {x - 2}\right) \left\lbrack  {{f}^{\prime }\left( x\right)  - f\left( x\right) }\right\rbrack  > 0$可知在$x>2$时$[f'(x)-f(x)] > 0$，在$x < 2$时$[f'(x)-f(x)] < 0$

            因此连续函数$F(x)$在$(2, +\infty)$严增，在$(-\infty, 2)$严减

            ${{e}}^{3}f\left( {\ln x}\right)  < {xf}\left( 3\right) \Longleftrightarrow F(lnx) < F(3)$，根据$F(x)$的性质可知不等式即$lnx \in (1,3)$，故$x \in \left( {{e},{{e}}^{3}}\right)$
        \end{jieda}


        \item 已知函数 $f\left( x\right)$ 的定义域为 $\left( {0, + \infty }\right) ,{f}^{\prime }\left( x\right)$ 为其导函数,若 $\forall x \in  \left( {0, + \infty }\right)$ , $f\left( x\right)  > \left\lbrack  {f\left( x\right)  - x{f}^{\prime }\left( x\right) }\right\rbrack  \ln x$ ，则不等式 $f\left( x\right) \left( {{{e}}^{x - 1} - 1}\right)  > 0$ 的解集是\blkda{$(1, +\infty)$}.
        \begin{jieda}
            设$F(x) = \dfrac{f(x)lnx}{x}$
            
            则$F'(x) = \dfrac{\left[f'(x)lnx + \dfrac{f(x)}{x}\right]x-f(x)lnx}{x^2} = \dfrac{xlnx\cdot f'(x)+(1-lnx)f(x)}{x^2}$，故$F(x)$在$(0, +\infty)$严增，$F(1) = 0$，故$f(x) > 0$，因此$f\left( x\right) \left( {{{e}}^{x - 1} - 1}\right)  > 0$ 的解集是$(1, +\infty)$

            \ps $f\left( x\right)  > \left\lbrack  {f\left( x\right)  - x{f}^{\prime }\left( x\right) }\right\rbrack  \ln x$中存在$\dfrac{f(x)}{x}$的导数结构和$lnx$
        \end{jieda}
        \item 填写下表
        
            \begin{center}
            \begin{tabular}{|c|c|}
            \hline
            不等式 & 构造函数 \\
            \cline{1-2}
            $x{f}^{\prime }\left( x\right)  + f\left( x\right)  > 0\left( {\text{ 或 } < 0}\right)$ & \blkda[6]{$F\left( x\right)  = {xf}\left( x\right)$} \\
            \cline{1-2}
            $x{f}^{\prime }\left( x\right)  - f\left( x\right)  > 0\left( {\text{ 或 } < 0}\right)$ & \blkda[6]{$F\left( x\right)  = \dfrac{f\left( x\right) }{x}$} \\
            \cline{1-2}
            ${f}^{\prime }\left( x\right)  + f\left( x\right)  > 0\left( {\text{ 或 } < 0}\right)$ & \blkda[6]{$F\left( x\right)  = {{e}}^{x}f\left( x\right)$} \\
            \cline{1-2}
            ${f}^{\prime }\left( x\right)  - f\left( x\right)  > 0$ (或 $< 0$ ) & \blkda[6]{$F\left( x\right)  = \dfrac{f\left( x\right) }{{{e}}^{x}}$} \\
            \cline{1-2}
            $x{f}^{\prime }\left( x\right)  + {nf}\left( x\right)  > 0\left( {\text{ 或 } < 0}\right)$ & \blkda[6]{$F\left( x\right)  = {x}^{n}f\left( x\right)$} \\
            \cline{1-2}
            $x{f}^{\prime }\left( x\right)  - {nf}\left( x\right)  > 0\left( {\text{ 或 } < 0}\right)$ & \blkda[6]{$F\left( x\right)  = \dfrac{f\left( x\right) }{{x}^{n}}$} \\
            \cline{1-2}
            ${f}^{\prime }\left( x\right)  + {nf}\left( x\right)  > 0$ (或 $< 0$ ) & \blkda[6]{$F\left( x\right)  = {e}^{nx}f\left( x\right)$} \\
            \cline{1-2}
            ${f}^{\prime }\left( x\right)  - {nf}\left( x\right)  > 0$ (或 $< 0$ ) & \blkda[6]{$F\left( x\right)  = \dfrac{f\left( x\right) }{{e}^{nx}}$} \\
            \cline{1-2}
            $f\left( x\right)  + {f}^{\prime }\left( x\right) \tan x > 0$ (或 $< 0$ ) & \blkda[6]{$F\left( x\right)  = \sin {xf}\left( x\right)$} \\
            \cline{1-2}
            $f\left( x\right)  - {f}^{\prime }\left( x\right) \tan x > 0$ (或 $< 0$ & \blkda[6]{$F\left( x\right)  = \dfrac{f\left( x\right) }{\sin x}$} \\
            \cline{1-2}
            ${f}^{\prime }\left( x\right)  - f\left( x\right) \tan x > 0\left( {\text{ 或 } < 0}\right)$ & \blkda[6]{$F\left( x\right)  = \cos {xf}\left( x\right)$} \\
            \cline{1-2}
            ${f}^{\prime }\left( x\right)  + f\left( x\right) \tan x > 0\left( {\text{ 或 } < 0}\right)$ & \blkda[6]{$F\left( x\right)  = \dfrac{f\left( x\right) }{\cos x}$} \\
            \cline{1-2}
            \hline
            \end{tabular}
            \end{center}
    \end{enumerate}
    
    \item 割线斜率
    \begin{enumerate}
        \item $f\left( x\right)  = {x}^{3} + a{x}^{2} + {2x} - {a}^{2}$ ,若 $\forall {x}_{1},{x}_{2} \in  \left\lbrack  {1,2}\right\rbrack$，${x}_{1} < {x}_{2}$ ,都有 $\dfrac{f\left( {x}_{1}\right)  - f\left( {x}_{2}\right) }{{x}_{1} - {x}_{2}} < 2$ 成立,则 $a$ 的最大值为\blkda{$-3$}
        \begin{jieda}
            \begin{enumerate}
                \item 利用拉格朗日中值定理逆变换

                $ {f}^{\prime }\left( x\right)  \leq  2$ 在 $\left\lbrack  {1,2}\right\rbrack$ 上恒成立，即$3{x}^{2} + {2ax} + 2 \leq  2$ 在 $x \in  \left\lbrack  {1,2}\right\rbrack$ 上恒成立，故$a \leq   - \dfrac{3}{2} x$ 在 $\left\lbrack  {1,2}\right\rbrack$ 上恒成立
                $\therefore a =  - 3$

                \item 构造函数
                
                $\dfrac{f\left( {x}_{1}\right)  - f\left( {x}_{2}\right) }{{x}_{1} - {x}_{2}} < 2 \Longleftrightarrow f(x_1) - f(x_2) > 2x_1 - 2x_2 \Longleftrightarrow f(x_1)-2x_1 > f(x_2)-2x_2$

                故$F(x) = f(x)-2x$在$[1, 2]$内为减函数，故$F'(x) = 3x^2+2ax \leq 0$在$[1, 2]$恒成立，即$a \leq \dfrac{3x}{2}$在 $\left\lbrack  {1,2}\right\rbrack$ 上恒成立
                $\therefore a =  - 3$
            \end{enumerate}
            \ps 拉格朗日中值定理：如果 $f\left( x\right)$ 满足在闭区间 $\left\lbrack  {a,b}\right\rbrack$ 上连续，在开区间$(a,b)$内可导，则至少有一点($a<\xi<b$),使等式 $\dfrac{f\left( b\right)  - f\left( a\right) }{b - a} = f'\left( \xi\right)$成立；即每一条割线都能找到与之平行的切线.

            拉格朗日中值定理不存在逆定理，即并非每一条切线都能找到与之平行的割线，如$y = x^3$在$x = 0$的切线则不存在与之平行的割线

            如果 $f\left( x\right)$ 满足在闭区间 $\left\lbrack  {a,b}\right\rbrack$ 上连续，在开区间$(a,b)$内可导（存在二阶导）则在$\left\lbrack  {a,b}\right\rbrack$内端点和拐点（凹凸性改变的点，通常为二阶导变号零点）处，使得导函数取得最值的位置的切线找不到与之平行的割线（不考虑局部线性的特殊情况）.

            故在填选中，可以进行如下转化：
            \begin{dingautolist}{172}
                \item $\dfrac{f\left( {x}_{1}\right)  - f\left( {x}_{2}\right) }{{x}_{1} - {x}_{2}} > m \Longrightarrow  f'\left( x\right)  \geq  m$
                
                \item $\dfrac{f\left( {x}_{1}\right)  - f\left( {x}_{2}\right) }{{x}_{1} - {x}_{2}} < m \Longrightarrow  {f}^{\prime }\left( x\right)  \leq  m$
            \end{dingautolist}
        \end{jieda}
    
        \item 设函数 $f\left( x\right)  = \dfrac{{x}^{2}}{2} - {2a}\ln x + \left( {a - 2}\right) x,a \in  R$ ,若果存在实数 $a$ ,对于不相等正数 ${x}_{1},{x}_{2}$ ,不等式 $\dfrac{f\left( {x}_{1}\right)  - f\left( {x}_{2}\right) }{{x}_{1} - {x}_{2}} > a$ 恒成立,则实数 $a$ 的取值范围为\blkda{$\left(-\infty, -\dfrac{1}{2}\right]$}
        \begin{jieda}
            \begin{enumerate}
                \item 利用拉格朗日中值定理逆变换

                $ {f}^{\prime }\left( x\right)  \geq  a$ 恒成立，即$x-\dfrac{2a}{x}+a-2 \geq  a$ 在$x > 0$恒成立，故$x^2-2x \geq 2a$在$x > 0$恒成立，解得$a \in \left(-\infty, -\dfrac{1}{2}\right]$

                \item 构造函数
                不妨设$x_1 > x_2$
                
                $\dfrac{f\left( {x}_{1}\right)  - f\left( {x}_{2}\right) }{{x}_{1} - {x}_{2}} > a \Longleftrightarrow f(x_1)-ax_1 > f(x_2)-ax_2$，即$F(x) = f(x)-ax$在$(0, +\infty)$为增函数，故$F'(x) = x-\dfrac{2a}{x}-2 \geq  0$在$x > 0$恒成立，故$x^2-2x \geq 2a$在$x > 0$恒成立，解得$a \in \left(-\infty, -\dfrac{1}{2}\right]$
            \end{enumerate}
        \end{jieda}

        \item $f\left( x\right)  = \left( {-x - 1}\right) \ln x + a\left( {x - 1}\right)$ ,若对 $x \in  \left( {1, + \infty }\right)$都有 $f\left( x\right)  < 0$ 成立,则 $a$ 的取值范围\blkda{$\left(-\infty, 2\right]$}
        \begin{jieda}
            $x \in  \left( {1, + \infty }\right)$ 都有 $f\left( x\right)  < 0$ 即 $a < \dfrac{\left( {x + 1}\right) \ln x}{x - 1} = \dfrac{\left( {x + 1}\right) \ln x - 0}{x - 1}$
            
            令 $g\left( x\right)  = \left( {x + 1}\right) \ln x$，可知 $a < \dfrac{g\left( x\right)  - g\left( 1\right) }{x - 1}$，故${g}^{\prime }\left( x\right)  \geq  a$ 对 $x \in  \left( {1, + \infty }\right)$ 恒成立 
            
            ${g}^{\prime }\left( x\right)  = \ln x + \dfrac{1}{x} + 1, g''\left( x\right)  = \dfrac{1}{x} - \dfrac{1}{{x}^{2}} = \dfrac{x - 1}{{x}^{2}} > 0$
            
            则 ${g}^{\prime }\left( x\right)$ 在 $\left( {1, + \infty }\right)$严增，故 ${g}^{\prime }\left( x\right)_{\min}  = {g}^{\prime }\left( 1\right)  = 2$，故$a \leq  2$ 
        \end{jieda}

        \item 已知$f(x), g(x)$都是定义域为$\mathbf{R}$的函数，且$f(x)$为奇函数，$g(x)$为偶函数，满足$f(x)+g(x) = ax^2+x+2$，若对任意的$1 < x_1 < x_2 < 2$，都有$\dfrac{g(x_1)-g(x_2)}{x_1-x_2} > -3$成立，则实数$a$的取值范围是\blkda{$\left[\dfrac{3}{4}, +\infty\right)$}
        \begin{jieda}
            $f(x) = x$，$g(x) = ax^2+2$
            \begin{enumerate}
                \item 方法1

                $g'(x) \geq -3$在$(1, 2)$恒成立，即$2ax \geq -3$在$(1, 2)$恒成立，故$a \geq \dfrac{3}{4}$

                \item 方法2

                $\dfrac{a(x_1^2-x_2^2)}{x_1-x_2} > -3$，即$a(x_1+x_2) > -3$故$a \geq -\dfrac{3}{4}$
            \end{enumerate}
        \end{jieda}

        
        \item 已知函数 $f\left( x\right)  = {x}^{2}\ln x$ .
        \begin{enumerate}
            \item 求函数 $f\left( x\right)$ 的极值
            
            \item 证明:对任意的 $x \in  \left( {0, + \infty }\right)$ ，有 $f\left( x\right)  \geq  x - 1$
            
            \item 若 ${x}_{1},{x}_{2} \in  \left( {0,1}\right)$ ，证明: $\left| {f\left( {x}_{1}\right)  - f\left( {x}_{2}\right) }\right|  \leq  \left| {{x}_{1} - {x}_{2}}\right|$ .
        \end{enumerate}

        \begin{jieda}
            \begin{enumerate}
                \item 因为 ${f}^{\prime }\left( x\right)  = {2x}\ln x + x = x\left( {2\ln x + 1}\right)$ ,
                
                令 ${f}^{\prime }\left( x\right)  = x\left( {2\ln x + 1}\right)  = 0,x = \dfrac{1}{\sqrt{{e}}}$ ,
                
                又因为 $x \in  \left( {0,\dfrac{1}{\sqrt{{e}}}}\right) ,{f}^{\prime }\left( x\right)  < 0,f\left( x\right)$ 单调递减; $x \in  \left( {\dfrac{1}{\sqrt{{e}}}, + \infty }\right) ,{f}^{\prime }\left( x\right)  > 0,f\left( x\right)$ 单调递增;
                
                所以 $f\left( x\right)$ 的极小值为 $f\left( \dfrac{1}{\sqrt{{e}}}\right)  =  - \dfrac{1}{2{e}}$ ,无极大值.

                \item 令 $t\left( x\right)  = x\ln x - 1 + \dfrac{1}{x}$ ，$f\left( x\right)  \geq  x - 1$即$t(x) \geq 0$
                
                可得 ${t}^{\prime }\left( x\right)  = \ln x + 1 - \dfrac{1}{{x}^{2}}$，$t''\left( x\right)  = \dfrac{1}{x} + \dfrac{2}{{x}^{3}} > 0$，故$t'\left( x\right)$ 单调递增, $t'\left( 1\right)  = 0$

                $t\left( x\right)$ 在$(0, 1)$单调递减，在$\left( {1, + \infty }\right)$ 单调递增，所以 $t{\left( x\right) }_{\min } = t\left( 1\right)  = 0$，得证.

                \item 不妨设$x_1 > x_2$，则$x_2-x_1 \leq f(x_1) - f(x_2) \leq x_1-x_2$，即$\left\{\begin{array}{l}
                     f(x_1)-x_1 \leq f(x_2)-x_2 \\
                     f(x_1)+x_2 \geq f(x_2)+x_2
                \end{array}\right.$

                即$F_1(x) = f(x)-x$在$(0, 1)$内为减函数，$F_2(x) = f(x)+x$在$(0, 1)$内为增函数

                即$F_1'(x) = {2x}\ln x + x-1 \leq 0$，$F_2'(x) = {2x}\ln x + x+1 \geq 0$

                ${2x}\ln x + x-1 \leq 0 \Longleftrightarrow 2lnx+1-\dfrac{1}{x} \leq 0$，令$G_1(x) = 2lnx+1-\dfrac{1}{x}$，则$G_1'(x) = \dfrac{2}{x}+\dfrac{1}{x^2} = \dfrac{2x+1}{x^2} > 0$，因此$G_1(x) < G_1(1) = 0$，故前者成立

                ${2x}\ln x + x+1 \leq 0 \Longleftrightarrow 2lnx+1+\dfrac{1}{x} \leq 0$，令$G_2(x) = 2lnx+1+\dfrac{1}{x}$，则$G_2'(x) = \dfrac{2}{x}-\dfrac{1}{x^2} = \dfrac{2x-1}{x^2}$，故$G_2'(x)$在$\left(0, \dfrac{1}{2}\right),G_2'(x) < 0$，在$\left(\dfrac{1}{2}, 1\right), G_2'(x) > 0$，故$G_2(x) \geq G_2\left(\dfrac{1}{2}\right) = -2ln2+3 > 0$，故后者成立.
            \end{enumerate}
        \end{jieda}
    \end{enumerate}
    
    \item 函数零点 (方程根) 的个数问题
    \begin{enumerate}
        \item 已知函数$f(x) = 4x^3+3tx^2-6t^2x+t-1, x \in \mathbf{R}$，其中$t \in \mathbf{R}$
        \begin{enumerate}
            \item 当$t \neq 0$时，求函数$y = f(x)$的单调区间
            \item 证明：对任意$t \in (0, +\infty), y = f(x)$在区间$(0, 1)$上均存在零点.
        \end{enumerate}
        \begin{jieda}
            \begin{enumerate}
                \item $f'(x) = 12x^2+6tx - 6t^2$，令$f'(x) = 0$，解得$x_1 = -t, x_2 = \dfrac{t}{2}$
                \begin{dingautolist}{172}
                    \item 若$t < 0$，则$f(x)$在$\left(-\infty, \dfrac{t}{2}\right)$和$\left(-t, +\infty\right)$内严格增，在$\left(\dfrac{t}{2}, -t\right)$严格减
                    \item 若$t > 0$，则$f(x)$在$\left(-\infty, -t\right)$和$\left(\dfrac{t}{2}, +\infty\right)$内严格增，在$\left(-t, \dfrac{t}{2}\right)$严格减
                \end{dingautolist}
                \item $f(0) = t-1, f(1) = 3+4t-6t^2$当$t \in (0, 1) \cup \left(\dfrac{2+\sqrt{22}}{6}, +\infty \right)$，端点处满足函数零点存在定理，因此$f(x)$在$(0, 1)$内必有零点. \\
                若$t \in \left[1, \dfrac{2+\sqrt{22}}{6}\right]$，则$f(x)$在$\left(0, \dfrac{t}{2}\right)$严格减，在$\left(\dfrac{t}{2}, 1\right)$严格增，故在$(0, 1)$存在极小值$f\left(\dfrac{t}{2}\right) = -\dfrac{7}{4}t^3+(t-1) < 0$，又因为此时$f(0) > 0$，因此在$\left(0, \dfrac{t}{2}\right)$内满足零点存在定理，因此$f(x)$在$(0, 1)$内必有零点.
            \end{enumerate}
        \end{jieda}
        \item 下列给出的四个函数中，零点的个数最多的是\dotfill{(\kh{B})}
        \xx{$y = \dfrac{2x}{x - 1}$}{$y = \ln x - \dfrac{1}{3}x$}{$y = {{e}}^{x} - \left( {x + 1}\right)$}{$y = {2}^{x} + \lg x$}

        \begin{jieda}
            A. $y = \dfrac{2x}{x - 1} = 0$ ,得 $x = 0$ ,有 1 个零点
            
            B. $f\left( x\right)  = \ln x - \dfrac{1}{3}x,{f}^{\prime }\left( x\right)  = \dfrac{1}{x} - \dfrac{1}{3} = 0$ ,得 $x = 3$ ,当 $0 < x < 3$ 时, ${y}^{\prime } > 0$ ,函数单调递增，当 $x > 3$ 时， ${y}^{\prime } < 0$ ，函数单调递减，当 $x = 3$ 时，函数取得最大值 $\ln 3 - 1 > 0$ ，且 $f\left( 1\right)  =  - \dfrac{1}{3} < 0$ ，所以 $x \in  \left( {1,3}\right)$ 只有 1 个零点，$f\left( {{e}}^{2}\right)  = 2 - \dfrac{{{e}}^{2}}{3} < 0$ ，所以在区间 $\left( {3,{{e}}^{2}}\right)$ 只有 1 个零点，所以函数 $y = \ln x - \dfrac{1}{3}x$ 有 2 个零点
            
            C. $y = {{e}}^{x} - \left( {x + 1}\right) ,{y}^{\prime } = {{e}}^{x} - 1 = 0$ ,得 $x = 0$，当 $x < 0$ 时， ${y}^{\prime } < 0$ ，函数单调递减，当 $x > 0$ 时， ${y}^{\prime } > 0$ ，函数单调递增，所以当 $x = 0$ 时，函数取得最小值 0，所以 $y = {{e}}^{x} - \left( {x + 1}\right)$ 只有 1 个零点，
            
            D. $f\left( x\right)  = {2}^{x} + \lg x$ 在定义域 $\left( {0, + \infty }\right)$ 单调递增， $f\left( \dfrac{1}{100}\right)  = {2}^{\frac{1}{100}} - 2 < 0,f\left( 1\right)  = 2 > 0$，所以函数 $y = {2}^{x} + \lg x$ 只有 1 个零点
        \end{jieda}
        
        \item 函数 $f\left( x\right)  = \dfrac{2}{3}{x}^{4} - x\ln x - x$ 的零点个数为\blkda{2}
        \begin{jieda}
            由题可得 $x \in  \left( {0, + \infty }\right)$ ,故令 $f\left( x\right)  = \dfrac{2}{3}{x}^{4} - x\ln x - x = 0$，即 $\dfrac{2}{3}{x}^{3} - \ln x - 1 = 0$ . 令 $g\left( x\right)  = \dfrac{2}{3}{x}^{3} - \ln x - 1$，则 ${g}^{\prime }\left( x\right)  = 2{x}^{2} - \dfrac{1}{x} = \dfrac{2{x}^{3} - 1}{x}$
            
            由 ${g}^{\prime }\left( x\right)  > 0 \Longrightarrow  x > {2}^{-\frac{1}{3}},{g}^{\prime }\left( x\right)  < 0 \Longrightarrow  0 < x < {2}^{-\frac{1}{3}}$ ,
            
            所以在 $\left( {{2}^{-\frac{1}{3}}, + \infty }\right)$ 单调递增,在 $\left( {0,{2}^{-\frac{1}{3}}}\right)$ 递减，又 $g\left( 1\right) < 0,g\left( \dfrac{1}{{e}}\right)   > 0,g\left( 3\right) > 0$ ,
            
            $g\left( x\right)$ 在 $\left( {0,{2}^{-\frac{1}{3}}}\right)$ 与 $\left( {{2}^{-\frac{1}{3}}, + \infty }\right)$ 分别有一个零点，所以 $g\left( x\right)$ 有两个零点，故 $f\left( x\right)$ 有两个零点
        \end{jieda}

        \item 若函数 $f\left( x\right)  = {e}^{x} - a{x}^{2}$ 有 3 个零点,则实数 $a$ 的取值范围为\blkda{$\left(\dfrac{e^2}{4}, +\infty\right)$}
        \begin{jieda}
            因为$f(0) \neq 0$，所以$f\left( x\right)  = {e}^{x} - a{x}^{2}$ 有 3 个零点即$y = a$和$g(x) = \dfrac{e^x}{x^2}$ 有 3 个交点，$g'(x) = \dfrac{e^x(x-2)}{x^3}$
            \begin{dingautolist}{172}
                \item 当$x < 0$时，$g'(x) > 0$，$g(x)$严增，值域为$(0, +\infty)$
                \item 当$x \in (0, 2)$时，$g'(x) < 0$，$g(x)$严减，值域为$\left(\dfrac{e^2}{4}, +\infty\right)$
                \item 当$x > 2$时，$g'(x) > 0$，$g(x)$严增，值域为$\left(\dfrac{e^2}{4}, +\infty\right)$
            \end{dingautolist}
            故可知$a \in \left(\dfrac{e^2}{4}, +\infty\right)$满足要求.
        \end{jieda}

        \item 已知函数 $f\left( x\right)  = \ln \dfrac{2 - x}{2 + x} - {2mx}$ 有且仅有 3 个零点,则 $m$ 的取值范围为\blkda[3]{$\left( {-\infty , - \dfrac{1}{2}}\right)$}
        \begin{jieda}
            由 $f\left( x\right)  = 0$ ,得 $ln\left(\dfrac{2-x}{2+x}\right) = {2mx}$ ,令函数 $h\left( x\right)  = \ln \dfrac{2 - x}{2 + x}$ ,其定义域为(-2,2),
            ${h}\left( {-{x}}\right)  + {h}\left( {x}\right)  = 0$ ,函数 ${h}\left( {x}\right)$ 为奇函数，$h(x) = ln\left(\dfrac{2-x}{2+x}\right) = ln\left(\dfrac{4}{2+x}-1\right)$严减
            
            依题意,直线 ${y} = 2{{mx}}$ 与函数 ${h}\left( {x}\right)$ 的图象有且仅有 3 个交点.
            求导可得曲线 ${y} = {h}\left( {x}\right)$ 在点(0,0)处的切线方程为 ${y} =  - {x}$，作出 ${h}\left( {x}\right)$ 的图象,如图:

            \begin{center}
                \includegraphics[width=0.25\linewidth]{bodyxb//src/lingdian.png}
            \end{center}
            
            观察图象知,当 $2{m} <  - 1$ 时,直线 ${y} = 2{m}{x}$ 与函数 ${h}\left( {x}\right)  = ln\left(\dfrac{2-x}{2+x}\right)$ 的图象有且仅有 3 个交点，所以 ${m}$ 的取值范围是 $\left( {-\infty , - \dfrac{1}{2}}\right)$ .
        \end{jieda}
        
        
        \item 已知函数 $f\left( x\right)  = \ln \left( {1 + x}\right)  + \cos x - x,x \in  \left( {-1,\pi }\right)$ .
        \begin{enumerate}
            \item 求曲线 $y = f\left( x\right)$ 在点 $\left( {0,f\left( 0\right) }\right)$ 处的切线方程
            \item 求 $f\left( x\right)$ 零点的个数
        \end{enumerate}
        \begin{jieda}
            \begin{enumerate}
                \item 因为 $f\left( x\right)  = \ln \left( {1 + x}\right)  + \cos x - x,x \in  \left( {-1,\pi }\right)$，所以 ${f}^{\prime }\left( x\right)  = \dfrac{1}{1 + x} - \sin x - 1,{f}^{\prime }\left( 0\right)  = 0,f\left( 0\right)  = 1$，所以切线方程为 $y - f\left( 0\right)  = {f}^{\prime }\left( 0\right) \left( {x - 0}\right)$ ,即 $y = 1$ .
                
                \item ${f}^{\prime }\left( 0\right)  = 0$ ，当 $x \in  \left( {0,\pi }\right)$ ， $\dfrac{1}{1 + \pi } < \dfrac{1}{1 + x} < 1$ ， $- 1 <  - \sin x < 0$ ，
                
                所以 ${f}^{\prime }\left( x\right)  = \dfrac{1}{1 + x} - \sin x - 1 < 0$ ,即 $f\left( x\right)$ 在 $\left( {0,\pi }\right)$ 上单调递减

                令 $h\left( x\right)  = \ln x - x,{h}^{\prime }\left( x\right)  = \dfrac{1}{x} - 1 = \dfrac{1 - x}{x}$ ,
                
                当 $x \in  \left( {0,1}\right)$ 时， ${h}^{\prime }\left( x\right)  > 0,h\left( x\right)$ 单调递增，且 $h\left( 1\right)  =  - 1$ ，
                
                当 $x \in  \left( {1, + \infty }\right)$ 时， ${h}^{\prime }\left( x\right)  < 0,h\left( x\right)$ 单调递减；
                
                所以 $h{\left( x\right) }_{\max } = h\left( 1\right)  =  - 1$ ,即 $h\left( x\right)  = \ln x - x < 0$ ,
                
                所以 $h\left( 1+\pi \right)  = \ln \left( {1 + \pi }\right)  - \left( {1 + \pi }\right)  <  - 1 < 0$ ,
                
                所以 $f\left( \pi \right)  = \ln \left( {1 + \pi }\right)  - \left( {1 + \pi }\right)  < 0$ ,
                
                又因为 $f\left( x\right)$ 在 $\left( {0,\pi }\right)$ 单调递减，根据零点存在定理 $f\left( x\right)$ 在 $\left( {0,\pi }\right)$ 有唯一零点.
                
                令 $f''\left( x\right)  =  - \dfrac{1}{{\left( 1 + x\right) }^{2}} - \cos x$，当 $x \in  \left( {-1,0}\right)$ 时， $f''\left( x\right)  < 0$ ，${f}^{\prime }\left( x\right)$ 在 $\left( {-1,0}\right)$ 上单调递减，又因为 ${f}^{\prime }\left( 0\right)  = 0$ ,当 $x \in  \left( {-1,0}\right)$ 时 ${f}^{\prime }\left( x\right)  > 0$ ,即 $f\left( x\right)$ 在 $\left( {-1,0}\right)$ 单调递增,
                
                因此 $f\left( x\right)$ 在 $\left( {-1,0}\right)$ 单调递增,在 $\left( {0,\pi }\right)$ 单调递减.
                
                当 $x \Longrightarrow   - 1$ 时, $\ln \left( {1 + x}\right)  \Longrightarrow   - \infty ,f\left( x\right)  < 0;f\left( 0\right)  = 1$ ,因为 $f\left( x\right)$ 在 $\left( {-1,0}\right)$ 单调递增,
                
                所以根据零点存在定理, $f\left( x\right)$ 在 $x \in  \left( {-1,0}\right)$ 有唯一零点;
                
                综上, $f\left( x\right)$ 在 $\left( {-1,\pi }\right)$ 上有 2 个零点.
            \end{enumerate}
        \end{jieda}
        
        \item 已知函数 $f\left( x\right)  = {{e}}^{x} - a\left( {x + 1}\right) ,a \in  {R}$ .
        \begin{enumerate}
            \item 讨论 $f\left( x\right)$ 的单调性
            \item 若 $a = 1$ ，求曲线 $f\left( x\right)$ 在 $x = 1$ 处的切线方程
            \item 当 $a > 0$ 时，试讨论函数 $f\left( x\right)$ 的零点个数
        \end{enumerate}

        \begin{jieda}
            \begin{enumerate}
                \item ${f}^{\prime }\left( x\right)  = {{e}}^{x} - a$
                
                当 $a \leq  0$ 时， ${f}^{\prime }\left( x\right)  > 0$ 恒成立，故 $f\left( x\right)$ 在 ${R}$ 上单调递增，
                
                当 $a > 0$ 时，令 ${f}^{\prime }\left( x\right)  = 0$ ，解得 $x = \ln a$ ，
                
                所以当 $x \in  \left( {\ln a, + \infty }\right)$ 时， ${f}^{\prime }\left( x\right)  > 0,f\left( x\right)$ 单调递增；
                
                当 $x \in  \left( {-\infty ,\ln a}\right)$ 时， ${f}^{\prime }\left( x\right)  < 0,f\left( x\right)$ 单调递减；
                
                综上,当 $a \leq  0$ 时, $f\left( x\right)$ 在 ${R}$ 上单调递增;
                
                当 $a > 0$ 时, $f\left( x\right)$ 在 $\left( {\ln a, + \infty }\right)$ 上单调递增,在 $\left( {-\infty ,\ln a}\right)$ 上单调递减
                
                \item 由 $a = 1$ ，则 $f\left( x\right)  = {{e}}^{x} - \left( {x + 1}\right)$ ，得 ${f}^{\prime }\left( x\right)  = {{e}}^{x} - 1$ ，
                
                曲线 $f\left( x\right)$ 在 $x = 1$ 处的切线的斜率为 ${f}^{\prime }\left( 1\right)  = {e} - 1,f\left( 1\right)  = {e} - 2$
                
                故曲线 $f\left( x\right)$ 在 $x = 1$ 处的切线的方程为 $y - {e} + 2 = \left( {{e} - 1}\right) \left( {x - 1}\right)$ ,
                
                即 $\left( {{e} - 1}\right) x - y - 1 = 0$ ;
                
                \item 由于 $f\left( x\right)  = 0$ ，即 ${{e}}^{x} - a\left( {x + 1}\right)  = 0$ ， $\therefore \dfrac{1}{a} = \dfrac{x + 1}{{{e}}^{x}}$ ，
                
                即 $f\left( x\right)$ 的零点个数可看作直线 $y = \dfrac{1}{a}$ 与曲线 $y = \dfrac{x + t}{{{e}}^{x}}$ 的交点个数问题;
                
                令 $g\left( x\right)  = \dfrac{x + 1}{{{e}}^{x}}$ ,则 ${g}^{\prime }\left( x\right)  = \dfrac{-x}{{{e}}^{x}}$ ,
                
                当 $x < 0$ 时， ${g}^{\prime }\left( x\right)  > 0,g\left( x\right)$ 在 $\left( {-\infty ,0}\right)$ 上单调递增，
                
                当 $x > 0$ 时, ${g}^{\prime }\left( x\right)  < 0,g\left( x\right)$ 在 $\left( {0, + \infty }\right)$ 上单调递减,
                
                故 $g{\left( x\right) }_{\max } = g\left( 0\right)  = 1$ ,
                
                当 $x <  - 1$ 时, $g\left( x\right)  < 0$ ,当 $x >  - 1$ 时, $g\left( x\right)  > 0$ ,
                
                当 $x$ 趋向于负无穷时, $g\left( x\right)$ 趋向于负无穷,当 $x$ 趋向于正无穷时, $g\left( x\right)$ 趋向于 0,
                
                作出函数 $g\left( x\right)$ 的图象如图:

                \begin{center}
                \includegraphics[width=0.2\textwidth]{bodyxb/src/bo_d5rb33bef24c73bob2k0_3_140_1609_327_214_0.jpg}
                \end{center}

                \begin{dingautolist}{172}
                    \item 当 $0 < \dfrac{1}{a} < 1$ ，即 $a > 1$ 时，直线 $y = \dfrac{1}{a}$ 与曲线 $y = \dfrac{x + 1}{{{e}}^{x}}$ 中的有 2 个交点，函数 $f\left( x\right)$ 的零点个数是 $2$
                    
                    \item 当 $\dfrac{1}{a} = 1$ ,即 $a = 1$ 时,直线 $y = \dfrac{1}{a}$ 与曲线 $y = \dfrac{x + 1}{{{e}}^{x}}$ 的有 1 个交点，函数 $f\left( x\right)$ 的零点个数是 1
                    
                    \item 当 $\dfrac{1}{a} > 1$ ,即 $0 < a < 1$ 时,直线 $y = \dfrac{1}{a}$ 与曲线 $y = \dfrac{x + 1}{{{e}}^{x}}$ 中的无交点，函数 $f\left( x\right)$ 的零点个数是0
                \end{dingautolist}
            \end{enumerate}
        \end{jieda}
    \end{enumerate}
    \item 函数零点 (方程根) 的参数范围问题
    \begin{enumerate}
        \item 已知函数 $f\left( x\right)  = \left\{  \begin{array}{l}  - {x}^{2} + \dfrac{1}{2}x,x < 0 \\  \ln \left( {x + 1}\right) ,x \geq  0 \end{array}\right.$ ，若函数 $y = f\left( x\right)  - {kx}$ 有且只有 3 个零点,则实数 $k$ 的取值范围为\blkda{$\left( {\dfrac{1}{2},1}\right)$}
        \begin{jieda}
            由函数 $f\left( x\right)  = \left\{  \begin{array}{l}  - {x}^{2} + \dfrac{1}{2}x,x < 0 \\  \ln \left( {x + 1}\right) ,x \geq  0 \end{array}\right.$ ,若 $y = f\left( x\right)  - {kx}$ 有且只有 3 个零点,
            
            当 $x = 0$ 时,可得 $f\left( 0\right)  = \ln 1 = 0$ ,可得 $x = 0$ 是 $y = f\left( x\right)  - {kx}$ 的一个零点,
            
            当 $x < 0$ 时，$y = f\left( x\right)  - {kx}$有零点即$- {x}^{2} + \dfrac{1}{2}x = {kx}$有解 ，可得 $x = \dfrac{1}{2} - k < 0$ ，解得 $k > \dfrac{1}{2}$，此时在$x < 0$ 时有一个零点
            
            当 $x > 0$ 时， $f\left( x\right)  = \ln \left( {x + 1}\right)$ ，可得 ${f}^{\prime }\left( x\right)  = \dfrac{1}{x + 1}$ ，可得 ${f}^{\prime }\left( 0\right)  = 1$ ，要使得函数 $y = f\left( x\right)  - {kx}$ 在 $x > 0$ 上有一个零点，即函数 $y = f\left( x\right)$ 与 $y = {kx}$ 的图象有一个公共点,则满足 $0 < k < 1$ 

            \begin{center}
            \includegraphics[width=0.3\textwidth]{bodyxb/src/bo_d5rb33bef24c73bob2k0_4_138_1082_375_301_0.jpg}
            \end{center}
            
            综上可得: $\dfrac{1}{2} < k < 1$ ，即函数 $f\left( x\right)  - {kx}$ 有三个零点时，实数 $k$ 的范围为 $\left( {\dfrac{1}{2},1}\right)$
        \end{jieda}
        
        \item 已知函数 $f\left( x\right)  = {x}^{3} + {2a}{x}^{2} + {bx} + a - 1$ 在 $x =  - 1$ 处取得极值0,其中 $a,b \in  {R}$ .
        \begin{enumerate}
            \item 求 $a,b$ 的值
            \item 当 $x \in  \left\lbrack  {-1,2}\right\rbrack$ 时，方程 $f\left( x\right)  = k$ 有两个不等实数根，求实数 $k$ 的取值范围
        \end{enumerate}

        \begin{jieda}
            \begin{enumerate}
                \item 由 $f\left( x\right)  = {x}^{3} + {2a}{x}^{2} + {bx} + a - 1$ 求导得 ${f}^{\prime }\left( x\right)  = 3{x}^{2} + {4ax} + b$ ,
                
                依题意可知 $\left\{  \begin{matrix} f\left( {-1}\right)  = 0 \\  {f}^{\prime }\left( {-1}\right)  = 0 \end{matrix}\right.$ ,即 $\left\{  \begin{matrix}  - 1 + {2a} - b + a - 1 = 0 \\  3 - {4a} + b = 0 \end{matrix}\right.$ ,解得 $a = b = 1$ ,
                
                此时 $f\left( x\right)  = {x}^{3} + 2{x}^{2} + x,{f}^{\prime }\left( x\right)  = 3{x}^{2} + {4x} + 1$ ,
                
                由 ${f}^{\prime }\left( x\right)  = 3{x}^{2} + {4x} + 1 = \left( {{3x} + 1}\right) \left( {x + 1}\right)  = 0$ 得 $x =  - 1$ 或 $x =  - \dfrac{1}{3}$
                
                当 $x <  - 1$ 时， ${f}^{\prime }\left( x\right)  > 0$ ，函数 $f\left( x\right)$ 递增，
                
                当 $- 1 < x <  - \dfrac{1}{3}$ 时， ${f}^{\prime }\left( x\right)  < 0$ ，函数 $f\left( x\right)$ 递减，
                
                故 $x =  - 1$ 时,函数取得极大值 $f\left( {-1}\right)  = 0$ ,故 $a = b = 1$ .
                
                \item 由 (1) 得 $f\left( x\right)  = {x}^{3} + 2{x}^{2} + x,{f}^{\prime }\left( x\right)  = 3{x}^{2} + {4x} + 1 = \left( {{3x} + 1}\right) \left( {x + 1}\right)$ ,
                
                令 ${f}^{\prime }\left( x\right)  = 0$ ,解得 $x =  - \dfrac{1}{3}$ 或 $x =  - 1$ ,因 $x \in  \left\lbrack  {-1,2}\right\rbrack$ ,
                
                故当 $- 1 < x <  - \dfrac{1}{3}$ 时 ${f}^{\prime }\left( x\right)  < 0$ ，函数 $f\left( x\right)$ 递减，当 $- \dfrac{1}{3} < x < 2$ 时 ${f}^{\prime }\left( x\right)  > 0$ ，函数 $f\left( x\right)$ 递增，
                
                当 $x =  - \dfrac{1}{3}$ 时, $f\left( x\right)$ 取得极小值,无极大值,所以 $f{\left( x\right) }_{\min } = f\left( {-\dfrac{1}{3}}\right)  =  - \dfrac{4}{27}$
                
                所以在区间 $\left\lbrack  {-1,2}\right\rbrack$ 上, $f\left( x\right)$ 的最大值为 $f\left( {-1}\right)$ 或 $f\left( 2\right)$ ,而 $f\left( {-1}\right)  = 0,f\left( 2\right)  = 8 + 8 + 2 = {18}$，所以 $f\left( x\right)$ 在区间 $\left\lbrack  {-1,2}\right\rbrack$ 上的最大值为 18,最小值为 $- \dfrac{4}{27}$
                
                作出函数 $y = f\left( x\right)$ 与直线 $y = k$ 的图像,如图,
                
                \begin{center}
                \includegraphics[width=0.2\textwidth]{bodyxb/src/bo_d5rb33bef24c73bob2k0_6_140_1813_349_284_0.jpg}
                \end{center}
                
                
                由图知 $- \dfrac{4}{27} < k \leq  0$ .
            \end{enumerate}
        \end{jieda}
        \item 已知函数$f(x) = x^3+ax^2+b(a, b \in \mathbf{R})$
        \begin{enumerate}
            \item 讨论$y = f(x)$的单调性
            \item 若$b = c-a$（$c \in \mathbf{R}$，与$a$无关），当函数$y = f(x)$有三个不同的零点时，$a$的取值范围恰好是$(-\infty, -3), \left(1, \dfrac{3}{2}\right) \cup \left(\dfrac{3}{2}, +\infty\right)$，求$c$的值.
        \end{enumerate}
        \begin{jieda}
            \begin{enumerate}
                \item $f'(x) = 3x^2+2ax$，令$f'(x) = 0$，解得$x_1 = 0, x_2 = -\dfrac{2a}{3}$
                \begin{dingautolist}{172}
                    \item 若$a < 0$，则$f(x)$在$\left(-\infty, 0\right)$和$\left(-\dfrac{2a}{3}, +\infty\right)$严格增，在$\left(0, -\dfrac{2a}{3}\right)$严格减
                    \item 若$a = 0$，则$f(x)$在$\mathbf{R}$上严格增
                    \item 若$a > 0$，则$f(x)$在$\left(-\infty, -\dfrac{2a}{3}\right)$和$\left(0, +\infty\right)$严格增，在$\left(-\dfrac{2a}{3}, 0\right)$严格减
                \end{dingautolist}
                \item 
                根据三次函数的性质可知应存在两个极值，且符号相反才可使得函数具有三个不同的零点，即$f(0) \cdot f\left(-\dfrac{2a}{3}\right) <0$，由题意可知$f(0) \cdot f\left(-\dfrac{2a}{3}\right) = (c-a)\left(\dfrac{4}{27}a^3+c-a\right) < 0$的解集为$(-\infty, -3), \left(1, \dfrac{3}{2}\right) \cup \left(\dfrac{3}{2}, +\infty\right)$，因此$(c-a)\left(\dfrac{4}{27}a^3+c-a\right) = 0$的解集为$\left\{-3, 1, \dfrac{3}{2}\right\}$ \\
                代入$a = 1$可知$c = 1$或$\dfrac{23}{27}$，代入$a = -3$可知$c = 1$或$-3$，取交集有$c = 1$. \\
                经检验$c = 1$符合要求.
        \end{enumerate}
        \end{jieda}
        
        \item 已知函数 $f\left( x\right)  = \left\{  \begin{matrix} {{e}}^{x} + 1,x \leq  1 \\   - {x}^{2} + {4x} - 1,x > 1 \end{matrix}\right.$ ,若方程 $f\left( x\right)  - k\left| {x + 2}\right|  = 0$ 恰有三个不同实数根，则实数 $k$ 的取值范围是\blkda[5]{$\left( {\dfrac{2}{3},8 - 2\sqrt{13}}\right)  \cup  \left( {1,\dfrac{{e} + 1}{3}}\right\rbrack$}
        \begin{jieda}
            如图所示，作出函数 $y = f\left( x\right) ,y = k |x + 2|$ 的图象
            
            \begin{center}
            \includegraphics[width=0.3\textwidth]{bodyxb/src/bo_d5rb33bef24c73bob2k0_5_139_185_418_324_0.jpg}
            \end{center}
            
            方程 $f\left( x\right)  - k\left| {x + 2}\right|  = 0$ 恰有三个不同实数根,等价于上述两个函数图象有三个交点
            
            显然$k \leq 0$不满足要求，当$k > 0$时$y = k\left| {x + 2}\right|$ 与 $y = f\left( x\right)$ 在$x \leq -2$时必然有且仅有一个交点
            
            所以要满足题意需 $y = k\left| {x + 2}\right|$ 与 $y = f\left( x\right)$ 在$x \geq -2$时有两个交点

            \begin{dingautolist}{172}
                \item 若$y = k\left| {x + 2}\right|$与 $y = {{e}}^{x} + 1$这一支在$x \geq -2$时有两个交点
                
                先求 $y = {kx} + {2k}\left( {k > 0}\right)$ 与 $y = {{e}}^{x} + 1$ 相切时 $k$ 的值，
                
                设切点为 $\left( {{x}_{0},{{e}}^{{x}_{0}} + 1}\right) \left( {{x}_{0} >  - 2}\right)$ ,则该处切线为$y = e^{x_0}(x-x_0)+(e^{x_0}+1)$，代入$(-2, 0)$，因此有$ \left( {{x}_{0} + 1}\right) {{e}}^{{x}_{0}} - 1 = 0$，存在唯一解$x_0 = 0$，故切点为$(0, 2)$，此时$k = 1$
                
                当 $y = {kx} + {2k}\left( {k > 0}\right)$ 过点 $\left( {1,e + 1}\right)$ 时， $k = \dfrac{{e} + 1 - 0}{1 - \left( {-2}\right) } = \dfrac{{e} + 1}{3} > 1$ ，
                
                此时满足条件的 $k \in  \left( {1,\dfrac{{e} + 1}{3}}\right\rbrack$
    
                \item 再求 $y = {kx} + {2k}\left( {k > 0}\right)$ 与 $y =  - {x}^{2} + {4x} - 1\left( {x > 1}\right)$ 相切时 $k$ 的值
                
                $\left\{  {\begin{matrix} y = {kx} + {2k} \\  y =  - {x}^{2} + {4x} - 1 \end{matrix} \Longrightarrow  {x}^{2} + \left( {k - 4}\right) x + {2k} + 1 = 0,\Delta  = {\left( k - 4\right) }^{2} - 4\left( {{2k} + 1}\right)  = 0 \Longrightarrow  k = 8 \pm  }\right.$
                
                $2\sqrt{13}$
                
                易知切点横坐标为 $\dfrac{4 - k}{2}$ ，显然 $k = 8 - 2\sqrt{13} < 1$ 时， $\dfrac{4 - k}{2} = \sqrt{13} - 2 > 1$ ，符合要求，
                
                当 $y = {kx} + {2k}\left( {k > 0}\right)$ 过点 $\left( {1,2}\right)$ 时， $k = \dfrac{2 - 0}{1 - \left( {-2}\right) } = \dfrac{2}{3} < 8 - 2\sqrt{13}$ ,
                
                此时满足条件的 $k \in  \left( {\dfrac{2}{3},8 - 2\sqrt{13}}\right)$
            \end{dingautolist}
            
            综上: $k \in  \left( {\dfrac{2}{3},8 - 2\sqrt{13}}\right)  \cup  \left( {1,\dfrac{{e} + 1}{3}}\right\rbrack$ .
        \end{jieda}

        \item 已知函数 $f\left( x\right)  = \ln \left( {x + 1}\right)  + \dfrac{ax}{x + 1}\left( {a \in  \mathbf{R}}\right)$ ；现有下述两个结论：
        \begin{dingautolist}{172}
            \item 若 $y = f\left( x\right)$ 在区间$(-1,0)$内恰有一个零点,则 $a$ 的取值范围是$(-1,0)$
            \item 若 $a > 0$ ,则方程 ${f}^{\prime }\left( x\right)  - f\left( x\right)  = \dfrac{\sqrt{5} - 1}{2} - \ln \dfrac{\sqrt{5} + 1}{2}$ 的解为 $x = \dfrac{-1 \pm  \sqrt{5}}{2}$
        \end{dingautolist}
        则下列说法正确的是\dotfill{(\kh{B})}
        \xx{结论\ding{172}和\ding{173}均正确}{结论\ding{172}正确,结论\ding{173}错误}{结论\ding{172}错误,结论\ding{173}正确}{结论\ding{172}和\ding{173}均错误}
        \begin{jieda}
            \begin{dingautolist}{172}
                \item 直曲分离
                
                $\ln \left( {x + 1}\right)  + \dfrac{ax}{x + 1} = 0$即$(x+1)ln(x+1) = -ax$，数形结合可知当$a \in (-1, 0)$时满足要求，故正确.
                
                \begin{center}
                    \includegraphics[width=0.2\textwidth]{bodyxb/src/26jiejinganyimo16ti.png}
                \end{center}
                
                \ps 对$y = (x+1)ln(x+1)$求导得$y' = ln(x+1)+1$，故在$\left(-1, \dfrac{1}{e}-1\right)$严减，在$\left(\dfrac{1}{e}-1, +\infty\right)$严增. 
                
                当$x \rightarrow -1$时利用洛必达法则可知$\mathop{\lim }\limits_{{x \rightarrow -1}}(x+1)ln(x+1) = \mathop{\lim }\limits_{{x \rightarrow -1}}\dfrac{ln(x+1)}{\dfrac{1}{x+1}} = \mathop{\lim }\limits_{{x \rightarrow -1}}\dfrac{\dfrac{1}{x+1}}{-\dfrac{1}{(x+1)^2}} = \mathop{\lim }\limits_{{x \rightarrow -1}}-(x+1) = 0$

                易知$y = (x+1)ln(x+1)$与$y = -ax$在原点处相交，故当$y = -ax$的斜率大于$y = (x+1)ln(x+1)$处切线斜率，且为负即满足要求.

                \item 易知$\dfrac{-1-\sqrt{5}}{2}$不在定义域内，故错误.
                
                \ps 下证明$x = \dfrac{\sqrt{5}-1}{2}$为该方程唯一解

                $f(x) = ln(x+1)+a-\dfrac{a}{x+1}, f'(x) = \dfrac{1}{x+1}+\dfrac{a}{(x+1)^2}$

                注意到$\left\{\begin{array}{l}
                    \dfrac{\sqrt{5}+1}{2} = \dfrac{\sqrt{5}-1}{2}+1 \\
                    \dfrac{\sqrt{5}+1}{2} \times \dfrac{\sqrt{5}-1}{2} = 1
                \end{array} \right.$

                因此，当$x = \dfrac{\sqrt{5}-1}{2}$时有$x = \dfrac{1}{x+1}$，故$LHS = \dfrac{1-ax}{x+1}+\dfrac{a}{(x+1)^2}-ln(x+1) = (1-ax)x+ax^2 - ln(x+1) = x - ln(x+1) = RHS$，故$x = \dfrac{\sqrt{5}-1}{2}$为原方程的解

                当$a > 0$时，$f'(x) > 0$，故$f(x)$严增，易知$f'(x)$严减，则$g(x) = f'(x) - f(x)$严减，故$x = \dfrac{\sqrt{5}-1}{2}$为原方程的唯一解
            \end{dingautolist}
        \end{jieda}
        
        \item 已知函数 $f\left( x\right)  = {2a}\ln x + \dfrac{1}{2}{x}^{2} - \left( {a + 2}\right) x$ ，其中 $a$ 为常数.
        \begin{enumerate}
            \item 当 $a > 0$ 时，试讨论 $f\left( x\right)$ 的单调性
            \item 若函数 $f\left( x\right)$ 有两个不相等的零点 ${x}_{1},{x}_{2}$
            \begin{dingautolist}{172}
                \item 求 $a$ 的取值范围
                \item 证明: ${x}_{1} + {x}_{2} > 4$
            \end{dingautolist}
        \end{enumerate}
        \begin{jieda}
            \begin{enumerate}
                \item 由题设 $f'\left( x\right)  = \dfrac{2a}{x} + x - \left( {a + 2}\right)  = \dfrac{{x}^{2} - \left( {a + 2}\right) x + {2a}}{x} = \dfrac{\left( {x - 2}\right) \left( {x - a}\right) }{x}$ ,且 $x \in  \left( {0, + \infty }\right)$
                \begin{dingautolist}{172}
                    \item 当 $0 < a < 2$ 时，在 $\left( {0,a}\right)$ 上 ${f}^{\prime }\left( x\right)  > 0$ ，在 $\left( {a,2}\right)$ 上 ${f}^{\prime }\left( x\right)  < 0$ ，在 $\left( {2, + \infty }\right)$ 上 ${f}^{\prime }\left( x\right)  > 0$，所以$f(x)$在 $\left( {0,a}\right)$ 、 $\left( {2, + \infty }\right)$ 上 $f\left( x\right)$ 单调递增，在 $\left( {a,2}\right)$ 上 $f\left( x\right)$ 单调递减
                    
                    \item 当 $a = 2$ 时，在 $\left( {0, + \infty }\right)$ 上 ${f}^{\prime }\left( x\right)  \geq  0$ 恒成立，故 $f\left( x\right)$ 在 $\left( {0, + \infty }\right)$ 上单调递增；
                    
                    \item 当 $a > 2$ 时，在 $\left( {0,2}\right)$ 上 ${f}^{\prime }\left( x\right)  > 0$ ，在 $\left( {2,a}\right)$ 上 ${f}^{\prime }\left( x\right)  < 0$ ，在 $\left( {a, + \infty }\right)$ 上 ${f}^{\prime }\left( x\right)  > 0$，所以$f(x)$在 $\left( {0,2}\right)$ 、 $\left( {a, + \infty }\right)$ 上 $f\left( x\right)$ 单调递增，在 $\left( {2,a}\right)$ 上 $f\left( x\right)$ 单调递减.
                \end{dingautolist}
                \item \begin{dingautolist}{172}
                    \item 由 $f\left( 2\right)  = {2a}\ln 2 + \dfrac{1}{2} \times  {2}^{2} - \left( {a + 2}\right)  \times  2 = 2\left( {a\ln 2 - a - 1}\right)  < 0$
                    
                    若 $a > 0$ 时，当$x \rightarrow +\infty$时，$f(x) \rightarrow +\infty$，结合(1)中单调性可知，$f(x)$必然在$(2, +\infty)$存在唯一零点
                    
                    $f\left( a\right)  = {2a}\ln a + \dfrac{1}{2}{a}^{2} - \left( {a + 2}\right) a = \dfrac{a}{2}\left( {4\ln a - a - 4}\right)$
                    
                    令 $y = 4\ln a - a - 4$ 且 $a > 0$ ，则 ${y}^{\prime } = \dfrac{4}{a} - 1 = \dfrac{4 - a}{a}$ ，
                    
                    所以 $0 < a < 4$ 时 ${y}^{\prime } > 0,a > 4$ 时 ${y}^{\prime } < 0$，故 $y = 4\ln a - a - 4$ 在 $\left( {0,4}\right)$ 上递增,在 $\left( {4, + \infty }\right)$ 上递减,则 ${y}_{\max } = 8\ln 2 - 8 < 0$，所以 $f\left( a\right)  < 0$，结合 (1) 中 $f\left( x\right)$ 的单调性可知 $a > 0$ 不可能出现两个不相等的零点,
                    
                    又 $a = 0$ 时， $f\left( x\right)  = \dfrac{1}{2}{x}^{2} - {2x}$ 在 $\left( {0, + \infty }\right)$ 上只有一个零点，不满足,
                    
                    所以必有 $a < 0$ ，此时，在 $\left( {0,2}\right)$ 上 ${f}^{\prime }\left( x\right)  < 0$ ，在 $\left( {2, + \infty }\right)$ 上 ${f}^{\prime }\left( x\right)  > 0$ ，
                    
                    故在 $\left( {0,2}\right)$ 上 $f\left( x\right)$ 单调递减，在 $\left( {2, + \infty }\right)$ 上 $f\left( x\right)$ 单调递增，则 $f{\left( x\right) }_{\min } = f\left( 2\right)  = 2\left( {a\ln 2 - a - 1}\right)$ ，
                    
                    又 $x$ 趋向于 0 或正无穷时, $f\left( x\right)$ 趋向正无穷,只需 $a\left( {\ln 2 - 1}\right)  - 1 < 0$ 成立，即所求范围为 $\dfrac{1}{\ln 2 - 1} < a < 0$

                    \item 由\ding{172},在 $\dfrac{1}{\ln 2 - 1} < a < 0$ 时, $f\left( x\right)$ 存在两个不相等的零点 ${x}_{1},{x}_{2}$ ,
                    
                    不妨令 $0 < {x}_{1} < 2 < {x}_{2}$ ,要证 ${x}_{1} + {x}_{2} > 4$ ,即证 ${x}_{2} > 4 - {x}_{1}$ ,而 $4 - {x}_{1} \in  \left( {2,4}\right)$ ,
                    
                    由 \ding{172} 知: 在 $\left( {2, + \infty }\right)$ 上 $f\left( x\right)$ 单调递增,只需证 $f\left( {4 - {x}_{1}}\right)  < f\left( {x}_{2}\right)  = f\left( {x}_{1}\right)  = 0$ ,
                    
                    由 ${2a}\ln {x}_{1} + \dfrac{1}{2}{x}_{1}^{2} - \left( {a + 2}\right) {x}_{1} = 0$ ，则 ${x}_{1}^{2} - 4{x}_{1} = {2a}\left( {{x}_{1} - 2\ln {x}_{1}}\right)$
                    
                    令 $h\left( x\right)  = f\left( {4 - x}\right)  - f\left( x\right)$ ，且 $0 < x < 2$ ，
                    
                    则 $h\left( x\right)  = {2a}\ln \left( {4 - x}\right)  + \dfrac{1}{2}{\left( 4 - x\right) }^{2} - \left( {a + 2}\right) \left( {4 - x}\right)  - {2a}\ln x - \dfrac{1}{2}{x}^{2} + \left( {a + 2}\right) x$
                    
                    $= {2a}\ln \left( {4 - x}\right)  - {2a}\ln x - {4a} + {2ax}$ ,
                    
                    所以，在 $\left( {0,2}\right)$ 上 ${h}^{\prime }\left( x\right)  =  - {2a} \cdot  \dfrac{{\left( x - 2\right) }^{2}}{x\left( {4 - x}\right) } > 0$ ，即 $h\left( x\right)$ 在 $\left( {0,2}\right)$ 上递增，
                    
                    所以 $h\left( x\right)  < h\left( 2\right)  = {2a}\ln 2 - {2a}\ln 2 - {4a} + {4a} = 0$ ,即 $f\left( {4 - {x}_{1}}\right)  < f\left( {x}_{1}\right)  = f\left( {x}_{2}\right)$ 成立,
                    
                    所以 ${x}_{1} + {x}_{2} > 4$ ,得证.
                \end{dingautolist}
            \end{enumerate}
        \end{jieda}
    \end{enumerate}
    
    \item 极值点偏移
    \begin{enumerate}
        \item 已知函数 $f\left( x\right)  = \left( {x - 2}\right) {{e}}^{x}$ ,若 $f\left( {x}_{1}\right)  = f\left( {x}_{2}\right)$ ,且 ${x}_{1} \neq  {x}_{2}$ , ${x}_{1} \cdot  {x}_{2} > 0$ ,则\dotfill{(\kh{D})}
        \xx{${x}_{1} > \dfrac{1}{2}$}{${x}_{2} < \dfrac{3}{2}$}{${x}_{1}{x}_{2} > 1$}{${x}_{1} + {x}_{2} < 2$}
        \begin{jieda}
            $f\left( x\right)  = \left( {x - 2}\right) {{e}}^{x}$ ,则 ${f}^{\prime }\left( x\right)  = \left( {x - 1}\right) {{e}}^{x}$ ,令 ${f}^{\prime }\left( x\right)  = 0 \Longrightarrow  x = 1$ ,
            
            当 $x \in  \left( {-\infty ,1}\right)$ 时 ${f}^{\prime }\left( x\right)  < 0,f\left( x\right)$ 单调递减,当 $x \in  \left( {1, + \infty }\right)$ 时 ${f}^{\prime }\left( x\right)  > 0,f\left( x\right)$ 单调递增,
            
            在 $x \in  \left( {-\infty ,2}\right)$ 上 $f\left( x\right)  < 0$ ，且 $f\left( 2\right)  = 0,f\left( 0\right)  =  - 2,f{\left( x\right) }_{\min } = f\left( 1\right)  =  - {e}$ ，即 $f\left( x\right)  \geq   - {e}$
            
            综上, $f\left( x\right)$ 的图象如下: 结合 $f\left( {x}_{1}\right)  = f\left( {x}_{2}\right)  = k,{x}_{1} \cdot  {x}_{2} > 0$ ,令 ${x}_{1} < {x}_{2}$ ,
            
            \begin{center}
            \includegraphics[width=0.3\textwidth]{bodyxb/src/bo_d5rb33bef24c73bob2k0_21_138_864_409_382_0.jpg}
            \end{center}
            
            如上图,若 $f\left( {x}_{0}\right)  = \left( {{x}_{0} - 2}\right) {{e}}^{{x}_{0}} =  - 2$ 且 ${x}_{0} \neq  0$ ,则 $0 < {x}_{1} < 1 < {x}_{2} < {x}_{0} < 2$ ,则 ${x}_{1} > \dfrac{1}{2}$ 不一定成立, ${A}$ 错误;
            
            又 $f\left( \dfrac{3}{2}\right)  = \left( {\dfrac{3}{2} - 2}\right) {{e}}^{\frac{3}{2}} =  - \dfrac{\sqrt{{{e}}^{3}}}{2} \in  \left( {-{e}, - 2}\right)$ ,故 $\dfrac{3}{2} < {x}_{0}$ ,则 ${x}_{2} < \dfrac{3}{2}$ 不一定成立, ${B}$ 错误;
            
            令 $g\left( x\right)  = f\left( {1 + x}\right)  - f\left( {1 - x}\right)  = \left( {x - 1}\right) {{e}}^{x + 1} + \left( {x + 1}\right) {{e}}^{1 - x}$ ,
            
            则 ${g}^{\prime }\left( x\right)  = x{{e}}^{x + 1} - x{{e}}^{1 - x} = x\left( {{{e}}^{x + 1} - {{e}}^{1 - x}}\right)$ ,
            
            当 $x \geq  0$ 时， $x + 1 \geq  1 - x$ ，得 ${{e}}^{x + 1} \geq  {{e}}^{1 - x}$ ，则 ${g}^{\prime }\left( x\right)  \geq  0$ ；
            
            当 $x < 0$ 时， $x + 1 < 1 - x$ ，得 ${{e}}^{x + 1} < {{e}}^{1 - x}$ ，则 ${g}^{\prime }\left( x\right)  > 0$ ，
            
            所以函数 $g\left( x\right)$ 在 ${R}$ 上单调递增,且 $g\left( 0\right)  = 0$ ,
            
            所以 $f\left( {1 + x}\right)  > f\left( {1 - x}\right)$ 在 ${R}$ 上恒成立,得 $f\left( {x}_{2}\right)  = f\left\lbrack  {1 + \left( {{x}_{2} - 1}\right) }\right\rbrack   > f\left\lbrack  {1 - \left( {{x}_{2} - 1}\right) }\right\rbrack$ ,
            
            即 $f\left( {x}_{2}\right)  > f\left( {2 - {x}_{2}}\right)$ ,又 $f\left( {x}_{1}\right)  = f\left( {x}_{2}\right)$ ,所以 $f\left( {x}_{1}\right)  > f\left( {2 - {x}_{2}}\right)$ ,
            
            由 $2 - {x}_{2} < 1$ ，且函数 $f\left( x\right)$ 在 $x \in  \left( {-\infty ,1}\right)$ 单调递减，得 ${x}_{1} < 2 - {x}_{2}$ ，即 ${x}_{1} + {x}_{2} < 2$ ， $D$ 正确.
            
            又 $0 < {x}_{1} < 1 < {x}_{2} < {x}_{0} < 2$ ,则 $2 > {x}_{1} + {x}_{2} > 2\sqrt{{x}_{1}{x}_{2}}$ ,即 $1 > \sqrt{{x}_{1}{x}_{2}}$ ,故 $0 < {x}_{1}{x}_{2} < 1,{C}$ 错误.

            \ps 数形结合，判断出极值点右偏可以快速选出答案.
        \end{jieda}
        \item 已知常数$a > 0$，函数$f(x) = \dfrac{1}{2}x^2-ax-2a^2lnx$，若$x_1, x_2$是$f(x)$的零点，且$x_1 \neq x_2$，证明：$x_1+x_2>4a$.
        \begin{jieda}
            $f'(x) = x-a-\dfrac{2a^2}{x}$，易知$f'(2a) = 0$，因此$f(x)$在$(0, 2a)$严减，在$(2a, +\infty)$严增，故不妨设$0 < x_1 < 2a < x_2$

            设$F(x) = f(x) - f(4a-x)$，则$F'(x) = f'(x) + f'(4a-x) = \left(x-a-\dfrac{2a^2}{x}\right)+\left(4a-x-a-\dfrac{2a^2}{4a-x}\right) = 2a - 2a^2\left(\dfrac{1}{x} + \dfrac{1}{4a-x}\right)$，由平均值不等式可知$F'(x) \leq 0 (x \in (0, 4a))$，故$F(x_1) > F(2a) = 0$

            因此$f(x_2) = f(x_1) > f(4a-x_1)$，根据$f(x)$单调性可知$x_2 > 4a-x_1$，即$x_1+x_2 > 4a$
        \end{jieda}

        \item 已知函数$f(x) = (1+lnx)e^{-ln(ax)}$，若$f(x) = 1$有两个根$x_1, x_2$，求实数$a$的取值范围，并证明\ding{172}$x_1x_2>1$\ding{173}$x_1^2+x_2^2>2$
        \begin{jieda}
            $(1+lnx)e^{-ln(ax)} = 1 \Longleftrightarrow 1+lnx = ax \Longleftrightarrow a = \dfrac{1+lnx}{x}$

            设$g(x) = \dfrac{1+lnx}{x}$，则$g'(x) = \dfrac{-lnx}{x^2}$，故$g(x)$在$(0, 1)$严增，在$(1, +\infty)$严减，易知$a \in (0, 1)$满足要求，故不妨设$0 < x_1 < 1 < x_2$

            \begin{dingautolist}{172}
                \item 设$G(x) = g(x) - g\left(\dfrac{1}{x}\right)$，则$G'(x) = g'(x) + \dfrac{g'\left(\dfrac{1}{x}\right)}{x^2} = \dfrac{-lnx}{x^2} - ln\left(\dfrac{1}{x}\right) = \dfrac{x^2-1}{x^2}lnx \geq 0$

                故$g(x_2) = g(x_1) < g\left(\dfrac{1}{x_1}\right)$，根据$g(x)$单调性可知$x_2 > \dfrac{1}{x_1}$，故$x_1x_2 > 1$

                \item 
                \begin{enumerate}[i]
                    \item 方法1

                    利用$x_1x_2 > 1$，结合平均值不等式$x_1^2+x_2^2 > 2x_1x_2 > 2$

                    \item 方法2

                    构造函数$H(x) = g(x) - g(2-x)$证明$x_1+x_2 > 2$，结合平均值不等式可知$\left\{\begin{array}{l}
                         x_1^2+1 > 2x_1  \\
                         x_2^2+1 > 2x_2
                    \end{array}\right.$，故$x_1^2+x_2^2 > 2(x_1+x_2) - 2 > 2$
                \end{enumerate}
                
            \end{dingautolist}

            
        \end{jieda}
    \end{enumerate}
    \item 函数单调性、极值与最值的综合应用
    \begin{enumerate}
        \item 已知函数$f(x) = \dfrac{1}{3}x^3 - \dfrac{a+1}{2}x^2+ax$
        \begin{enumerate}
            \item 讨论函数$y = f(x)$的单调性
            \item 对于任意$x_1, x_2 \in [0, 2]$，都有$|f(x_1) - f(x_2)| \leq \dfrac{2}{3}$，求$a$的取值范围
        \end{enumerate}
        \begin{jieda}
            \begin{enumerate}
                \item $f'(x) = x^2 - (a+1)x + a$，$f'(a) = f'(1) = 0$
                \begin{dingautolist}{172}
                    \item $a < 1$，此时$f(x)$在$(-\infty, a)$和$(1, +\infty)$内严格增，在$(a, 1)$严格减.
                    \item $a > 1$，此时$f(x)$在$(-\infty, 1)$和$(a, +\infty)$内严格增，在$(1, a)$严格减.
                    \item $a = 1$，此时$f(x)$在$\mathbf{R}$上严格增
                \end{dingautolist}
                \item 对于任意$x_1, x_2 \in [0, 2]$，都有$|f(x_1) - f(x_2)| \leq \dfrac{2}{3}$即$f(x)$在$[0, 2]$内$f_{max} - f_{min} \leq \dfrac{2}{3}$  \\
                方法1：分类讨论 
                \begin{dingautolist}{172}
                    \item $a \in (-\infty, 0]$ \\
                    此时$f(x)$在$[0, 1]$内严减，在$[1, 2]$内严增，因此最大值为$f(2) = \dfrac{2}{3}$，最小值为$f(1) = \dfrac{a}{2}-\dfrac{1}{6}$，不满足要求.
                    \item $a \in (0, 1)$ \\
                    此时$f(x)$在$[0, a]$内严增，在$[a, 1]$内严减，在$[1, 2]$内严增，因此最大值为$f(2) = \dfrac{2}{3}$，当$a \in \left[\dfrac{1}{3}, 1\right)$时最小值为$f(0) = 0$，满足要求，否则不满足.
                    \item $a = 1$ \\
                    此时$f(x)$在$[0, 2]$上严增，因此最大值为$f(2) = \dfrac{2}{3}$，最小值为$f(0) = 0$，满足要求.
                    \item $a \in (1, 2)$ \\
                    此时$f(x)$在$[0, 1]$严增，在$[1, a]$内严减，在$[a, 2]$内严增，最大值为$f(2) = \dfrac{2}{3}$，当$a \in \left(1, \dfrac{5}{3}\right]$时，最小值为$f(0)$，满足要求，否则不满足.
                    \item $a \in [2, +\infty)$ \\
                    此时$f(x)$在$[0, 1]$严增，在$[1, 2]$内严减，最大值为$f(1) = \dfrac{a}{2}-\dfrac{1}{6}$，最小值为$f(0)=0$，不满足要求.
                \end{dingautolist}
                方法2：利用必要条件 \\
                $f(0) = 0, f(1) = \dfrac{1}{3} - \dfrac{a+1}{2} + a = \dfrac{a}{2}-\dfrac{1}{6}, f(2) = \dfrac{8}{3} - 2(a+1)+2a = \dfrac{2}{3}, f(a) = \dfrac{a^3}{3}-\dfrac{(a+1)a^2}{2} + a^2 = -\dfrac{a^3}{6} + \dfrac{a^2}{2}$ \\
                若满足要求，则必有$\dfrac{2}{3}$为最大值，$0$为最小值. 因此首先可以确定$f(1)$必然要满足$f(1) \in \left[0, \dfrac{2}{3}\right]$，故$\dfrac{a}{2} - \dfrac{1}{6} \in \left[0, \dfrac{2}{3}\right]$，解得$a \in \left[\dfrac{1}{3}, \dfrac{5}{3}\right]$，故有$a \in [0, 2]$，$f(a)$也应满足$f(a) \in \left[0, \dfrac{2}{3}\right]$，故$-\dfrac{a^3}{6} + \dfrac{a^2}{2} \in \left[0, \dfrac{2}{3}\right]$，即$-a^3+3a^2 \in [0, 4]$，设$h(a) = -a^3+3a^2$，$h'(a) = -3a^2+6a \geq 0$在$\left[\dfrac{1}{3}, \dfrac{5}{3}\right]$内恒成立，因此$h(a)$在$\left[\dfrac{1}{3}, \dfrac{5}{3}\right]$内严增，可以判断$0 < h\left(\dfrac{1}{3}\right) \leq h\left(\dfrac{5}{3}\right) < 4$，满足要求.
                
                综上，当$a \in \left[\dfrac{1}{3}, \dfrac{5}{3}\right]$时满足要求
            \end{enumerate}
        \end{jieda}
        
        \item 关于函数$f(x) = (x^2+ax-1)e^x$有以下三个结论
        \begin{dingautolist}{172}
            \item 函数恒有两个零点，且两个零点之积为$-1$
            \item 函数的极值点不可能是$-1$
            \item 函数必有最小值
        \end{dingautolist}
        其中正确的结论是\blkda{\ding{172}\ding{173}\ding{174}}
        \begin{jieda}
            \begin{dingautolist}{172}
                \item 易知$f(x)$的零点即$y = x^2+ax-1$的零点，由于$\Delta = a^2+4 > 0$，所以$y = x^2+ax-1$有两个不同的零点$x_1, x_2$，根据韦达定理可知$x_1x_2 = 1$，因此\ding{172}正确
                \item $f'(x) = e^x(2x+a + x^2+ax-1) = e^x(x^2+(a+2)x+(a-1))$，令$f'(x) = 0$，则$x^2+(a+2)x+(a-1) = 0$，代入$x = -1$等式必不成立，因此\ding{173}正确
                \item $f'(x) = e^x(2x+a + x^2+ax-1) = e^x(x^2+(a+2)x+(a-1))$，令$f'(x) = 0$，则$x^2+(a+2)x+(a-1) = 0$，因为$\Delta = (a+2)^2-4(a-1) = a^2+4a+4-4a+4 = a^2+8 > 0$，因此$f'(x) = 0$必有两个相异实根$x_3, x_4(x_3 < x_4)$ \\
                判断$f'(x)$的符号有：$f'(x)$在$(-\infty, x_3) \cup (x_4, +\infty)$内为正，在$(x_3, x_4)$内为负，因此$f(x)$在$(-\infty, x_3)$内严格增，在$(x_3, x_4)$内严格减，在$(x_4, +\infty)$内严格增，故存在极小值$x_4$. \\
                通过作差法可以确定$x_1, x_2, x_3, x_4$的大小为$x_3 < x_1 < x_4 < x_2$后可以确定$f(x_4)$是最小值点；但事实上，通过数形结合即可以确定$f(x_4)$为最小值点，故\ding{174}正确.
            \end{dingautolist}
        \end{jieda}
        
        \item 若 $x = \dfrac{7}{2}$ 是函数 $f\left( x\right)  = a{x}^{2} + \ln \left( {3x}\right)$ 的一个极值点，则当 $x \in  \left[\dfrac{1}{e}, e\right]$ 时， $f\left( x\right)$ 的最小值为\blkda[3]{$- 2{{e}}^{2} + \ln 3 + 1$}
        \begin{jieda}
            由 $f\left( x\right)  = a{x}^{2} + \ln \left( {3x}\right)$ 求导得, ${f}^{\prime }\left( x\right)  = {2ax} + \dfrac{1}{x}$ ,
            
            依题意， ${f}^{\prime }\left( \dfrac{1}{2}\right)  = {2a} \times  \dfrac{1}{2} + \dfrac{1}{\dfrac{1}{2}} = a + 2 = 0$ ，解得 $a =  - 2$ ，
            
            此时， $f\left( x\right)  =  - 2{x}^{2} + \ln \left( {3x}\right)$ ，则 ${f}^{\prime }\left( x\right)  =  - {4x} + \dfrac{1}{x} = \dfrac{\left( {1 + {2x}}\right) \left( {1 - {2x}}\right) }{x}$ ，因 $x > 0$ ，
            
            故当 $0 < x < \dfrac{1}{2}$ 时， ${f}^{\prime }\left( x\right)  > 0$ ，当 $x > \dfrac{1}{2}$ 时， ${f}^{\prime }\left( x\right)  < 0$ ，
            
            即函数 $f\left( x\right)$ 在 $\left( {0,\dfrac{1}{2}}\right)$ 上递增,在 $\left( {\dfrac{1}{2}, + \infty }\right)$ 上递减,即 $\dfrac{1}{2}$ 是 $f\left( x\right)$ 的极大值点.
            
            又因 $x \in  \left[\dfrac{1}{e}, e\right]$ ，故得函数 $f\left( x\right)$ 在 $\left( {\dfrac{1}{{e}},\dfrac{1}{2}}\right)$ 上递增，在 $\left( {\dfrac{1}{2},{e}}\right)$ 上递减.
            
            因 $f\left( \dfrac{1}{{e}}\right)  =  - \dfrac{2}{{{e}}^{2}} + \ln 3 - 1,f\left( {e}\right)  =  - 2{{e}}^{2} + \ln 3 + 1$ ,显然 $f\left( \dfrac{1}{{e}}\right)  > f\left( {e}\right)$ ,
            
            故 $f\left( x\right)$ 的最小值为 $f\left( {e}\right)  =  - 2{{e}}^{2} + \ln 3 + 1$ .
        \end{jieda}
        
        \item 若函数 $f\left( x\right)  = x{{e}}^{x} - {ax}$ ，则根据下列说法正确的是\blkda{\ding{172}\ding{173}\ding{174}}
        \begin{dingautolist}{172}
            \item 当 $a \in  \left( {-\infty , - {{e}}^{-2}}\right\rbrack$ 时， $f\left( x\right)$ 在 $x \in  {R}$ 上单调递增
            \item 当 $a \in  \left( {-{{e}}^{-2},0}\right)$ 时， $f\left( x\right)$ 有两个极值点
            \item 当 $a \in  \left( {-\infty , - {{e}}^{-2}}\right\rbrack$ 时， $f\left( x\right)$ 没有最小值.
        \end{dingautolist}
        \begin{jieda}
            ${f}^{\prime }\left( x\right)  = \left( {1 + x}\right) {{e}}^{x} - a$ ,
            
            设 $g\left( x\right)  = \left( {1 + x}\right) {{e}}^{x} - a,{g}^{\prime }\left( x\right)  = \left( {2 + x}\right) {{e}}^{x}$ ,
            
            当 $x <  - 2$ 时， ${g}^{\prime }\left( x\right)  < 0,g\left( x\right)$ 单调递减，
            
            $x >  - 2$ 时, ${g}^{\prime }\left( x\right)  > 0,g\left( x\right)$ 单调递增,
            
            所以 $g{\left( x\right) }_{\min } = g\left( {-2}\right)  =  - {{e}}^{-2} - a =  - \dfrac{1}{{{e}}^{2}} - a$ ,
            
            当 $a \leq   - {{e}}^{-2}$ 时, $g\left( x\right)  \geq  0$ ,即 ${f}^{\prime }\left( x\right)  \geq  0$ ,
            
            所以函数 $f\left( x\right)$ 在 $x \in  R$ 上单调递增，则没有最小值， \ding{172}  \ding{174} 正确；
            
            当 $- {{e}}^{-2} < a < 0$ 时， $g\left( x\right)  = \left( {1 + x}\right) {{e}}^{x} - a = 0$ ，即 $\left( {1 + x}\right) {{e}}^{x} = a$ ，
            
            设 $h\left( x\right)  = \left( {1 + x}\right) {{e}}^{x}$ ，由上面的研究可知，
            
            当 $x <  - 2$ 时， ${h}^{\prime }\left( x\right)  < 0,h\left( x\right)$ 单调递减，
            
            $x >  - 2$ 时, ${h}^{\prime }\left( x\right)  > 0,h\left( x\right)$ 单调递增,
            
            所以 $h{\left( x\right) }_{\min } = h\left( {-2}\right)  =  - {{e}}^{-2} =  - \dfrac{1}{{{e}}^{2}}$ ,
            
            且当 $x <  - 2$ 时, $h\left( x\right)  < 0$ ,且 $h\left( x\right)  \in  \left( {-\dfrac{1}{{{e}}^{2}},0}\right) ,x >  - 2$ 时, $h\left( {-1}\right)  = 0$ ,
            
            所以此时方程 $\left( {1 + x}\right) {{e}}^{x} = a$ 有两个解,即 $g\left( x\right)  = \left( {1 + x}\right) {{e}}^{x} - a$ 有两个零点,
            
            所以 $f\left( x\right)$ 有两个极值点, \ding{173}  正确，
            
            所以正确答案是 \ding{172}  \ding{173}  \ding{174} .
        \end{jieda}
        
        \item 已知函数 $f\left( x\right)  = {x}^{2} - \left( {{2a} + 1}\right) x + a\ln x,a \in  R$
        \begin{enumerate}
            \item 若 $a = 0$ ，求曲线 $y = f\left( x\right)$ 在点 $P\left( {2,f\left( 2\right) }\right)$ 处的切线方程
            
            \item 若 $f\left( x\right)$ 在 $x = 1$ 处取得极值，求 $f\left( x\right)$ 的极值.
            
            \item 若 $f\left( x\right)$ 在 $\left\lbrack  {1,e}\right\rbrack$ 上的最小值为 $- {2a}$ ，求 $a$ 的取值范围.
        \end{enumerate}

        \begin{jieda}
            \begin{enumerate}
                \item 若 $a = 0$ ，则 $f\left( x\right)  = {x}^{2} - x$ ，则 ${f}^{\prime }\left( x\right)  = {2x} - 1$ ，
                
                故 $f\left( 2\right)  = 2,{f}^{\prime }\left( 2\right)  = 3$ ,
                
                故曲线 $y = f\left( x\right)$ 在点 $P\left( {2,f\left( 2\right) }\right)$ 处的切线方程为 $y - 2 = 3\left( {x - 2}\right)$ ,即 ${3x} - y - 4 = 0$ ;
                
                \item  $f\left( x\right)  = {x}^{2} - \left( {{2a} + 1}\right) x + a\ln x,a \in  R$ 定义域为 $\left( {0, + \infty }\right)$ ,
                
                则 ${f}^{\prime }\left( x\right)  = {2x} - \left( {{2a} + 1}\right)  + \dfrac{a}{x}$ ,
                
                由于 $f\left( x\right)$ 在 $x = 1$ 处取得极值,故 ${f}^{\prime }\left( 1\right)  = 2 - \left( {{2a} + 1}\right)  + a = 0,\therefore a = 1$ ,
                
                则 ${f}^{\prime }\left( x\right)  = {2x} - 3 + \dfrac{1}{x} = \dfrac{2{x}^{2} - {3x} + 1}{x} = \dfrac{\left( {{2x} - 1}\right) \left( {x - 1}\right) }{x}$ ,
                
                令 ${f}^{\prime }\left( x\right)  > 0$ ,则 $0 < x < \dfrac{1}{2}$ 或 $x > 1$ ,函数 $f\left( x\right)$ 在 $\left( {0,\dfrac{1}{2}}\right) ,\left( {1, + \infty }\right)$ 上均单调递增,
                
                令 ${f}^{\prime }\left( x\right)  < 0$ ，则 $\dfrac{1}{2} < x < 1$ ，函数 $f\left( x\right)$ 在 $\left( {\dfrac{1}{2},1}\right)$ 上单调递减，
                
                故当 $x = \dfrac{1}{2}$ 时， $f\left( x\right)$ 取到极大值 $f\left( \dfrac{1}{2}\right)  = \dfrac{1}{4} - \dfrac{3}{2} + \ln \dfrac{1}{2} =  - \dfrac{5}{4} - \ln 2$ ，
                
                当 $x = 1$ 时， $f\left( x\right)$ 取到极小值 $f\left( 1\right)  = 1 - 3 =  - 2$ ；
                
                \item  由于 ${f}^{\prime }\left( x\right)  = {2x} - \left( {{2a} + 1}\right)  + \dfrac{a}{x} = \dfrac{\left( {{2x} - 1}\right) \left( {x - a}\right) }{x},x \in  \left\lbrack  {1,{e}}\right\rbrack$

                \begin{dingautolist}{172}
                    \item 当 $a \leq  1$ 时， ${f}^{\prime }\left( x\right)  \geq  0$ ，仅在 $a = 1$ ， $x = 1$ 时等号取得， $f\left( x\right)$ 在 $\left\lbrack  {1,{e}}\right\rbrack$ 上单调递增，则 $f{\left( x\right) }_{\min } = f\left( 1\right)  =  - {2a}$ ,符合题意
                    
                    \item 当 $1 < a < {e}$ 时,则 $1 < x < a$ 时, ${f}^{\prime }\left( x\right)  < 0,f\left( x\right)$ 在 $\left\lbrack  {1,a}\right\rbrack$ 上单调递减, $a < x < {e}$ 时, ${f}^{\prime }\left( x\right)  > 0,f\left( x\right)$ 在 $\left\lbrack  {a,{e}}\right\rbrack$ 上单调递增,故 $f{\left( x\right) }_{\min } = f\left( a\right)  < f\left( 1\right)  =  - {2a}$ ,不符合题意
                    
                    \item 当 $a \geq  {e}$ 时， ${f}^{\prime }\left( x\right)  < 0,f\left( x\right)$ 在 $\left\lbrack  {1,{e}}\right\rbrack$ 上单调递减，故 $f{\left( x\right) }_{\min } = f\left( {e}\right)  < f\left( 1\right)  =  - {2a}$ ,不符合题意
                \end{dingautolist}
                
                综上,可知 $a$ 的取值范围为 $( - \infty ,1\rbrack$ .
            \end{enumerate}
        \end{jieda}
        
        \item 已知函数 $f\left( x\right)  =  - {x}^{3} + 3{x}^{2} + a\left( {x > 0}\right) ,g\left( x\right)  = x\ln x +  a{x}^{2} - {2x}$ .
        \begin{enumerate}
            \item 求函数 $f\left( x\right)$ 的极值
            
            \item 若 $f\left( x\right) ,g\left( x\right)$ 的导数分别为 ${f}^{\prime }\left( x\right) ,{g}^{\prime }\left( x\right)$ ，且 $\left\{  {x|{f}^{\prime }\left( x\right)  < 0}\right\}   \subseteq  \left\{  {x|{g}^{\prime }\left( x\right)  < 0}\right\}$ ，求 $a$ 的取值范围
            \item 用 $\min \{ m,n\}$ 表示 $m,n$ 中的最小值,设 $h\left( x\right)  = \min \{ f\left( x\right) ,g\left( x\right) \}$ ,若 $\left| a\right|  > 1$ ,判断函数 $h\left( x\right)$ 的零点个数.
        \end{enumerate}
        \begin{jieda}
            \begin{enumerate}
                \item ${f}^{\prime }\left( x\right)  =  - 3{x}^{2} + {6x}$ ，令 ${f}^{\prime }\left( x\right)  = 0$ 得， $x = 2$ 或 $x = 0$ (不合题意舍去)
                
                当 $x \in  \left( {0,2}\right)$ 时， ${f}^{\prime }\left( x\right)  > 0$ ， $f\left( x\right)$ 在 $\left( {0,2}\right)$ 上单调递增，
                
                当 $x \in  \left( {2, + \infty }\right)$ 时， ${f}^{\prime }\left( x\right)  < 0$ ， $f\left( x\right)$ 在 $\left( {2, + \infty }\right)$ 上单调递减，
                
                所以 $f\left( x\right)$ 的极大值为 $f\left( 2\right)  = 4 + a$ ,无极小值.
                
                \item 由 (1) 得, ${f}^{\prime }\left( x\right)  < 0$ 时, $x \in  \left( {2, + \infty }\right)$
                
                对 $g\left( x\right)$ 求导得, ${g}^{\prime }\left( x\right)  = \ln x + {2ax} - 1$ ,
                
                $\left\{  {x|{f}^{\prime }\left( x\right)  < 0}\right\}   \subseteq  \left\{  {x|{g}^{\prime }\left( x\right)  < 0}\right\}   \Longleftrightarrow  x > 2$ 时, ${g}^{\prime }\left( x\right)  = \ln x + {2ax} - 1 < 0$ 恒成立,
                
                所以 $x > 2$ 时, $a < \dfrac{1 - \ln x}{2x}$ 恒成立,设 $F\left( x\right)  = \dfrac{1 - \ln x}{2x},x > 2$ ,
                
                ${F}^{\prime }\left( x\right)  = \dfrac{\ln x - 2}{2{x}^{2}},$ 令 ${F}^{\prime }\left( x\right)  = 0,$ 得 $x = {{e}}^{2},$
                
                当 $x \in  \left( {2,{{e}}^{2}}\right)$ 时， ${F}^{\prime }\left( x\right)  < 0,F\left( x\right)$ 在 $\left( {2,{{e}}^{2}}\right)$ 上单调递减，
                
                当 $x \in  \left( {{{e}}^{2}, + \infty }\right)$ 时， ${F}^{\prime }\left( x\right)  > 0,F\left( x\right)$ 在 $\left( {{{e}}^{2}, + \infty }\right)$ 上单调递增，
                
                所以 $F{\left( x\right) }_{\min } = F\left( {{e}}^{2}\right)  =  - \dfrac{1}{2{{e}}^{2}}$ ,
                
                所以 $a <  - \dfrac{1}{2{{e}}^{2}}$ ,即 $a$ 的取值范围是 $\left( {-\infty , - \dfrac{1}{2{{e}}^{2}}}\right)$ .
                
                \item 因为 $g\left( x\right)  = x\left( {\ln x + {ax} - 2}\right)$ ,设 $m\left( x\right)  = \ln x + {ax} - 2$ ,则 ${m}^{\prime }\left( x\right)  = \dfrac{1}{x} + a$ 
                \begin{dingautolist}{172}
                    \item 若 $a <  - 1$ ，令 ${m}^{\prime }\left( x\right)  = 0$ ，解得 $x =  - \dfrac{1}{a}$
                
                    当 $x \in  \left( {0, - \dfrac{1}{a}}\right)$ 时， ${m}^{\prime }\left( x\right)  > 0$ ，所以 $m\left( x\right)$ 在 $\left( {0, - \dfrac{1}{a}}\right)$ 上单调递增，
                    
                    当 $x \in  \left( {-\dfrac{1}{a}, + \infty }\right)$ 时， ${m}^{\prime }\left( x\right)  < 0$ ，所以 $m\left( x\right)$ 在 $\left( {-\dfrac{1}{a}, + \infty }\right)$ 上单调递减，
                    
                    所以 $m\left( x\right)  \leq  m\left( {-\dfrac{1}{a}}\right)  = \ln \left( {-\dfrac{1}{a}}\right)  - 3 < 0$ ,
                    
                    所以 $a <  - 1$ 时, $m\left( x\right)  < 0 \Longrightarrow  g\left( x\right)  < 0 \Longrightarrow  h\left( x\right)  \leq  g\left( x\right)  < 0,h\left( x\right)$ 没有零点
                    
                    \item 若 $a > 1$ ，由(1)知，
                    
                    当 $x \in  \left( {0,2}\right)$ 时， ${f}^{\prime }\left( x\right)  > 0,f\left( x\right)$ 在 $\left( {0,2}\right)$ 上单调递增，
                    
                    又 $f\left( 0\right)  = a > 1$ ，所以 $x \in  \left( {0,2}\right)$ 时， $f\left( x\right)  > 0$ ，
                    
                    当 $x \in  \left( {0,2}\right)$ 时， ${m}^{\prime }\left( x\right)  = \dfrac{1}{x} + a > 0$ ，所以 $m\left( x\right)$ 在 $\left( {0,2}\right)$ 上单调递增，且 $m\left( \dfrac{l}{a}\right)  =  - \ln a - 1 < 0,m\left( 2\right)  =$
                    
                    $\ln 2 + {2a} - 2 > 0$ ,
                    
                    存在唯一 ${x}_{1} \in  \left( {0,2}\right)$ ,使得 $m\left( {x}_{1}\right)  = 0$ ,则 $g\left( {x}_{1}\right)  = 0,h\left( {x}_{1}\right)  = 0$ ,
                    
                    当 $x \in  \lbrack 2, + \infty )$ 时， ${m}^{\prime }\left( x\right)  = \dfrac{1}{x} + a > 0$ ，即 $m\left( x\right)$ 在 $\lbrack 2, + \infty )$ 单调递增，
                    
                    所以 $m\left( x\right)  = \ln x + {ax} - 2 > \ln 2 + {2a} - 2 > 0,g\left( x\right)  > 0$ ,
                    
                    当 $x \in  \lbrack 2, + \infty )$ 时， ${f}^{\prime }\left( x\right)  < 0,f\left( x\right)$ 在 $\lbrack 2, + \infty )$ 上单调递减，且 $f\left( 2\right)  = a + 4 > 0,f\left( {4a}\right)  =  - {64}{a}^{3} + {48}{a}^{2} + a <  - {64}{a}^{3} + {48}{a}^{3} + {a}^{3} =  - {15}{a}^{3} < 0,$
                    
                    所以存在唯一 ${x}_{2} \in  \left( {2, + \infty }\right)$ ,使得 $f\left( {x}_{2}\right)  = 0,h\left( {x}_{2}\right)  = 0$ 
                \end{dingautolist}
                综上所述， $a <  - 1$ 时， $h\left( x\right)$ 无零点，当 $a > 1$ 时 $h\left( x\right)$ 有 2 个零点.
            \end{enumerate}
        \end{jieda} 
        
        \item 设函数$f(x) = |x^3 - 6x^2+ax+b|$，若对任意的$a, b \in \mathbf{R}$，总存在$x_0 \in [0, 3]$，使得$f(x_0) \geq m$，则$m$的最大值为\blkda{2}
        \begin{jieda}
            根据题意可知对任意的$a, b \in \mathbf{R}$，$f(x)$在$[0, 3]$上的最大值$\geq m$ \\
            设$g_1(x) = x^3-6x^2, h_1(x) = -ax-b$，则$f(x) = |g_1(x) - h_1(x)|$，$g_1(x)$为一确定的曲线，$h_1(x)$为任意直线，故应有在$[0, 3]$内$g_1(x)$上的点与任意直线上横坐标相同的点在竖直方向上的距离的最大值$\geq m$ \\
            进一步可以考虑在保持该结构的基础上改变定曲线以使得其性质更好 \\
            设$g_2(x) = x^3-6x^2+9x, h_2(x) = (9-a)x-b$，则$f(x) = |g_2(x) - h_2(x)|$，此时$g_2(x)$在$[0, 1]$上单调递增，在$[1, 3]$上单调递减，且$g_2(0) = g_2(3) = 0$，$g_2(1) = 4$ \\
            结合图像可知当$h_2(x) = 2$时可以使得$f(x)$的最大值取得最小值2，因此$m \leq 2$
        \end{jieda}
    \end{enumerate}
\end{enumerate}